% chapters/problems.tex
% 问题 (Problems)
% 源: 原书 PDF 第 439--524 页 (印刷页 435--520)

\chapter*{问题}
\addcontentsline{toc}{chapter}{问题}
\markboth{问题}{问题}

括号中的数字表示求解该问题所需掌握的原书各章知识.

\begin{problem}[第~
\ref{chap:1}~章, 仅第 9 问需第~
\ref{chap:4}~章]
\index{KreinMilmandingli@Krein--Milman 定理 (Krein--Milman theorem)}
\index{jizhidian@极值点 (extreme point)}
\label{prob:1}
\runin{极值点; Krein--Milman 定理}
设 $E$ 为赋范向量空间, $K\subset E$ 为凸子集.若对满足 $x\ne y$ 的一切 $x,y\in K$ 及
一切 $t\in(0,1)$,
\begin{equation*}
tx+(1-t)y\ne a,
\end{equation*}
则称点 $a\in K$ 为\emph{极值点}.

1. 验证: $a\in K$ 是极值点当且仅当集合 $K\setminus\{a\}$ 是凸集.

2. 设 $a$ 为 $K$ 的极值点.设 $(x_i)_{1\le i\le n}$ 为 $K$ 中的有限列,
$(\alpha_i)_{1\le i\le n}$ 为满足对一切 $i$ 有 $\alpha_i>0$, $\sum\alpha_i=1$ 且
$\sum\alpha_ix_i=a$ 的实数有限列.证明对一切 $i$, $x_i=a$.

在下文中, 我们假设 $K\subset E$ 是 $E$ 的非空紧凸子集.若子集 $M\subset K$
非空, 闭, 且只要对某个 $t\in(0,1)$ 有 $tx+(1-t)y\in M$ ($x,y\in K$), 就有
$x\in M$ 且 $y\in M$, 则称 $M$ 为\emph{极端集}.

3. 设 $a\in K$.验证: $a$ 是极值点当且仅当 $\{a\}$ 是极端集.

我们的第一个目标是证明每个极端集包含至少一个极值点.

4. 设 $A\subset K$ 为极端集, $f\in E^\star$.令
\begin{equation*}
B=\left\{x\in A;\ \dual{f}{x}=\max_{y\in A}\dual{f}{y}\right\}.
\end{equation*}
证明 $B$ 是 $K$ 的极端子集.

5. 设 $M\subset K$ 为 $K$ 的极端集.考虑包含在 $M$ 内的 $K$ 的所有极端集的族
$\mathcal{F}$; $\mathcal{F}$ 配以下述偏序:
\begin{equation*}
A\le B\iff B\subset A.
\end{equation*}
证明 $\mathcal{F}$ 有极大元 $M_0$.

6. 证明 $M_0$ 退化为单点.\\
{[提示: 使用 Hahn--Banach 定理与第 4 问.]}

7. 得出结论.

8. 证明 $K$ 与其所有极值点的闭凸包重合.\\
{[提示: 用反证法并使用 Hahn--Banach 定理.]}

9. 确定在下述情形下 $B_E$ ($=E$ 的闭单位球) 的所有极值点的集合 $\mathcal{E}$:
\begin{enumerate}[label=(\alph*), itemsep=0.1em]
	\item $E=\ell^\infty$,
	\item $E=c$,
	\item $E=c_0$,
	\item $E=\ell^1$,
	\item $E=\ell^p$, 其中 $1<p<\infty$,
	\item $E=L^1(\R)$.
\end{enumerate}
[关于记号参见 \S \ref{sec:11.3}].
\end{problem}

\begin{problem}[第~
\ref{chap:1}~章, 仅第 B4 问需第~
\ref{chap:2}~章]
\index{ciweifen@次微分 (subdifferential)}
\label{prob:2}
\runin{凸函数的次微分}
设 $E$ 为赋范向量空间, $\varphi:E\to(-\infty,+\infty]$ 为满足
$\varphi\not\equiv+\infty$ 的凸函数.对每个 $x\in E$, $\varphi$ 的次微分定义为
\begin{equation*}
\begin{cases}
\partial\varphi(x)=\{f\in E^\star;\ \varphi(y)-\varphi(x)\ge\dual{f}{y-x}\ \forall y\in E\} & \text{若}\ x\in D(\varphi),\\
\partial\varphi(x)=\emptyset & \text{若}\ x\notin D(\varphi),
\end{cases}
\end{equation*}
我们令 $D(\partial\varphi)=\{x\in E;\ \partial\varphi(x)\ne\emptyset\}$, 从而
$D(\partial\varphi)\subset D(\varphi)$.构造一个使此包含关系为严格包含的例子.

\runin{-- A --}

1. 证明 $\partial\varphi(x)$ 是 $E^\star$ 的闭凸子集.

2. 设 $x_1,x_2\in D(\partial\varphi)$, $f_1\in\partial\varphi(x_1)$,
$f_2\in\partial\varphi(x_2)$.证明
\begin{equation*}
\dual{f_1-f_2}{x_1-x_2}\ge0.
\end{equation*}

3. 证明
\begin{equation*}
f\in\partial\varphi(x)\iff\varphi(x)+\varphi^\star(f)=\dual{f}{x}.
\end{equation*}

4. 在下述情形下确定 $\partial\varphi$:
\begin{enumerate}[label=(\alph*), itemsep=0.1em]
	\item $\varphi(x)=\frac12\norm{x}^2$,
	\item $\varphi(x)=\norm{x}$,
	\item $\varphi(x)=I_K(x)$ ($K$ 的示性函数), 其中 $K\subset E$ 为非空凸集 (相应地线性子空间),
	\item $\varphi(x)$ 为 $E$ 上的可微凸函数.
\end{enumerate}
[提示: 在情形 (a), (b) 中, $\partial\varphi$ 与第~
\ref{chap:1}~章注记 2 中定义的对偶映射 $F$
有关; 另见习题 1.1.]

5. 设 $\psi:E\to(-\infty,+\infty]$ 为另一个满足 $\psi\not\equiv+\infty$ 的凸函数.
假设 $D(\varphi)\cap D(\psi)\ne\emptyset$.证明对一切 $x\in E$,
\begin{equation*}
\partial\varphi(x)+\partial\psi(x)\subset\partial(\varphi+\psi)(x)
\end{equation*}
(约定若 $A=\emptyset$ 或 $B=\emptyset$, 则 $A+B=\emptyset$).构造一个使此包含
关系为严格包含的例子.

\runin{-- B --}

在整个第 B 部分中, 我们假设 $x_0\in E$ 满足条件
\begin{equation}
\exists M\in\R\ \text{与}\ \exists R>0\ \text{使得}\ \varphi(x)\le M
\quad\forall x\in E\ \text{满足}\ \norm{x-x_0}\le R.
\label{eq:prob2.1}
\end{equation}

1. 证明 $\partial\varphi(x_0)\ne\emptyset$.\\
{[提示: 在 $E\times\R$ 中使用 Hahn--Banach 定理.]}

2. 证明对一切 $f\in\partial\varphi(x_0)$,
$\norm{f}\le\frac1R(M-\varphi(x_0))$.

3. 由此推出: 对一切 $r<R$, 存在 $L\ge0$ 使得对一切满足 $\norm{x_i-x_0}\le r$
($i=1,2$) 的 $x_1,x_2\in E$,
\begin{equation*}
\abs{\varphi(x_1)-\varphi(x_2)}\le L\norm{x_1-x_2}.
\end{equation*}
[另一种证明另见习题 2.1.]

4. 这里假设 $E$ 为 Banach 空间且 $\varphi$ 下半连续.证明
\begin{equation*}
\operatorname{Int}D(\partial\varphi)=\operatorname{Int}D(\varphi).
\end{equation*}

5. 证明对每个 $y\in E$,
\begin{equation*}
\lim_{t\downarrow0}\frac{\varphi(x_0+ty)-\varphi(x_0)}{t}
=\max_{f\in\partial\varphi(x_0)}\dual{f}{y}.
\end{equation*}
[提示: 参看习题 1.25 第 5 问.]

6. 设 $\psi:E\to(-\infty,+\infty]$ 为满足 $x_0\in D(\psi)$ 的凸函数.证明对一切
$x\in E$,
\begin{equation*}
\partial\varphi(x)+\partial\psi(x)=\partial(\varphi+\psi)(x).
\end{equation*}
[提示: 给定 $f_0\in\partial(\varphi+\psi)(x)$, 对函数
$\tilde\varphi(y)=\varphi(y)-\varphi(x)-\dual{f_0}{y-x}$ 与
$\tilde\psi(y)=\psi(y)-\psi(x)$ 应用定理  \ref{thm:1.12}.]

\runin{-- C --}

1. 设 $\varphi:E\to\R$ 为满足对某个常数 $k\ge0$ 与 $C$, 对一切 $x\in E$
有 $\varphi(x)\le k\norm{x}+C$ 的凸函数.证明对一切 $x_1,x_2\in E$,
\begin{equation*}
\abs{\varphi(x_1)-\varphi(x_2)}\le k\norm{x_1-x_2}.
\end{equation*}
关于 $D(\varphi^\star)$ 能得出什么结论?

2. 设 $A\subset\R^n$ 为开凸集.设 $\varphi:A\to\R$ 为凸函数.证明 $\varphi$
在 $A$ 上连续.

\runin{-- D --}

设 $\varphi:E\to\R$ 为连续凸函数, 设 $C=\{x\in E;\ \varphi(x)\le0\}$.假设存在
某个 $x_0\in E$ 使得 $\varphi(x_0)<0$.给定 $x\in C$, 证明
$f\in\partial I_C(x)$ 当且仅当存在某个 $\lambda\in\R$ 使得
$f\in\lambda\partial\varphi(x)$, 其中若 $\varphi(x)<0$ 则 $\lambda=0$,
若 $\varphi(x)=0$ 则 $\lambda\ge0$.
\end{problem}

\begin{problem}[第~
\ref{chap:1}~章]
\index{Ekelandbianfenyuanli@Ekeland 变分原理 (Ekeland variational principle)}
\index{BishopPhelpsdingli@Bishop--Phelps 定理 (Bishop--Phelps theorem)}
\label{prob:3}
\runin{Ekeland, Brönsted--Rockafellar 与 Bishop--Phelps 定理; $\varepsilon$-次微分}

\runin{-- A --}

设 $M$ 为配以距离 $d(x,y)$ 的非空完备度量空间.设 $\psi:M\to(-\infty,+\infty]$
为有下界的, 下半连续的且满足 $\psi\not\equiv+\infty$ 的函数.我们的目标是
证明存在某个 $a\in M$ 使得
\begin{equation*}
\psi(x)-\psi(a)+d(x,a)\ge0\quad\forall x\in M.
\end{equation*}
给定 $x\in M$, 令
\begin{equation*}
S(x)=\{y\in M;\ \psi(y)-\psi(x)+d(x,y)\le0\}.
\end{equation*}

1. 验证 $x\in S(x)$, 且 $y\in S(x)\Rightarrow S(y)\subset S(x)$.

2. 固定实数的任意序列 $(\varepsilon_n)$, 满足对一切 $n$, $\varepsilon_n>0$,
且 $\varepsilon_n\to0$.给定 $x_0\in M$, 由归纳法构造序列 $(x_n)$ 如下: 一旦
已知 $x_n$, 选取满足下式的任意元素 $x_{n+1}$:
\begin{equation*}
\begin{cases}
x_{n+1}\in S(x_n),\\
\psi(x_{n+1})\le\inf_{x\in S(x_n)}\psi(x)+\varepsilon_{n+1}.
\end{cases}
\end{equation*}
验证对一切 $n$, $S(x_{n+1})\subset S(x_n)$, 且对一切 $n$ 与一切 $p$,
\begin{equation*}
\psi(x_{n+p})-\psi(x_n)+d(x_n,x_{n+p})\le0.
\end{equation*}
由此推出 $(x_n)$ 是 Cauchy 序列, 从而它收敛于一个极限, 记作 $a$.

3. 证明 $a$ 满足所要求的性质.\\
{[提示: 给定 $x\in M$, 考虑两种情形: 或者对一切 $n$, $x\in S(x_n)$; 或者存在
$N$ 使得 $x\notin S(x_N)$.]}

4. 给出几何解释.

\runin{-- B --}

设 $E$ 为 Banach 空间, $\varphi:E\to(-\infty,+\infty]$ 为满足
$\varphi\not\equiv+\infty$ 的凸下半连续函数.给定 $\varepsilon>0$ 与
$x\in D(\varphi)$, 令
\begin{equation*}
\partial_\varepsilon\varphi(x)=\{f\in E^\star;\ \varphi(x)+\varphi^\star(f)-\dual{f}{x}\le\varepsilon\}.
\end{equation*}
验证 $\partial_\varepsilon\varphi(x)\ne\emptyset$.

我们的目的是证明: 给定任何 $x_0\in D(\varphi)$ 与任何
$f_0\in\partial_\varepsilon\varphi(x_0)$, 下述性质成立:
\begin{equation*}
\forall\lambda>0,\ \exists x_1\in D(\varphi)\ \text{与}\ \exists f_1\in E^\star\
\text{满足}\ f_1\in\partial\varphi(x_1),\ \text{使得}\
\norm{x_1-x_0}\le\varepsilon/\lambda\ \text{且}\ \norm{f_1-f_0}\le\lambda.
\end{equation*}
(次微分 $\partial\varphi$ 在问题 2 中定义; 建议在做本题之前先做问题 2.)

1. 考虑由 $\psi(x)=\varphi(x)+\varphi^\star(f_0)-\dual{f_0}{x}$ 定义的函数 $\psi$.
证明存在某个 $x_1\in E$ 使得 $\norm{x_1-x_0}\le\varepsilon/\lambda$ 且对一切
$x\in E$,
\begin{equation*}
\psi(x)-\psi(x_1)+\lambda\norm{x-x_1}\ge0.
\end{equation*}
[提示: 在集合 $M=\{x\in E;\ \psi(x)\le\psi(x_0)-\lambda\norm{x-x_0}\}$ 上应用
第 A 部分的结果.]

2. 得出结论.\\
{[提示: 使用问题 2 第 B6 问的结果.]}

3. 由此推出 $\overline{D(\partial\varphi)}=\overline{D(\varphi)}$ 且
$\overline{R(\partial\varphi)}=\overline{D(\varphi^\star)}$, 其中
$R(\partial\varphi)=\{f\in E^\star;\ \exists x\in D(\partial\varphi)\ \text{使得}\ f\in\partial\varphi(x)\}$.

\runin{-- C --}

设 $E$ 为 Banach 空间, $C\subset E$ 为非空闭凸集.

1. 假设 $C$ 也是有界的, 证明集合
\begin{equation*}
\left\{f\in E^\star;\ \sup_{x\in C}\dual{f}{x}\ \text{可达到}\right\}
\end{equation*}
在 $E^\star$ 中稠密.\\
{[提示: 对函数 $\varphi=I_C$ 应用第 B 部分的结果.]}

2. 若 $E$ 的闭超平面 $H$ 分离 $C$ 与 $\{x\}$, 则称 $H$ 为 $C$ 在点 $x\in C$
处的\emph{支撑超平面}.证明: $C$ 中允许支撑超平面的点集在 $C$ 的边界
$\partial C$ ($=C\setminus\operatorname{Int}C$) 中稠密.
\end{problem}

\begin{problem}[第~
\ref{chap:1}~章]
\label{prob:4}
\runin{Asplund 定理与严格凸范数}
设 $E$ 为赋范向量空间, $\varphi_0,\psi_0:E\to[0,\infty)$ 为两个凸函数, 满足
$\varphi_0(0)=\psi_0(0)=0$ 以及对一切 $x\in E$, $0\le\psi_0(x)\le\varphi_0(x)$.
从 $\varphi_0$ 与 $\psi_0$ 出发, 由归纳法定义两个函数序列 $(\varphi_n)$ 与
$(\psi_n)$ 如下:
\begin{align*}
\varphi_{n+1}(x)&=\frac12(\varphi_n(x)+\psi_n(x)),\\
\psi_{n+1}(x)&=\frac12\inf_{y\in E}\{\varphi_n(x+y)+\psi_n(x-y)\}
=\frac12(\varphi_n\mathbin{\triangledown}\psi_n)(2x).
\end{align*}
[在开始本题之前, 先做涉及下卷积 $\triangledown$ 的习题 1.23.]

\runin{-- A --}

1. 验证对一切 $x\in E$ 与一切 $n$, $0\le\psi_n(x)\le\varphi_n(x)$, 且
$\varphi_n(0)=\psi_n(0)=0$.

2. 验证 $\varphi_n$ 与 $\psi_n$ 是凸的.

3. 证明序列 $(\varphi_n)$ 单调递减, 序列 $(\psi_n)$ 单调递增.由此推出
$(\varphi_n)$ 与 $(\psi_n)$ 有公共的极限, 记作 $\theta$, 满足
$\psi_0\le\theta\le\varphi_0$, 且 $\theta$ 是凸的.

4. 证明 $\varphi_n^\star\uparrow\theta^\star$.

5. 证明 $\psi_{n+1}^\star=\frac12(\varphi_n^\star+\psi_n^\star)$, 并由此推出
当 $D(\psi_0^\star)=E^\star$ 时 $\psi_n^\star\downarrow\theta^\star$.

6. 假设存在某个 $x_0\in D(\varphi_0)$ 使得 $\varphi_0$ 在 $x_0$ 处连续.
证明 $\varphi_n$ 与 $\psi_n$ 在 $x_0$ 处也连续.\\
{[提示: 应用习题 2.1 第 2 问.]}\\
由此推出 $\varphi_{n+1}^\star(f)=\frac12\inf_{g\in E^\star}\{\varphi_n^\star(f+g)+\psi_n^\star(f-g)\}$.

\runin{-- B --}

设 $\varphi:E\to[0,+\infty)$ 为二阶齐次的凸函数, 即对一切 $\lambda\in\R$ 与
一切 $x\in E$, $\varphi(\lambda x)=\lambda^2\varphi(x)$.证明对一切
$x,y\in E$ 与一切 $t\in(0,1)$,
\begin{equation*}
\varphi(x+y)\le\frac1t\varphi(x)+\frac{1}{1-t}\varphi(y).
\end{equation*}
由此推出函数 $x\mapsto\sqrt{\varphi(x)}$ 是半范数, 反之亦然.并建立
\begin{equation}
4\varphi(x)\le\frac1t\varphi(x+y)+\frac{1}{1-t}\varphi(x-y)
\quad\forall x,y\in E,\ \forall t\in(0,1).
\label{eq:prob4.1}
\end{equation}

在下文中, 我们额外假设 $\varphi_0$ 与 $\psi_0$ 是二阶齐次的, 且存在常数
$C>0$ 使得对一切 $x\in E$, $\varphi_0(x)\le(1+C)\psi_0(x)$.

1. 验证 $\varphi_n,\psi_n$ 与 $\theta$ 是二阶齐次的.

2. 证明对每个 $n$, 对一切 $x\in E$ 有
\begin{equation*}
\varphi_n(x)\le\left(1+\frac{C}{4^n}\right)\psi_n(x).
\end{equation*}
[提示: 用数学归纳法并使用 \eqref{eq:prob4.1}.]

3. 假设 $\varphi_0$ 或 $\psi_0$ 严格凸, 证明 $\theta$ 严格凸 (关于严格凸函数的
定义见习题 1.26).\\
{[提示: 使用第 B2 问中建立的不等式.把 $\varphi_n$ 分解为
$\varphi_n=\theta_n+\frac{1}{2^n}\varphi_0$ 是方便的, 其中 $\theta_n$ 是某个
凸函数, 不应试图把它显式写出.注意
$\theta_n+\left(\frac{1}{2^n}-\frac{C}{4^n}\right)\varphi_0\le\theta\le\theta_n+\frac{1}{2^n}\varphi_0$.]}

\runin{-- C --}

假设在 $E$ 上存在两个等价范数, 记作 $\norm{\cdot}_1$ 与 $\norm{\cdot}_2$.
用 $\norm{\cdot}_1^\star$ 与 $\norm{\cdot}_2^\star$ 记 $E^\star$ 上的相应
对偶范数.假设范数 $\norm{\cdot}_1$ 与 $\norm{\cdot}_2^\star$ 严格凸.
利用上述结果, 证明存在与 $\norm{\cdot}_1$ (从而与 $\norm{\cdot}_2$) 等价的
第三个范数 $\norm{\cdot}$, 它与其对偶范数 $\norm{\cdot}^\star$ 都严格凸.
\end{problem}

\begin{problem}[第 1, 2 章]
\label{prob:5}
\runin{正线性泛函}
设 $E$ 为赋范向量空间, $P$ 为以 $0$ 为顶点的凸锥, 即对一切 $x,y\in P$ 及一切
$\lambda,\mu>0$, $\lambda x+\mu y\in P$.令 $F=P-P$, 从而 $F$ 是线性子空间.
考虑下述两条性质:
\begin{enumerate}[label=(\roman*), itemsep=0.1em]
	\item $E$ 上满足对一切 $x\in P$ 有 $f(x)\ge0$ 的每个线性泛函 $f$ 在 $E$ 上连续.
	\item $F$ 是有限余维的闭子空间.
\end{enumerate}
本问题的目标是证明 (i) $\Rightarrow$ (ii); 反之, 当 $E$ 为 Banach 空间且 $P$
闭时, (ii) $\Rightarrow$ (i).

\runin{-- A --}

在整个第 A 部分中, 我们假设 (i) 成立.

1. 证明 $F$ 是闭的.\\
{[提示: 给定任何 $x_0\notin F$, 在 $E$ 上构造一个线性泛函 $f$, 满足 $f(x_0)=1$
且在 $F$ 上 $f=0$.]}

2. 设 $M$ 为 $E$ 的满足 $M\cap F=\{0\}$ 的任意线性子空间.证明 $\dim M<+\infty$.\\
{[提示: 使用习题 1.5.]}

3. 由此推出 (i) $\Rightarrow$ (ii).

\runin{-- B --}

在整个第 B 部分中, 我们假设 $E$ 为 Banach 空间且 $P$ 是闭的.

1. 这里额外假设 $P-P=E$.证明存在常数 $C>0$ 使得每个 $x\in E$ 有分解
$x=y-z$, 其中 $y,z\in P$, $\norm{y}\le C\norm{x}$ 且 $\norm{z}\le C\norm{x}$.\\
{[提示: 考虑集合 $K=\{x=y-z;\ y,z\in P,\ \norm{y}\le1,\ \norm{z}\le1\}$, 遵循
开映射定理 (定理  \ref{thm:2.6}) 证明的思想.]}

2. 由此推出 $P-P=E\Rightarrow\text{(i)}$.\\
{[提示: 用反证法, 考虑 $E$ 中的序列 $(x_n)$ 满足 $\norm{x_n}\le1/2^n$ 且
$f(x_n)\ge1$.然后使用第 B1 问的结果.]}

3. 证明 (ii) $\Rightarrow$ (i).

\runin{-- C --}

在下述例子中确定 $F=P-P$ 并检验 (i) 或 (ii) 是否成立:
\begin{enumerate}[label=(\alph*), itemsep=0.1em]
	\item 配以通常范数的 $E=C([0,1])$ 与 $P=\{u\in E;\ u(t)\ge0\ \forall t\in[0,1]\}$.
	\item 配以通常范数的 $E=C([0,1])$ 与 $P=\{u\in E;\ u(t)\ge0\ \forall t\in[0,1]\ \text{且}\ u(0)=u(1)=0\}$.
	\item 配以通常范数的 $E=\{u\in C^1([0,1]);\ u(0)=u(1)=0\}$ 与 $P=\{u\in E;\ u(t)\ge0\ \forall t\in[0,1]\}$.
\end{enumerate}
\end{problem}

\begin{problem}[第 1, 2 章]
\label{prob:6}
设 $E$ 为 Banach 空间, 设 $A:D(A)\subset E\to E^\star$ 为满足对一切
$x\in D(A)$, $\dual{Ax}{x}\ge0$ 的闭无界算子.

\runin{-- A --}

我们的目的是证明下述性质等价:
\begin{enumerate}[label=(\roman*), itemsep=0.1em]
	\item 对一切 $x\in D(A)$, 存在 $C(x)\in\R$ 使得对一切 $y\in D(A)$, $\dual{Ay}{y-x}\ge C(x)$,
	\item 存在 $k\ge0$ 使得对一切 $x,y\in D(A)$,
	\begin{equation*}
	\abs{\dual{Ay}{x}}\le k(\norm{x}+\norm{Ax})\sqrt{\dual{Ay}{y}}.
	\end{equation*}
\end{enumerate}

1. 证明 (ii) $\Rightarrow$ (i).

反之, 假设 (i) 成立.

2. 证明存在两个常数 $R>0$ 与 $M\ge0$ 使得对一切 $y\in D(A)$ 及一切满足
$\norm{x}+\norm{Ax}\le R$ 的 $x\in D(A)$, $\dual{Ay}{x-y}\le M$.\\
{[提示: 考虑函数 $\varphi(x)=\sup_{y\in D(A)}\dual{Ay}{x-y}$ 并应用习题 2.1.]}

3. 由此推出对一切 $y\in D(A)$ 及一切满足 $\norm{x}+\norm{Ax}\le R$ 的
$x\in D(A)$, $\abs{\dual{Ay}{x}}^2\le4M\dual{Ay}{y}$.

4. 得出结论.

\runin{-- B --}

在下文中假设 $D(A)=E$.设 $\alpha>0$.

1. 证明下述性质等价:
\begin{equation}
\norm{Ay}\le\alpha\sqrt{\dual{Ay}{y}}\quad\forall y\in E,
\tag{iii}
\label{eq:prob6.iii}
\end{equation}
\begin{equation}
\dual{Ay}{y-x}\ge-\frac1{4\alpha^2}\norm{x}^2\quad\forall x,y\in E.
\tag{iv}
\label{eq:prob6.iv}
\end{equation}
[提示: 使用与第 A 部分相同的方法.]

2. 设 $A^\star\in\LL(E^{\star\star},E^\star)$ 为 $A$ 的伴随算子.证明 \eqref{eq:prob6.iv} 等价于
\begin{equation}
\dual{A^\star y}{y-x}\ge-\frac1{4\alpha^2}\norm{x}^2\quad\forall x,y\in E.
\tag{iv$^\star$}
\label{eq:prob6.ivstar}
\end{equation}

3. 由此推出 \eqref{eq:prob6.iii} 等价于
\begin{equation}
\norm{A^\star y}\le\alpha\sqrt{\dual{A^\star y}{y}}\quad\forall y\in E.
\tag{iii$^\star$}
\label{eq:prob6.iiistar}
\end{equation}
\end{problem}

\begin{problem}[第 1, 2 章]
\label{prob:7}
\runin{两个无界线性算子之和的伴随算子}
设 $E$ 为 Banach 空间.给定 $E$ 中的两个闭线性子空间 $M$ 与 $N$, 令
\begin{equation*}
\rho(M,N)=\sup_{\substack{x\in M\\\norm{x}\le1}}\dist(x,N).
\end{equation*}

\runin{-- A --}

1. 验证 $\rho(M,N)\le1$; 若此外 $N\subset M$ 且 $N\ne M$, 证明 $\rho(M,N)=1$.\\
{[提示: 使用引理  \ref{lem:6.1}.]}

2. 设 $L,M$ 与 $N$ 为三个闭线性子空间.令 $a=\rho(M,N)$ 与 $b=\rho(N,L)$.
证明 $\rho(M,L)\le a+b+ab$.由此推出若 $L\subset M$, $a\le1/3$ 且 $b\le1/3$,
则 $L=M$.

3. 证明 $\rho(M,N)=\rho(N^\perp,M^\perp)$.\\
{[提示: 借助定理  \ref{thm:1.12} 验证对一切 $x\in E$ 与一切 $f\in E^\star$,
$\dist(x,N)=\sup_{g\in N^\perp,\norm{g}\le1}\dual{g}{x}$, 以及
$\dist(f,M^\perp)=\sup_{y\in M,\norm{y}\le1}\dual{f}{y}$.]}

\runin{-- B --}

设 $E$ 与 $F$ 为两个 Banach 空间; $E\times F$ 配以范数
$\norm{[u,v]}_{E\times F}=\norm{u}_E+\norm{v}_F$.给定稠密定义且闭的两个
无界算子 $A:D(A)\subset E\to F$ 与 $B:D(B)\subset E\to F$, 令
\begin{equation*}
\rho(A,B)=\rho(G(A),G(B)).
\end{equation*}

1. 证明 $\rho(A,B)=\rho(B^\star,A^\star)$.

2. 证明若 $D(A)\cap D(B)$ 在 $E$ 中稠密, 则
\begin{equation*}
A^\star+B^\star\subset(A+B)^\star.
\end{equation*}
[回忆 $D(A+B)=D(A)\cap D(B)$ 以及 $D(A^\star+B^\star)=D(A^\star)\cap D(B^\star)$.]
可能出现包含关系 $A^\star+B^\star\subset(A+B)^\star$ 是严格包含的情形——构造
这样的一个例子.我们的目的是证明在某些附加假设下等号成立.

3. 假设
\begin{equation}
D(A)\subset D(B)\ \text{且存在常数}\ k\in[0,1)\ \text{与}\ C\ge0\ \text{使得}\
\norm{Bu}\le k\norm{Au}+C\norm{u}\quad\forall u\in D(A).
\tag{H}
\label{eq:prob7.H}
\end{equation}
证明 $A+B$ 是闭的, 且 $\rho(A,A+B)\le k+C$.

4. 除 \eqref{eq:prob7.H} 外, 假设还成立
\begin{equation}
D(A^\star)\subset D(B^\star)\ \text{且存在常数}\ k^\star\in[0,1)\ \text{与}\
C^\star\ge0\ \text{使得}\ \norm{B^\star v}\le k^\star\norm{A^\star v}+C^\star\norm{v}\quad\forall v\in D(A^\star).
\tag{H$^\star$}
\label{eq:prob7.Hstar}
\end{equation}
设 $\varepsilon>0$ 使得 $\varepsilon(k+C)\le1/3$ 且 $\varepsilon(k^\star+C^\star)\le1/3$.
证明 $(A+\varepsilon B)^\star=A^\star+\varepsilon B^\star$.

5. 在假设 \eqref{eq:prob7.H} 与 \eqref{eq:prob7.Hstar} 下证明
$(A+B)^\star=A^\star+B^\star$.\\
{[提示: 使用相继步骤.验证下述不等式对一切 $t\in[0,1]$ 成立:
$\norm{Bu}\le\dfrac{k}{1-k}\norm{Au+tBu}+\dfrac{C}{1-k}\norm{u}\ \forall u\in D(A)$.]}
\end{problem}

\begin{problem}[第 2, 3 章, 仅第 6 问需第~
\ref{chap:4}~章]
\label{prob:8}
\runin{$\ell^1$ 中的弱收敛. Schur 定理}
设 $E=\ell^1$, 从而 $E^\star=\ell^\infty$ (见 \S \ref{sec:11.3}).给定 $x\in E$,
记 $x=(x_1,x_2,\dots,x_i,\dots)$, $\norm{x}_1=\sum_{i=1}^\infty\abs{x_i}$;
给定 $f\in E^\star$, 记 $f=(f_1,f_2,\dots,f_i,\dots)$,
$\norm{f}_\infty=\sup_i\abs{f_i}$.

设 $(x_n)$ 为 $E$ 中的序列, 使得 $x_n\rightharpoonup0$ 弱收敛于 $\sigma(E,E^\star)$.
我们的目标是证明 $\norm{x_n}_1\to0$.

1. 给定 $f,g\in B_{E^\star}$ (即 $\norm{f}_\infty\le1$ 且 $\norm{g}_\infty\le1$), 令
\begin{equation*}
d(f,g)=\sum_{i=1}^\infty\frac{1}{2^i}\abs{f_i-g_i}.
\end{equation*}
验证 $d$ 是 $B_{E^\star}$ 上的度量, 且 $B_{E^\star}$ 关于相应的拓扑是紧的.

2. 给定 $\varepsilon>0$, 令
\begin{equation*}
F_k=\{f\in B_{E^\star};\ \abs{\dual{f}{x_n}}\le\varepsilon\ \forall n\ge k\}.
\end{equation*}
证明存在某个 $f^0\in B_{E^\star}$, 常数 $\rho>0$, 以及整数 $k_0$, 使得
\begin{equation*}
[f\in B_{E^\star}\ \text{且}\ d(f,f^0)<\rho]\implies[f\in F_{k_0}].
\end{equation*}
[提示: 使用 Baire 纲定理.]

3. 固定满足 $(1/2^{N-1})<\rho$ 的整数 $N$.证明对一切 $n\ge k_0$,
\begin{equation*}
\norm{x_n}_1\le\varepsilon+2\sum_{i=1}^N\abs{x_i^n}.
\end{equation*}

4. 得出结论.

5. 用类似的方法证明: 若 $(x_n)$ 是 $\ell^1$ 中的序列使得对每个 $f\in\ell^\infty$,
序列 $(\dual{f}{x_n})$ 收敛于某个极限, 则 $(x_n)$ 在 $\ell^1$ 中强收敛于一个极限.

6. 考虑 $E=L^1(0,1)$, 从而 $E^\star=L^\infty(0,1)$.在 $E$ 中构造一个序列
$(u_n)$ 使得 $u_n\rightharpoonup0$ 弱收敛于 $\sigma(E,E^\star)$, 且对一切 $n$,
$\norm{u_n}_1=1$.
\end{problem}

\begin{problem}[第 1, 2, 3 章]
\label{prob:9}
\runin{弱$^\star$拓扑的 Hahn--Banach 定理及其应用}
设 $E$ 为 Banach 空间.

\runin{-- A --}

1. 设 $A\subset E^\star$ 与 $B\subset E^\star$ 为满足 $A\cap B=\emptyset$ 的两个
非空凸集.假设 $A$ 在拓扑 $\sigma(E^\star,E)$ 中开.证明存在某个 $x\in E$,
$x\ne0$, 以及常数 $\alpha$, 使得超平面 $\{f\in E^\star;\ \dual{f}{x}=\alpha\}$
分离 $A$ 与 $B$.

2. 假设 $A\subset E^\star$ 在 $\sigma(E^\star,E)$ 中闭且 $B\subset E^\star$ 在
$\sigma(E^\star,E)$ 中紧.证明 $A+B$ 在 $\sigma(E^\star,E)$ 中闭.

3. 设 $A\subset E^\star$ 与 $B\subset E^\star$ 为满足 $A\cap B=\emptyset$ 的两个
非空凸集.假设 $A$ 在 $\sigma(E^\star,E)$ 中闭且 $B$ 在 $\sigma(E^\star,E)$ 中紧.
证明存在某个 $x\in E$, $x\ne0$, 以及常数 $\alpha$, 使得超平面
$\{f\in E^\star;\ \dual{f}{x}=\alpha\}$ 严格分离 $A$ 与 $B$.

4. 设 $A\subset E^\star$ 为凸集.证明 $A$ 在 $\sigma(E^\star,E)$ 中的闭包
$\overline A^{\sigma(E^\star,E)}$ 是凸的.

\runin{-- B --}

下面是上述结果的若干应用:

1. 设 $N\subset E^\star$ 为线性子空间.回忆
\begin{equation*}
N^\perp=\{x\in E;\ \dual{f}{x}=0\ \forall f\in N\}\quad\text{且}\quad
N^{\perp\perp}=\{f\in E^\star;\ \dual{f}{x}=0\ \forall x\in N^\perp\}.
\end{equation*}
证明 $N^{\perp\perp}=\overline N^{\sigma(E^\star,E)}$.若 $E$ 自反, 能得出什么结论?
由此推出 $c_0$ 在拓扑 $\sigma(\ell^\infty,\ell^1)$ 下在 $\ell^\infty$ 中稠密.

2. 设 $\varphi:E\to(-\infty,+\infty]$ 为凸下半连续函数, $\varphi\not\equiv+\infty$.
证明 $\psi=\varphi^\star$ 在拓扑 $\sigma(E^\star,E)$ 下下半连续.

反之, 给定关于拓扑 $\sigma(E^\star,E)$ 下半连续的凸函数
$\psi:E^\star\to(-\infty,+\infty]$ ($\psi\not\equiv+\infty$), 证明存在凸下半
连续函数 $\varphi:E\to(-\infty,+\infty]$ ($\varphi\not\equiv+\infty$) 使得
$\psi=\varphi^\star$.

3. 设 $F$ 为另一个 Banach 空间, 设 $A:D(A)\subset E\to F$ 为稠密定义且闭的
无界线性算子.证明:
\begin{enumerate}[label=(\roman*), itemsep=0.1em]
	\item $\overline{R(A^\star)}^{\sigma(E^\star,E)}=N(A)^\perp$,
	\item $\overline{D(A^\star)}^{\sigma(F^\star,F)}=F^\star$.
\end{enumerate}
若 $E$ (相应地 $F$) 自反, 能得出什么结论?

4. 不借助引理~
\ref{lem:3.3} 证明: $J(B_E)$ 在拓扑 $\sigma(E^{\star\star},E^\star)$ 下在
$B_{E^{\star\star}}$ 中稠密 (见引理  \ref{lem:3.4}).

5. 设 $A:B_E\to E^\star$ 为单调映射, 即对一切 $x,y\in B_E$,
$\dual{Ax-Ay}{x-y}\ge0$.令 $S_E=\{x\in E;\ \norm{x}=1\}$.证明
$A(B_E)\subset\overline{\operatorname{conv}A(S_E)}^{\sigma(E^\star,E)}$.
\end{problem}

\begin{problem}[第~
\ref{chap:3}~章]
\index{Eberleinmuliandingli@Eberlein--Šmulian 定理 (Eberlein--Smulian theorem)}
\label{prob:10}
\runin{Eberlein--Šmulian 定理}
设 $E$ 为 Banach 空间, $A\subset E$.令 $B=\overline A^{\sigma(E,E^\star)}$.
本问题的目标是证明下述性质等价:
\begin{enumerate}[label=(\Alph*), itemsep=0.1em]
	\item[(P)] $B$ 在拓扑 $\sigma(E,E^\star)$ 中是紧的.
	\item[(Q)] $A$ 中的每个序列 $(x_n)$ 都有弱收敛的子序列.
\end{enumerate}
此外, (P) (或 (Q)) 蕴含下述性质:
\begin{enumerate}[label=(\Alph*), itemsep=0.1em]
	\item[(R)] 对每个 $y\in B$, 存在序列 $(y_n)\subset A$ 使得 $y_n\rightharpoonup y$ 弱收敛于 $\sigma(E,E^\star)$.
\end{enumerate}

\runin{-- A --}

断言 (P) $\Rightarrow$ (Q) 的证明.

1. 在 $E$ 可分的额外假设下证明 (P) $\Rightarrow$ (Q).\\
{[提示: 考虑 $B_{E^\star}$ 中的一个可数且对拓扑 $\sigma(E^\star,E)$ 在
$B_{E^\star}$ 中稠密的集合 $(b_k)$ (为什么这样的集合存在?).验证量
$d(x,y)=\sum_{k=1}^\infty\frac{1}{2^k}\abs{\dual{b_k}{x-y}}$ 是度量, 并由此推出
$B$ 对 $\sigma(E,E^\star)$ 是可度量化的.]}

2. 在一般情形下证明 (P) $\Rightarrow$ (Q).\\
{[提示: 使用第 A1 问.]}

\runin{-- B --}

后文将需要下述引理:

\runin{引理.} 设 $F$ 为赋范向量空间, $M\subset F^\star$ 为有限维向量空间.则
在 $B_F$ 中存在有限子集 $(a_i)_{1\le i\le k}$ 使得对一切 $g\in M$,
\begin{equation*}
\max_{1\le i\le k}\dual{g}{a_i}\ge\frac12\norm{g}.
\end{equation*}
[提示: 先在 $S_M$ 中选取点 $(g_i)_{1\le i\le k}$ 使得
$S_M\subset\bigcup_{i=1}^kB(g_i,1/4)$, 其中 $S_M=\{g\in M;\ \norm{g}=1\}$.]

\runin{-- C --}

设 $\xi\in E^{\star\star}$ 使得 $\xi\in\overline A^{\sigma(E^{\star\star},E^\star)}$.
利用假设 (Q), 我们将证明 $\xi\in B$ 且存在序列 $(y_k)\subset A$ 使得在
$\sigma(E,E^\star)$ 中 $y_k\rightharpoonup\xi$.

1. 令 $n_1=1$, 固定任何 $f_1\in B_{E^\star}$.证明存在某个 $x_1\in A$ 使得
$\abs{\dual{\xi}{f_1}-\dual{f_1}{x_1}}<1$.

2. 设 $M_1=[\xi,x_1]$ 为由 $\xi$ 与 $x_1$ 张成的线性空间.证明在 $B_{E^\star}$
中存在 $(f_i)_{1<i\le n_2}$ 使得对一切 $\eta\in M_1$,
$\max_{1<i\le n_2}\dual{\eta}{f_i}\ge\frac12\norm{\eta}$.证明存在某个
$x_2\in A$ 使得对一切 $1\le i\le n_2$, $\abs{\dual{\xi}{f_i}-\dual{f_i}{x_2}}<1/2$.

3. 重复上述构造, 我们得到两个序列 $(x_k)\subset A$ 与 $(f_i)\subset B_{E^\star}$,
以及递增的整数序列 $(n_k)$, 使得
\begin{align*}
\max_{n_k<i\le n_{k+1}}\dual{\eta}{f_i}&\ge\frac12\norm{\eta}\quad\forall\eta\in M_k=[\xi,x_1,x_2,\dots,x_k], \tag{a}\\
\abs{\dual{\xi}{f_i}-\dual{f_i}{x_{k+1}}}&<\frac{1}{k+1}\quad\forall i,\ 1\le i\le n_{k+1}. \tag{b}
\end{align*}

4. 由 (a) 推出对一切 $\eta\in\bigcup_{k=1}^\infty M_k=\mathcal{M}$,
$\sup_{i\ge1}\dual{\eta}{f_i}\ge\frac12\norm{\eta}$, 进而对一切 $\eta\in\overline{\mathcal{M}}$,
$\sup_{i\ge1}\dual{\eta}{f_i}\ge\frac12\norm{\eta}$, 其中 $\overline{\mathcal{M}}$
表示 $\mathcal{M}$ 在 $E^{\star\star}$ 中关于强拓扑的闭包.

5. 利用 (b) 与假设 (Q), 证明存在某个 $x\in B\cap\overline{\mathcal{M}}$ 使得
对一切 $i\ge1$, $\dual{\xi}{f_i}=\dual{f_i}{x}$.由此推出 $\xi=x$ 并得出结论.

\runin{-- D --}

1. 证明 (Q) $\Rightarrow$ (P).

2. 证明 (Q) $\Rightarrow$ (R).
\end{problem}
\begin{problem}[第~
\ref{chap:3}~章]
\label{prob:11}
\runin{Banach--Dieudonné--Krein--Šmulian 定理}
设 $E$ 为 Banach 空间, $C\subset E^\star$ 为凸集.假设对每个整数 $n$, 集合
$C\cap(nB_{E^\star})$ 在拓扑 $\sigma(E^\star,E)$ 中是闭的.本问题的目标是证明 $C$
在拓扑 $\sigma(E^\star,E)$ 中是闭的.

\runin{-- A --}

额外假设 $0\notin C$.我们将证明存在 $E$ 中的序列 $(x_n)$ 使得
\begin{equation}
\norm{x_n}\to0\quad\text{且}\quad\sup_n\dual{f}{x_n}>1\quad\forall f\in C.
\label{eq:prob11.1}
\end{equation}
令 $d=\dist(0,C)$, 考虑满足 $d_1>d$ 的递增序列 $d_n\uparrow+\infty$.令
$C_k=\{f\in C;\ \norm{f}\le d_k\}$.

1. 验证集合 $C_k$ 在拓扑 $\sigma(E^\star,E)$ 中是紧的.证明存在某个 $f_0\in C$
使得 $d=\norm{f_0}>0$.

2. 证明存在某个 $x_1\in E$ 使得对一切 $f\in C_1$, $\dual{f}{x_1}>1$.\\
{[提示: 使用弱$^\star$拓扑的 Hahn--Banach 定理; 见问题 9 第 A3 问.]}

3. 令 $A_1=\{x_1\}$.证明存在有限子集 $A_2\subset E$ 使得
$A_2\subset\frac{1}{d_1}B_E$ 且对一切 $f\in C_2$, $\sup_{x\in A_1\cup A_2}\dual{f}{x}>1$.\\
{[提示: 对满足 $A\subset\frac{1}{d_1}B_E$ 的每个有限子集 $A\subset E$, 考虑集合
$Y_A=\left\{f\in C_2;\ \sup_{x\in A_1\cup A}\dual{f}{x}\le1\right\}$, 先证明
$\bigcap_A Y_A=\emptyset$.]}

4. 由归纳法构造有限子集 $A_k\subset E$, 使得
$A_k\subset\frac{1}{d_{k-1}}B_E$ 且对一切 $f\in C_k$, $\sup_{x\in\bigcup_{i=1}^kA_i}\dual{f}{x}>1$.

5. 构造满足 \eqref{eq:prob11.1} 的序列 $(x_n)$.

\runin{-- B --}

1. 再次假设 $0\notin C$.证明存在某个 $x\in E$ 使得对一切 $f\in C$, $\dual{f}{x}\ge1$.\\
{[提示: 设 $(x_n)$ 为满足 \eqref{eq:prob11.1} 的序列.考虑由 $T(f)=(\dual{f}{x_n})_n$
定义的算子 $T:E^\star\to c_0$, 并在 $c_0$ 中分离 $T(C)$ 与 $c_0$ 的开单位球.]}

2. 得出结论.
\end{problem}

\begin{problem}[第 1, 2, 3 章]
\label{prob:12}
在开始本题之前, 必须先做习题 1.23.

设 $E$ 为自反 Banach 空间, $\varphi,\psi:E\to(-\infty,+\infty]$ 为满足
$D(\varphi)\cap D(\psi)\ne\emptyset$ 的凸下半连续函数.令
$\theta=\varphi^\star\mathbin{\triangledown}\psi^\star$.

\runin{-- A --}

我们断言 $D((\varphi+\psi)^\star)=D(\varphi^\star)+D(\psi^\star)$.

1. 证明 $D(\varphi^\star)+D(\psi^\star)\subset D((\varphi+\psi)^\star)$.

2. 证明 $\theta$ 将 $E^\star$ 映入 $(-\infty,+\infty]$, $\theta$ 是凸的,
$D(\theta)=D(\varphi^\star)+D(\psi^\star)$, 且 $\theta^\star=\varphi+\psi$.

3. 由此推出 $D((\varphi+\psi)^\star)=D(\theta)$ 并得出结论.

\runin{-- B --}

额外假设 $\varphi$ 与 $\psi$ 满足
\begin{equation}
\bigcup_{\lambda\ge0}\lambda(D(\varphi)-D(\psi))=E.
\tag{H}
\label{eq:prob12.H}
\end{equation}
我们断言:
\begin{enumerate}[label=(\roman*), itemsep=0.1em]
	\item $(\varphi+\psi)^\star=\varphi^\star\mathbin{\triangledown}\psi^\star$,
	\item $\inf_{x\in E}\{\varphi(x)+\psi(x)\}=\max_{g\in E^\star}\{-\varphi^\star(-g)-\psi^\star(g)\}$,
	\item $D((\varphi+\psi)^\star)=D(\varphi^\star)+D(\psi^\star)$.
\end{enumerate}

1. 证明对每个固定的 $f\in E^\star$ 与 $\alpha\in\R$, 集合
$M=\{g\in E^\star;\ \varphi^\star(f-g)+\psi^\star(g)\le\alpha\}$ 是有界的.\\
{[提示: 使用假设 \eqref{eq:prob12.H} 与推论  \ref{cor:2.5}.]}

2. 设 $\alpha\in\R$ 固定.设 $(f_n)$ 与 $(g_n)$ 为 $E^\star$ 中的两个序列使得
$(f_n)$ 有界且 $\varphi^\star(f_n-g_n)+\psi^\star(g_n)\le\alpha$.证明 $(g_n)$ 有界.

3. 由此推出 $\theta$ 是下半连续的.

4. 证明 (i), (ii) 与 (iii).将这些结果与习题 1.23 第 3 问以及定理  \ref{thm:1.12} 作比较.
\end{problem}

\begin{problem}[第 1, 3 章]
\label{prob:13}
\runin{对偶映射的性质. 一致凸性. 范数的可微性}
设 $E$ 为 Banach 空间.回忆对偶映射的定义 (见第~
\ref{chap:1}~章注记 2): 对每个 $x\in E$,
\begin{equation*}
F(x)=\{f\in E^\star;\ \norm{f}=\norm{x}\ \text{且}\ \dual{f}{x}=\norm{x}^2\}.
\end{equation*}
在开始本题之前, 先做习题 1.1 与 1.25 是有益的.

\runin{-- A --}

假设 $E^\star$ 严格凸, 从而 $F(x)$ 只由单个元素组成.

1. 验证对一切 $x,y\in E$,
\begin{equation*}
\lim_{\lambda\to0}\frac{1}{2\lambda}(\norm{x+\lambda y}^2-\norm{x}^2)=\dual{Fx}{y}.
\end{equation*}
[提示: 应用习题 1.25 的结果; 区分 $\lambda>0$ 与 $\lambda<0$ 的情形.]

2. 证明对一切 $x,y\in E$, 映射 $t\in\R\mapsto\dual{F(x+ty)}{y}$ 在 $t=0$ 处连续.\\
{[提示: 对 $u=x+ty$ 与 $v=x+\lambda y$ 使用不等式 $\frac12(\norm{v}^2-\norm{u}^2)\ge\dual{Fu}{v-u}$.]}

3. 由此推出 $F$ 从 $E$ 强到 $E^\star$ 弱$^\star$连续.\\
{[提示: 使用习题 3.11 的结果.]}\\
当 $E$ 自反或可分时, 用简捷的直接方法证明同样的结果.

4. 验证对一切 $x,y\in E$,
\begin{equation*}
\dual{Fx+Fy}{x+y}+\dual{Fx-Fy}{x-y}=2(\norm{x}^2+\norm{y}^2).
\end{equation*}
由此推出对一切满足 $\norm{x}=\norm{y}=1$ 的 $x,y\in E$,
$\norm{Fx+Fy}+\norm{x-y}\ge2$.

5. 额外假设 $E$ 自反且严格凸.证明 $F$ 是由 $E$ 到 $E^\star$ 上的双射.
验证 $F^{-1}$ 与 $E^\star$ 的对偶映射重合.

\runin{-- B --}

在本部分中假设 $E^\star$ 一致凸.

1. 证明 $F$ 从 $E$ 强到 $E^\star$ 强连续.

2. 更精确地说, 证明 $F$ 在 $E$ 的有界集上一致连续.\\
{[提示: 用反证法并应用第 A4 问.]}

3. 由此推出函数 $\varphi(x)=\frac12\norm{x}^2$ 可微且其微分是 $F$, 即对每个 $x_0\in E$,
\begin{equation*}
\lim_{\substack{x\to x_0\\x\ne x_0}}\frac{\varphi(x)-\varphi(x_0)-\dual{Fx_0}{x-x_0}}{\norm{x-x_0}}=0.
\end{equation*}

\runin{-- C --}

反之, 假设对每个 $x\in E$, 集合 $Fx$ 只包含单个元素, 且 $F$ 在 $E$ 的有界集上一致连续.
证明 $E^\star$ 一致凸.\\
	{[提示: 先证明对一切 $y\in E$ 与一切 $f,g\in E^\star$,
	\begin{equation*}
	\norm{f+g}\le\frac12\norm{f}^2+\frac12\norm{g}^2-\dual{f-g}{y}+\sup_{\substack{x\in E\\\norm{x}\le1}}\{\varphi(x+y)+\varphi(x-y)\}.
	\end{equation*}]}
\end{problem}

\begin{problem}[第 1, 3 章]
\label{prob:14}
\runin{通过下卷积对凸函数作正则化}
设 $E$ 为使得 $E^\star$ 一致凸的 Banach 空间.假设 $\varphi:E\to(-\infty,+\infty]$
为凸下半连续函数且 $\varphi\not\equiv+\infty$.本问题的目标是证明存在可微凸函数序列
$(\varphi_n)$ 使得当 $n\uparrow+\infty$ 时 $\varphi_n\uparrow\varphi$.

\runin{-- A --}

对每个固定的 $x\in E$, 考虑由
\begin{equation*}
\Phi_x(f)=\frac12\norm{f}^2+\varphi^\star(f)-\dual{f}{x},\quad f\in E^\star
\end{equation*}
定义的函数 $\Phi_x:E^\star\to(-\infty,+\infty]$.

1. 验证存在唯一的元素 $f_x\in E^\star$ 使得
$\Phi_x(f_x)=\inf_{f\in E^\star}\Phi_x(f)$.令 $Sx=f_x$.

2. 证明映射 $x\mapsto Sx$ 从 $E$ 强到 $E^\star$ 强连续.\\
{[提示: 先证明 $S$ 从 $E$ 强到 $E^\star$ 弱$^\star$连续.]}

\runin{-- B --}

考虑由 $\psi(x)=\inf_{y\in E}\left\{\frac12\norm{x-y}^2+\varphi(y)\right\}$ 定义的
函数 $\psi:E\to\R$.我们断言 $\psi$ 是凸的, 可微的, 且其微分与 $S$ 重合.

1. 验证 $\psi$ 是凸的, 且对一切 $x\in E$,
\begin{equation}
\psi(x)=-\min_{f\in E^\star}\left\{\frac12\norm{f}^2+\varphi^\star(f)-\dual{f}{x}\right\},
\tag{i}
\end{equation}
对一切 $f\in E^\star$,
\begin{equation}
\psi^\star(f)=\frac12\norm{f}^2+\varphi^\star(f).
\tag{ii}
\end{equation}
[提示: 应用定理  \ref{thm:1.12}.]

2. 由此推出对一切 $x\in E$, $\psi(x)+\psi^\star(Sx)=\dual{Sx}{x}$, 且对一切 $x,y\in E$,
\begin{equation*}
\abs{\psi(y)-\psi(x)-\dual{Sx}{y-x}}\le\norm{Sy-Sx}\norm{y-x}.
\end{equation*}

3. 得出结论.

\runin{-- C --}

对每个整数 $n\ge1$ 与每个 $x\in E$, 令
\begin{equation*}
\varphi_n(x)=\inf_{y\in E}\left\{\frac n2\norm{x-y}^2+\varphi(y)\right\}.
\end{equation*}
证明 $\varphi_n$ 是凸的, 可微的, 且对每个 $x\in E$, 当 $n\uparrow+\infty$ 时
$\varphi_n(x)\uparrow\varphi(x)$.\\
{[提示: 使用与习题 1.24 相同的方法.]}
\end{problem}

\begin{problem}[第~
\ref{chap:1}~章, 仅第 B6 问需第~
\ref{chap:5}~章]
\label{prob:15}
\runin{Chebyshev 意义下的集合中心. 正规结构. Edelstein 意义下的序列渐近中心. Kirk, Browder, Göhde 与 Edelstein 的紧致映射不动点定理}
在开始本题之前, 先做习题 3.29 是有益的.

设 $E$ 为一致凸 Banach 空间, $C\subset E$ 为非空闭凸集.

\runin{-- A --}

设 $A\subset C$ 为非空有界集.对每个 $x\in E$, 定义
\begin{equation*}
\phi(x)=\sup_{y\in A}\norm{x-y}.
\end{equation*}

1. 验证 $\phi$ 是凸函数, 且对一切 $x_1,x_2\in E$,
$\abs{\phi(x_1)-\phi(x_2)}\le\norm{x_1-x_2}$.

2. 证明存在唯一的元素 $c\in C$ 使得 $\phi(c)=\inf_{x\in C}\phi(x)$.点 $c$
称为 $A$ 的\emph{中心}, 记作 $c=\sigma(A)$.

3. 证明若 $A$ 不退化为单点, 则
\begin{equation*}
\phi(\sigma(A))<\operatorname{diam}A=\sup_{x,y\in A}\norm{x-y}.
\end{equation*}

\runin{-- B --}

设 $(a_n)$ 为 $C$ 中的有界序列; 令 $A_n=\bigcup_{i=n}^\infty\{a_i\}$,
对 $x\in E$ 令 $\phi_n(x)=\sup_{y\in A_n}\norm{x-y}$.

1. 对每个 $x\in E$, 考虑 $\phi(x)=\lim_{n\to+\infty}\phi_n(x)$.证明这个极限存在,
且 $\phi$ 在 $E$ 上是凸的, 连续的.

2. 证明存在唯一的元素 $\sigma\in C$ 使得 $\phi(\sigma)=\inf_{x\in C}\phi(x)$.
点 $\sigma$ 称为序列 $(a_n)$ 的\emph{渐近中心}.

3. 设 $\sigma_n=\sigma(A_n)$ 为在第 A2 问意义下集合 $A_n$ 的中心.证明
\begin{equation*}
\lim_{n\to\infty}\phi_n(\sigma_n)=\lim_{n\to\infty}\phi(\sigma_n)=\phi(\sigma),
\end{equation*}
且 $\sigma_n\rightharpoonup\sigma$ 弱收敛于 $\sigma(E,E^\star)$.

4. 由此推出 $\sigma_n\to\sigma$ 强收敛.\\
{[提示: 用反证法并应用习题 3.29 的结果.]}

5. 假设 $a_n\to a$ 强收敛.确定序列 $(a_n)$ 的渐近中心.

6. 这里假设 $E$ 为 Hilbert 空间且 $a_n\rightharpoonup a$ 弱收敛于 $\sigma(E,E^\star)$.
计算 $\phi(x)$ 并确定序列 $(a_n)$ 的渐近中心.\\
{[提示: 展开范数的平方.]}

\runin{-- C --}

假设 $T:C\to C$ 是压缩映射, 即对一切 $x,y\in C$, $\norm{Tx-Ty}\le\norm{x-y}$.
设 $a\in C$ 给定, 设 $a_n=T^na$ 为其迭代序列.假设序列 $(a_n)$ 有界.设 $\sigma$
为序列 $(a_n)$ 的渐近中心.

1. 证明 $\sigma$ 是 $T$ 的不动点, 即 $T\sigma=\sigma$.

2. 验证 $T$ 的不动点集是闭凸集.
\end{problem}
\begin{problem}[第 2, 3 章]
\label{prob:16}
\runin{线性极大单调算子的表征}
设 $E$ 为 Banach 空间, 设 $A:D(A)\subset E\to E^\star$ 为满足单调性条件
\begin{equation}
\dual{Au}{u}\ge0\quad\forall u\in D(A)
\tag{M}
\label{eq:prob16.M}
\end{equation}
的无界线性算子.我们用 (P) 记下述性质:
\begin{equation}
\text{若}\ x\in E\ \text{且}\ f\in E^\star\ \text{满足}\
\dual{Au-f}{u-x}\ge0\quad\forall u\in D(A),
\ \text{则}\ x\in D(A)\ \text{且}\ Ax=f.
\tag{P}
\label{eq:prob16.P}
\end{equation}

\runin{-- A --}

1. 证明若 (P) 成立, 则 $D(A)$ 在 $E$ 中稠密.\\
{[提示: 证明若 $f\in E^\star$ 且对一切 $u\in D(A)$ 有 $\dual{f}{u}=0$, 则 $f=0$.]}

2. 证明若 (P) 成立, 则 $A$ 是闭的.

3. 证明函数 $u\in D(A)\mapsto\dual{Au}{u}$ 是凸的.

4. 证明 $N(A)\subset R(A)^\perp$.由此推出若 $D(A)$ 在 $E$ 中稠密, 则 $N(A)\subset N(A^\star)$.

5. 证明若 $D(A)=E$, 则 (P) 成立.

在整个本题的其余部分中, 我们额外假设:
\begin{enumerate}[label=(\roman*), itemsep=0.1em]
	\item $E$ 自反, $E$ 与 $E^\star$ 严格凸,
	\item $D(A)$ 在 $E$ 中稠密且 $A$ 是闭的,
\end{enumerate}
从而 $A^\star:D(A^\star)\subset E\to E^\star$ 且 $D(A^\star)$ 在 $E$ 中稠密 (为什么?).
本问题的目标是建立等价性 (P) $\iff$ (M$^\star$), 其中 (M$^\star$) 记下述性质:
\begin{equation}
\dual{A^\star v}{v}\ge0\quad\forall v\in D(A^\star).
\tag{M$^\star$}
\label{eq:prob16.Mstar}
\end{equation}

\runin{-- B --}

在本节中假设 (P) 成立.

1. 证明对一切 $v\in D(A)\cap D(A^\star)$, $\dual{A^\star v}{v}\ge0$.

2. 设 $v\in D(A^\star)$ 满足 $v\notin D(A)$.证明对一切 $f\in E^\star$, 存在 $u\in D(A)$
使得 $\dual{Au-f}{u-v}<0$.选取 $f=-A^\star v$, 证明 $\dual{A^\star v}{v}>0$.
由此推出 (M$^\star$) 成立.

3. 证明 $N(A)=N(A^\star)$ 且 $R(A)=R(A^\star)$.

\runin{-- C --}

在本部分中假设 (M$^\star$) 成立.

1. 验证配以图范数 $\norm{u}_{D(A)}=\norm{u}_E+\norm{Au}_{E^\star}$ 的空间 $D(A)$
是自反的.

2. 给定 $x\in E$ 与 $f\in E^\star$, 考虑在 $D(A)$ 上由
\begin{equation*}
\phi(u)=\frac12\norm{Au-f}^2+\frac12\norm{u-x}^2+\dual{Au-f}{u-x}
\end{equation*}
定义的函数 $\phi$.证明 $\phi$ 在 $D(A)$ 上是凸的, 连续的.证明当 $\norm{u}_{D(A)}\to\infty$
时 $\phi(u)\to+\infty$.

3. 由此推出存在某个 $u_0\in D(A)$ 使得对一切 $u\in D(A)$, $\phi(u_0)\le\phi(u)$.
通过选取 $u=u_0+tv$ ($v\in D(A),t>0$) 并令 $t\to0$, 能得到关于 $A$ 与 $A^\star$ 的
什么方程? \\
{[提示: 应用问题 13 第 A 部分的结果.]}

4. 证明 (M$^\star$) $\Rightarrow$ (P).

5. 由此推出 $A^\star$ 也满足性质 (P).
\end{problem}

\begin{problem}[第 1, 3, 4 章]
\label{prob:17}
\runin{-- A --}

设 $E$ 为自反 Banach 空间, $M$ 为 $E$ 的闭线性子空间.设 $C$ 为 $E^\star$ 的凸子集.
对每个 $u\in E$ 令
\begin{equation*}
\phi(u)=\sup_{g\in C}\dual{g}{u}.
\end{equation*}

1. 证明对每个 $f\in M^\perp+C$, 对一切 $u\in M$ 有 $\phi(u)\ge\dual{f}{u}$.\\
{[提示: 先从 $f\in M^\perp+C$ 的情形开始.]}

2. 反之, 设 $f\in E^\star$ 满足对一切 $u\in M$, $\phi(u)\ge\dual{f}{u}$.证明
$f\in M^\perp+C$.\\
{[提示: 使用 Hahn--Banach 定理.]}

3. 假设 $C$ 是闭且有界的, 证明 $M^\perp+C$ 是闭的.

\runin{-- B --}

在本节中假设 $E=L^p(\Omega)$ ($1<p<\infty$),
\begin{equation*}
M=\left\{u\in L^p(\Omega);\ \int_\Omega ju=0\right\},
\end{equation*}
以及
\begin{equation*}
C=\{g\in L^{p'}(\Omega);\ \abs{g(x)}\le k(x)\ \text{在}\ \Omega\ \text{上几乎处处}\},
\end{equation*}
其中 $j$ 与 $k\ge0$ 是 $L^{p'}(\Omega)$ 中的给定函数.

1. 验证 $M$ 是闭线性子空间, 且 $C$ 是凸集, 闭集和有界集.

2. 确定 $M^\perp$.

3. 对每个 $u\in L^p(\Omega)$ 确定 $\phi(u)$.

4. 由此推出若 $f\in L^{p'}(\Omega)$ 满足对一切 $u\in M$, $\int_\Omega k\abs{u}\ge\int_\Omega fu$,
则存在常数 $\lambda\in\R$ 与函数 $g\in C$ 使得 $f=\lambda j+g$.

5. 证明逆命题也成立.

\runin{-- C --}

设 $M\subset L^1(\Omega)$ 为线性子空间.设 $f,g\in L^\infty(\Omega)$ 满足在 $\Omega$
上几乎处处 $f\le g$.证明下述性质等价:
\begin{enumerate}[label=(\roman*), itemsep=0.1em]
	\item 存在 $\phi\in M^\perp$ 使得在 $\Omega$ 上几乎处处 $f\le\phi\le g$,
	\item 对一切 $u\in M$, $\int_\Omega(fu^+-gu^-)\le0$,
\end{enumerate}
其中 $u^+=\max\{u,0\}$, $u^-=\max\{-u,0\}$.\\
{[提示: 假设 (ii), 验证对一切 $u\in M$, $\int_\Omega(g+f)u\le\int_\Omega(g-f)\abs{u}$,
并应用定理  \ref{thm:1.12} 找到满足 $\abs{\psi}\le g-f$ 的某个 $\psi\in L^\infty(\Omega)$
使得 $\psi-(g+f)\in M^\perp$.取 $\phi=\frac12(g+f-\psi)$.]}
\end{problem}

\begin{problem}[第 3, 4 章]
\label{prob:18}
设 $\Omega$ 为有限测度的测度空间.设 $1<p<\infty$.设 $g:\R\to\R$ 为连续单调不减函数,
满足对某个常数 $C$, 对一切 $t\in\R$, $\abs{g(t)}\le C(\abs{t}^{p-1}+1)$.令
$G(t)=\int_0^tg(s)ds$.

1. 验证对每个 $u\in L^p(\Omega)$, 有 $g(u)\in L^{p'}(\Omega)$ 且 $G(u)\in L^1(\Omega)$.

设 $(u_n)$ 为 $L^p(\Omega)$ 中的序列, $u\in L^p(\Omega)$ 满足:
(i) $u_n\rightharpoonup u$ 弱收敛于 $\sigma(L^p,L^{p'})$, 且
(ii) $\limsup_{n\to\infty}\int_\Omega G(u_n)\le\int_\Omega G(u)$.
本问题的目的是建立下述性质:
\begin{equation}
\text{对每个}\ q\in[1,p'),\ g(u_n)\to g(u)\ \text{在}\ L^q\ \text{中强收敛},
\tag{1}
\label{eq:prob18.1}
\end{equation}
\begin{equation}
\text{若此外假设}\ g\ \text{严格递增, 则对每个}\ q\in[1,p),\ u_n\to u\ \text{在}\ L^q\ \text{中强收敛}.
\tag{2}
\label{eq:prob18.2}
\end{equation}

2. 验证对一切 $a,b\in\R$, $G(a)-G(b)-g(b)(a-b)\ge0$.若 $G(a)-G(b)-g(b)(a-b)=0$,
能得出什么结论?

3. 设 $(a_n)$ 为 $\R$ 中的序列, $b\in\R$ 使得
$\lim_{n\to\infty}[G(a_n)-G(b)-g(b)(a_n-b)]=0$.证明 $g(a_n)\to g(b)$.

4. 证明 $\int_\Omega\abs{G(u_n)-G(u)-g(u)(u_n-u)}\to0$.由此推出存在子序列 $(u_{n_k})$
使得在 $\Omega$ 上几乎处处 $G(u_{n_k})-G(u)-g(u)(u_{n_k}-u)\to0$, 从而在 $\Omega$
上几乎处处 $g(u_{n_k})\to g(u)$.

5. 证明 \eqref{eq:prob18.1} 成立 (对整个序列验证 \eqref{eq:prob18.1}, 而不仅对子序列).

6. 证明 \eqref{eq:prob18.2} 成立.

在下文中, 我们额外假设存在常数 $\alpha>0$ 与 $C$ 使得
\begin{equation}
\abs{g(t)}\ge\alpha\abs{t}^{p-1}-C\quad\forall t\in\R.
\tag{3}
\label{eq:prob18.3}
\end{equation}

7. 证明在 $L^{p'}$ 中 $g(u_n)\to g(u)$ 强收敛.

8. 没有假设 \eqref{eq:prob18.3} 能得出同样的结论吗?

9. 若此外 $g$ 严格递增, 证明在 $L^p$ 中 $u_n\to u$ 强收敛.
\end{problem}

\begin{problem}[第 3, 4 章]
\label{prob:19}
设 $E$ 为配以范数 $\norm{u}_E=\norm{u}_1+\norm{u}_2$ 的空间 $L^1(\R)\cap L^2(\R)$.

1. 验证 $E$ 是 Banach 空间.设 $f(x)=f_1(x)+f_2(x)$, 其中 $f_1\in L^\infty(\R)$
且 $f_2\in L^2(\R)$.验证映射 $u\mapsto\int_\R f(x)u(x)dx$ 是 $E$ 上的连续线性泛函.

2. 设 $0<\alpha<1/2$; 验证映射 $u\mapsto\int_\R\dfrac{1}{\abs{x}^\alpha}u(x)dx$
是 $E$ 上的连续线性泛函.\\
{[提示: 把积分拆成两部分: $[\abs{x}>M]$ 与 $[\abs{x}\le M]$.]}

3. 令
\begin{equation*}
K=\left\{u\in E;\ u\ge0\ \text{几乎处处在}\ \R\ \text{上且}\ \int_\R u(x)dx\le1\right\}.
\end{equation*}
验证 $K$ 是 $E$ 的闭凸子集.

4. 设 $(u_n)$ 为 $K$ 中的序列, 设 $u\in K$ 使得在 $L^2(\R)$ 中 $u_n\rightharpoonup u$
弱收敛.验证 $u\in K$ 并证明
\begin{equation*}
\int_\R\frac{1}{\abs{x}^\alpha}u_n(x)dx\to\int_\R\frac{1}{\abs{x}^\alpha}u(x)dx.
\end{equation*}

考虑对每个 $u\in E$ 由 $J(u)=\int_\R u^2(x)dx-\int_\R\dfrac{1}{\abs{x}^\alpha}u(x)dx$
定义的函数 $J$.

5. 验证存在常数 $C$ 使得对一切 $u\in K$, $J(u)\ge C$.我们断言 $m=\inf_{u\in K}J(u)$
可以达到.

6. 设 $(u_n)$ 为 $K$ 中的序列使得 $J(u_n)\to m$.证明 $\norm{u_n}_E$ 有界.

7. 设 $(u_{n_k})$ 为使得在 $L^2(\R)$ 中 $u_{n_k}\rightharpoonup u$ 弱收敛的子序列.
证明 $J(u)=m$.

8. $E$ 是自反空间吗?
\end{problem}

\begin{problem}[第~
\ref{chap:4}~章]
\label{prob:20}
\runin{Clarkson 不等式. $L^p$ 的一致凸性}

\runin{-- A --}

在本部分中假设 $2\le p<\infty$, 我们将建立下述不等式:
\begin{equation}
\abs{x+y}^p+\abs{x-y}^p\le2(\abs{x}^{p'}+\abs{y}^{p'})^{p/p'}\quad\forall x,y\in\R,
\label{eq:prob20.1}
\end{equation}
\begin{equation}
2(\abs{x}^{p'}+\abs{y}^{p'})^{p/p'}\le2^{p-1}(\abs{x}^p+\abs{y}^p)\quad\forall x,y\in\R.
\label{eq:prob20.2}
\end{equation}

1. 证明 \eqref{eq:prob20.2}.\\
{[提示: 利用函数 $g(t)=\abs{t}^{p/p'}$ 的凸性.]}

2. 令 $f(x)=(1+x^{1/p})^p+(1-x^{1/p})^p$, $x\in(0,1)$.证明对一切 $x\in(0,1)$,
$f''(x)\le0$.由此推出
\begin{equation}
f(x)\le f(y)+(x-y)f'(y)\quad\forall x,y\in(0,1).
\label{eq:prob20.3}
\end{equation}

3. 证明对一切 $x\in(0,1)$, $f(x)\le2(1+x^{p'/p})^{p/p'}$.\\
{[提示: 在 \eqref{eq:prob20.3} 中取 $y=x^{p'}$.]}

4. 推出 \eqref{eq:prob20.1}.

在下文中, $\Omega$ 表示 $\sigma$-有限测度空间.

\runin{-- B --}

在本部分中再次假设 $2\le p<\infty$.

1. 证明下述不等式:
\begin{equation}
\norm{f+g}_p^p+\norm{f-g}_p^p\le2(\norm{f}_p^{p'}+\norm{g}_p^{p'})^{p/p'}\quad\forall f,g\in L^p(\Omega),
\label{eq:prob20.4}
\end{equation}
\begin{equation}
2(\norm{f}_p^{p'}+\norm{g}_p^{p'})^{p/p'}\le2^{p-1}(\norm{f}_p^p+\norm{g}_p^p)\quad\forall f,g\in L^p(\Omega).
\label{eq:prob20.5}
\end{equation}

2. 推出 Clarkson 第一不等式 (见定理  \ref{thm:4.10}).

\runin{-- C --}

在本部分中假设 $1<p\le2$.

1. 建立下述不等式:
\begin{equation}
\norm{f+g}_p^{p'}+\norm{f-g}_p^{p'}\le2(\norm{f}_p^p+\norm{g}_p^p)^{p'/p}\quad\forall f,g\in L^p(\Omega).
\label{eq:prob20.6}
\end{equation}
不等式 \eqref{eq:prob20.6} 称为 Clarkson 第二不等式.\\
{[提示: 有两种不同的方法:
(i) 由 \eqref{eq:prob20.4} 通过对偶性, 注意到
$\sup_{\phi,\psi\in L^{p'}}\dfrac{\int(u\phi+v\psi)}{[\norm{\phi}_{p'}^p+\norm{\psi}_{p'}^p]^{1/p}}=(\norm{u}_p^{p'}+\norm{v}_p^{p'})^{1/p'}$.
(ii) 直接由 \eqref{eq:prob20.1} 结合习题 4.11 的结果.]}

2. 由此推出对 $1<p\le2$, $L^p(\Omega)$ 是一致凸的.
\end{problem}
\begin{problem}[第~
\ref{chap:4}~章]
\index{Marcinkiewiczkongjian@Marcinkiewicz 空间 (Marcinkiewicz space)}
\index{fenbuhanshu@分布函数 (distribution function)}
\label{prob:21}
\runin{分布函数. Marcinkiewicz 空间}
在整个本题中, $\Omega$ 表示具有有限测度 $\mu$ 的测度空间.给定可测函数
$f:\Omega\to\R$, 我们定义其\emph{分布函数} $\alpha$ 为
\begin{equation*}
\alpha(t)=\abs{[\abs{f}>t]}=\operatorname{meas}\{x\in\Omega;\ \abs{f(x)}>t\}\quad\forall t\ge0.
\end{equation*}

\runin{-- A --}

1. 验证 $\alpha$ 单调不增.证明对一切 $t\ge0$, $\alpha(t+0)=\alpha(t)$.构造一个简单的例子,
使得对某个 $t>0$, $\alpha(t-0)\ne\alpha(t)$.

2. 设 $(f_n)$ 为可测函数序列, 使得在 $\Omega$ 上几乎处处 $f_n\to f$.用 $(\alpha_n)$ 与
$\alpha$ 记相应的分布函数.证明对一切 $t>0$,
\begin{equation*}
\alpha(t)\le\liminf_{n\to\infty}\alpha_n(t)\le\limsup_{n\to\infty}\alpha_n(t)\le\alpha(t-0).
\end{equation*}
由此推出 $\alpha_n(t)\to\alpha(t)$ 几乎处处成立.

\runin{-- B --}

1. 设 $g\in L^1_{\text{loc}}(\R)$ 为满足几乎处处 $g\ge0$ 的函数.令 $G(t)=\int_0^tg(s)ds$.
证明对每个可测函数 $f$,
\begin{equation*}
\int_\Omega G(\abs{f(x)})d\mu<\infty\iff\int_0^\infty\alpha(t)g(t)dt<\infty,
\end{equation*}
且
\begin{equation*}
\int_\Omega G(\abs{f(x)})d\mu=\int_0^\infty\alpha(t)g(t)dt.
\end{equation*}
[提示: 使用 Fubini 与 Tonelli 定理.]

2. 更一般地, 证明对一切 $\lambda\ge0$,
\begin{equation*}
\int_{[\abs{f}>\lambda]}G(\abs{f(x)})d\mu=\alpha(\lambda)G(\lambda)+\int_\lambda^\infty\alpha(t)g(t)dt.
\end{equation*}

3. 由此推出对 $1\le p<\infty$,
\begin{equation*}
f\in L^p(\Omega)\iff\int_0^\infty\alpha(t)t^{p-1}dt<\infty,
\end{equation*}
且对一切 $\lambda\ge0$,
\begin{equation*}
\int_{[\abs{f}>\lambda]}\abs{f(x)}^pd\mu=\alpha(\lambda)\lambda^p+p\int_\lambda^\infty\alpha(t)t^{p-1}dt.
\end{equation*}
验证若 $f\in L^p(\Omega)$, 则 $\lim_{t\to+\infty}\alpha(t)t^p=0$.

\runin{-- C --}

设 $1<p<\infty$.对每个 $f\in L^1(\Omega)$, 定义
\begin{equation*}
[f]_p=\sup\left\{\abs{A}^{-1/p'}\int_A\abs{f};\ A\subset\Omega\ \text{可测},\ \abs{A}>0\right\}\le\infty,
\end{equation*}
并考虑集合 $M^p(\Omega)=\{f\in L^1(\Omega);\ [f]_p<\infty\}$, 称为 $p$ 阶
\emph{Marcinkiewicz 空间}.空间 $M^p$ 也称为\emph{弱 $L^p$ 空间}, 但这个术语
容易引起混淆, 因为``弱''字已经用于弱拓扑.

1. 验证 $M^p(\Omega)$ 是线性空间, 且 $[\cdot]_p$ 是范数.证明 $L^p(\Omega)\subset M^p(\Omega)$,
且对每个 $f\in L^p(\Omega)$, $[f]_p\le\norm{f}_p$.

2. 证明配以范数 $[\cdot]_p$ 的 $M^p(\Omega)$ 是 Banach 空间.验证对 $1<q\le p$,
$M^p(\Omega)\subset M^q(\Omega)$ 连续嵌入.

我们断言: $f\in M^p(\Omega)$ 当且仅当 $f$ 可测且 $\sup_{t>0}t^p\alpha(t)<\infty$.

3. 证明若 $f\in M^p(\Omega)$, 则对一切 $t>0$, $t^p\alpha(t)\le[f]_p^p$.

4. 反之, 设 $f$ 为满足 $\sup_{t>0}t^p\alpha(t)<\infty$ 的可测函数.证明存在常数 $C_p$
(只依赖于 $p$) 使得
\begin{equation*}
[f]_p^p\le C_p\sup_{t>0}t^p\alpha(t).
\end{equation*}
[提示: 使用第 B3 问并写出 $\int_A\abs{f}=\int_{A\cap[\abs{f}>\lambda]}\abs{f}+\int_{A\cap[\abs{f}\le\lambda]}\abs{f}$; 然后变动 $\lambda$.]

5. 证明对 $1\le q<p$, $M^p(\Omega)\subset L^q(\Omega)$ 连续嵌入.

6. 设 $1<q<r<\infty$ 与 $\theta\in(0,1)$; 令 $\dfrac1p=\dfrac\theta q+\dfrac{1-\theta}{r}$.
证明存在只依赖于 $q,r$ 与 $\theta$ 的常数 $C$ 使得对一切 $f\in M^r(\Omega)$,
\begin{equation*}
\norm{f}_p\le C[f]_q^\theta[f]_r^{1-\theta}.
\end{equation*}

7. 设 $\Omega=\{x\in\R^N;\ \abs{x}<1\}$ 配以 Lebesgue 测度, 设 $f(x)=\abs{x}^{-N/p}$
($1<p<\infty$).验证 $f\in M^p(\Omega)$, 而 $f\notin L^p(\Omega)$.
\end{problem}

\begin{problem}[第~
\ref{chap:4}~章]
\label{prob:22}
\runin{一个插值定理 (Schur, Riesz, Thorin, Marcinkiewicz)}
设 $\Omega$ 为有限测度的测度空间.设 $T:L^1(\Omega)\to L^1(\Omega)$ 为有界线性算子,
其范数记作 $N_1=\norm{T}_{\LL(L^1,L^1)}$.我们假设 $T(L^\infty(\Omega))\subset L^\infty(\Omega)$.

1. 证明 $T$ 是由 $L^\infty(\Omega)$ 到自身中的有界算子.令 $N_\infty=\norm{T}_{\LL(L^\infty,L^\infty)}$.

本问题的目标是证明对每个 $1<p<\infty$,
\begin{equation*}
T(L^p(\Omega))\subset L^p(\Omega),
\end{equation*}
且 $T:L^p(\Omega)\to L^p(\Omega)$ 是有界算子, 其范数 $N_p=\norm{T}_{\LL(L^p,L^p)}$
满足不等式 $N_p\le2N_1^{1/p}N_\infty^{1/p'}$.

为简单起见, 先设 $N_\infty=1$.给定函数 $u\in L^1(\Omega)$, 对每个 $\lambda>0$, 我们令
$u=v_\lambda+w_\lambda$, 其中 $v_\lambda=u\chi_{[\abs{u}>\lambda]}$ 且
$w_\lambda=u\chi_{[\abs{u}\le\lambda]}$, $f=Tu$, $g_\lambda=Tv_\lambda$, 以及
$h_\lambda=Tw_\lambda$, 从而 $f=g_\lambda+h_\lambda$.

2. 验证对一切 $\lambda>0$, $\norm{g_\lambda}_1\le N_1\int_{[\abs{u}>\lambda]}\abs{u(x)}d\mu$,
且 $\norm{h_\lambda}_\infty\le\lambda$.

3. 考虑分布函数 $\alpha(t)=\abs{[\abs{u}>t]}$, $\beta(t)=\abs{[\abs{f}>t]}$, 以及
$\gamma_\lambda(t)=\abs{[\abs{g_\lambda}>t]}$.证明对一切 $\lambda>0$,
\begin{equation*}
\int_0^\infty\gamma_\lambda(t)dt\le N_1\left[\alpha(\lambda)\lambda+\int_\lambda^\infty\alpha(t)dt\right],
\end{equation*}
且对一切 $\lambda>0$ 及一切 $t>\lambda$, $\beta(t)\le\gamma_\lambda(t-\lambda)$.\\
{[提示: 应用问题 21 第 B 部分的结果.]}

4. 假设 $u\in L^p(\Omega)$, 证明 $f\in L^p(\Omega)$ 且 $\norm{f}_p\le2N_1^{1/p}\norm{u}_p$.

5. 在 $N_\infty\ne1$ 的一般情形下得出结论.

\runin{注记.} 事实上通过不同的论证可以证明 $N_p\le N_1^{1/p}N_\infty^{1/p'}$;
参见例如 Bergh--Löfström \bref{BerghLofstrom}{1} 及其书第~
\ref{chap:1}~章注记中的文献.
\end{problem}

\begin{problem}[第 3, 4 章]
\label{prob:23}
\runin{$L^1$ 的弱紧子集与等度可积族. Hahn--Vitali--Saks, Dunford--Pettis 与 de la Vallée-Poussin 定理}
设 $\Omega$ 为 $\sigma$-有限测度空间.回忆 (见习题 4.36), 若子集 $\mathcal{F}\subset L^1(\Omega)$
满足下述性质, 则称它是\emph{等度可积}的:
\begin{enumerate}[label=(\alph*), itemsep=0.1em]
	\item $\mathcal{F}$ 在 $L^1(\Omega)$ 中有界,
	\item 对一切 $\varepsilon>0$, 存在 $\delta>0$ 使得对满足 $\abs{A}<\delta$ 的一切可测集 $A\subset\Omega$ 及一切 $f\in\mathcal{F}$, $\int_A\abs{f}<\varepsilon$,
	\item 对一切 $\varepsilon>0$, 存在满足 $\abs{\omega}<\infty$ 的可测集 $\omega\subset\Omega$ 使得对一切 $f\in\mathcal{F}$, $\int_{\Omega\setminus\omega}\abs{f}<\varepsilon$.
\end{enumerate}
本问题的第一个目标是建立给定集合 $\mathcal{F}\subset L^1(\Omega)$ 的下述性质等价:
\begin{enumerate}[label=(\roman*), itemsep=0.1em]
	\item $\mathcal{F}$ 包含在 $L^1(\Omega)$ 的弱 (关于 $\sigma(L^1,L^\infty)$) 紧子集中,
	\item $\mathcal{F}$ 是等度可积的.
\end{enumerate}

\runin{-- A --}

蕴含关系 (i) $\Rightarrow$ (ii).

1. 设 $(f_n)$ 为 $L^1(\Omega)$ 中的序列, 使得对满足 $\abs{A}<\infty$ 的一切可测集
$A\subset\Omega$, $\int_A f_n\to0$.证明 $(f_n)$ 满足性质 (b).\\
{[提示: 考虑由 $X=\{\chi_A;\ A\subset\Omega\ \text{可测且}\ \abs{A}<\infty\}$ 定义的
子集 $X\subset L^1(\Omega)$.验证 $X$ 在 $L^1(\Omega)$ 中闭, 并对序列
$X_n=\{\chi_A\in X;\ \abs{\int_A f_k}\le\varepsilon\ \forall k\ge n\}$ 应用 Baire 纲定理,
其中 $\varepsilon>0$ 固定.]}

2. 设 $(f_n)$ 为 $L^1(\Omega)$ 中的序列, 使得对满足 $\abs{A}\le\infty$ 的一切可测集
$A\subset\Omega$, $\int_A f_n\to0$.证明 $(f_n)$ 满足性质 (c).\\
{[提示: 设 $(\Omega_i)$ 为满足 $\Omega=\bigcup_i\Omega_i$ 的测度有限的可测集的单调递增序列.
在 $L^\infty(\Omega)$ 上考虑由 $d(f,g)=\sum_i\frac{1}{2^i\abs{\Omega_i}}\int_{\Omega_i}\abs{f-g}$
定义的度量 $d$.令 $Y=\{\chi_A;\ A\subset\Omega\ \text{可测}\}$.验证 $Y$ 关于度量 $d$
完备, 并对序列 $Y_n=\{\chi_A\in Y;\ \abs{\int_A f_k}\le\varepsilon\ \forall k\ge n\}$
应用 Baire 纲定理.]}

3. 由此推出: 若 $(f_n)$ 是 $L^1(\Omega)$ 中的序列且 $f_n\rightharpoonup f$ 弱收敛于
$\sigma(L^1,L^\infty)$, 则 $(f_n)$ 是等度可积的.

4. 证明 (i) $\Rightarrow$ (ii).\\
{[提示: 用反证法并应用 Eberlein--Šmulian 定理; 见问题 10.]}

5. 仅假设对满足 $\abs{A}<\infty$ (相应地 $\abs{A}\le\infty$) 的每个可测集 $A\subset\Omega$,
$\int_A f_n$ 收敛于有限极限 $\ell(A)$, 重新讨论第 1 问 (相应地第 2 问).

\runin{-- B --}

蕴含关系 (ii) $\Rightarrow$ (i).

1. 设 $E$ 为 Banach 空间, $F\subset E$.假设对一切 $\varepsilon>0$, 存在弱
(关于 $\sigma(E,E^\star)$) 紧集 $F_\varepsilon\subset E$ 使得
$F\subset F_\varepsilon+\varepsilon B_E$.证明 $F$ 包含在 $E$ 的弱紧子集中.\\
{[提示: 考虑 $G=\overline F^{\sigma(E^{\star\star},E^\star)}$.]}

2. 由此推出 (ii) $\Rightarrow$ (i).\\
{[提示: 考虑族 $(\chi_\omega T_nf)_{f\in\mathcal{F}}$, 其中 $\abs{\omega}<\infty$,
$T_n$ 是定理  \ref{thm:4.12} 证明中的截断算子.]}

\runin{-- C --}

若干应用.

1. 设 $(f_n)$ 为 $L^1(\Omega)$ 中的序列使得 $f_n\rightharpoonup f$ 弱收敛于
$\sigma(L^1,L^\infty)$ 且 $f_n\to f$ 几乎处处.证明 $\norm{f_n-f}_1\to0$.\\
{[提示: 应用习题 4.14.]}

2. 设 $u_1,u_2\in L^1(\Omega)$ 满足几乎处处 $u_1\le u_2$.证明集合
$K=\{f\in L^1(\Omega);\ u_1\le f\le u_2\ \text{几乎处处}\}$ 在弱拓扑
$\sigma(L^1,L^\infty)$ 中是紧的.

3. 设 $(f_n)$ 为 $L^1(\Omega)$ 中的等度可积序列.证明存在弱收敛于 $\sigma(L^1,L^\infty)$
的子序列 $(f_{n_k})$.

4. 设 $(f_n)$ 为 $L^1(\Omega)$ 中的有界序列使得对每个可测集 $A\subset\Omega$,
$\int_A f_n$ 收敛于有限极限 $\ell(A)$.证明存在某个 $f\in L^1(\Omega)$ 使得
$f_n\rightharpoonup f$ 弱收敛于 $\sigma(L^1,L^\infty)$.

5. 设 $g:\R\to\R$ 为满足对一切 $t\in\R$ 有 $\abs{g(t)}\le C$ 的连续单调递增函数.
令 $G(t)=\int_0^tg(s)ds$.设 $(f_n)$ 为 $L^1(\Omega)$ 中的序列使得
$f_n\rightharpoonup f$ 弱收敛于 $\sigma(L^1,L^\infty)$ 且
$\limsup\int_\Omega G(f_n)\le\int_\Omega G(f)$.证明 $\norm{f_n-f}_1\to0$.\\
{[提示: 参看问题 18.]}

\runin{-- D --}

在本部分中假设 $\abs{\Omega}<\infty$.设 $\mathcal{F}\subset L^1(\Omega)$.

1. 设 $G:[0,+\infty)\to[0,+\infty)$ 为满足 $\lim_{t\to+\infty}G(t)/t=+\infty$
的连续函数.假设存在常数 $C$ 使得对一切 $f\in\mathcal{F}$, $\int_\Omega G(\abs{f})\le C$.
证明 $\mathcal{F}$ 是等度可积的.

2. 反之, 假设 $\mathcal{F}$ 是等度可积的.证明存在凸单调递增函数
$G:[0,+\infty)\to[0,+\infty)$ 满足 $\lim_{t\to+\infty}G(t)/t=+\infty$, 使得对某个
常数 $C$, 对一切 $f\in\mathcal{F}$ 有 $\int_\Omega G(\abs{f})\le C$.\\
{[提示: 使用分布函数; 见问题 21.]}
\end{problem}

\begin{problem}[第 1, 3, 4 章]
\index{Radonceduo@Radon 测度 (Radon measure)}
\label{prob:24}
\runin{Radon 测度}
设 $K$ 为带有距离 $d$ 的紧度量空间, 设 $E=C(K)$ 配以通常的范数
$\norm{f}=\max_{x\in K}\abs{f(x)}$.对偶空间 $E^\star$ 记作 $\mathcal{M}(K)$,
称为 $K$ 上的 \emph{Radon 测度空间}.空间 $\mathcal{M}(K)$ 配以对偶范数, 记作
$\norm{\cdot}_{\mathcal{M}}$ 或简记为 $\norm{\cdot}$.本问题的目的是介绍
$\mathcal{M}(K)$ 的若干性质.

\runin{-- A --}

这里我们证明 $C(K)$ 可分.给定 $\delta>0$, 设 $\bigcup_{j\in J}B(a_j,\delta/2)$ 为
$K$ 的有限覆盖.令
\begin{align*}
q_j(x)&=\max\{0,\delta-d(x,a_j)\},\quad j\in J,\ x\in K,\\
q(x)&=\sum_{j\in J}q_j(x).
\end{align*}

1. 验证函数 $(q_j)_{j\in J}$ 与 $q$ 在 $K$ 上连续.证明对一切 $x\in K$,
$q(x)\ge\delta/2$.

2. 令 $\theta_j(x)=\dfrac{q_j(x)}{q(x)}$ ($j\in J,x\in K$).证明函数
$(\theta_j)_{j\in J}$ 在 $K$ 上连续, $[\theta_j(x)\ne0]\iff[d(x,a_j)<\delta]$,
且对一切 $x\in K$, $\sum_{j\in J}\theta_j(x)=1$.函数族 $(\theta_j)_{j\in J}$
称为 (从属于开覆盖 $\bigcup_{j\in J}B(a_j,2\delta)$ 的) \emph{单位分解}.

3. 给定 $f\in C(K)$, 令 $\tilde f(x)=\sum_{j\in J}f(a_j)\theta_j(x)$.证明
\begin{equation*}
\norm{f-\tilde f}\le\sup_{\substack{x,y\in K\\d(x,y)<\delta}}\abs{f(x)-f(y)}.
\end{equation*}

4. 选取 $\delta=1/n$, 上述构造给出一个有限集 $J_n$, 点的有限族 $(a_j)_{j\in J_n}$,
以及函数的有限族 $(\theta_j)_{j\in J_n}$.证明由函数 $(\theta_j)$ ($j\in J_n,n=1,2,3,\dots$)
张成的向量空间在 $C(K)$ 中稠密.

5. 由此推出 $C(K)$ 可分.

\runin{-- B --}

在本部分中假设 $K=\overline\Omega$, 其中 $\Omega$ 是 $\R^N$ 中的有界开集.通过
由对一切 $u\in L^1(\Omega)$ 与一切 $f\in C(\overline\Omega)$ 由
$\dual{Tu}{f}=\int_\Omega uf$ 定义的嵌入 $T:L^1(\Omega)\to\mathcal{M}(\overline\Omega)$,
把 $L^1(\Omega)$ 视为 $\mathcal{M}(\overline\Omega)$ 的子空间是方便的.

1. 验证对一切 $u\in L^1(\Omega)$, $\norm{Tu}_{\mathcal{M}}=\norm{u}_{L^1}$.\\
{[提示: 使用习题 4.26.]}

2. 设 $(v_n)$ 为 $L^1(\Omega)$ 中的有界序列.证明存在子序列 $(v_{n_k})$ 与某个
$\mu\in\mathcal{M}(\overline\Omega)$ 使得在 $\mathcal{M}(\overline\Omega)$ 中
$v_{n_k}\overset{*}{\rightharpoonup}\mu$ 且
$\norm{\mu}_{\mathcal{M}}\le\liminf_{k\to\infty}\norm{v_{n_k}}_{L^1}$.\\
{[提示: 使用推论  \ref{cor:3.30}.]}

现在的目标是证明: 给定任何 $\mu_0\in\mathcal{M}(\overline\Omega)$, 存在 $C_c^\infty(\Omega)$
中的序列 $(u_n)$ 使得
\begin{equation}
\int_\Omega u_nf\to\dual{\mu_0}{f}\quad\forall f\in C(\overline\Omega)
\label{eq:prob24.1}
\end{equation}
且
\begin{equation}
\norm{u_n}_{L^1}=\norm{\mu_0}_{\mathcal{M}}\quad\forall n.
\label{eq:prob24.2}
\end{equation}
不失一般性可以假设 $\norm{\mu_0}_{\mathcal{M}}=1$.令
$A=\{u\in C_c^\infty(\Omega);\ \norm{u}_{L^1}\le1\}$.

3. 证明 $\mu_0\in\overline A^{\sigma(E^\star,E)}$.\\
{[提示: 在 $E^\star$ 中对弱$^\star$拓扑 $\sigma(E^\star,E)$ 应用 Hahn--Banach 定理;
见问题 9.然后使用推论  \ref{cor:4.23}.]}

4. 由此推出存在 $A$ 中的序列 $(v_n)$ 使得在 $\sigma(E^\star,E)$ 中
$v_n\overset{*}{\rightharpoonup}\mu_0$.验证 $\lim_{n\to\infty}\norm{v_n}_{L^1}=1$.

5. 得出序列 $u_n=v_n/\norm{v_n}_{L^1}$ 满足 \eqref{eq:prob24.1} 与 \eqref{eq:prob24.2}.

若对满足 $f\ge0$ 的一切 $f\in C(\overline\Omega)$ 有 $\dual{\mu}{f}\ge0$, 则称
$\mu\ge0$.

6. 验证若 $\mu\ge0$, 则 $\dual{\mu}{1}=\norm{\mu}$, 其中 $1$ 表示恒等于 $1$ 的函数.

7. 假设 $\mu_0\in\mathcal{M}(\overline\Omega)$ 满足 $\mu_0\ge0$ 且 $\norm{\mu_0}=1$
(这样的测度称为\emph{概率测度}).在 $C_c^\infty(\Omega)$ 中构造满足 \eqref{eq:prob24.1}, \eqref{eq:prob24.2} 且此外对一切 $n$ 与一切 $x\in\Omega$ 满足 $u_n(x)\ge0$ 的序列 $(u_n)$.

8. 计算 $\norm{u+\delta_a}_{\mathcal{M}}$, 其中 $u\in L^1$, $\delta_a$ ($a\in\Omega$)
在下文定义.

\runin{-- C --}

现在回到一般框架, 用 $\delta_a$ 记点 $a\in K$ 处的 Dirac 质量, 即由
$\dual{\delta_a}{f}=f(a)$ ($\forall f\in C(K)$) 定义的测度.令
\begin{equation*}
\mathcal{D}=\left\{\mu=\sum_{j\in J}\alpha_j\delta_{a_j};\ J\ \text{有限},\ \alpha_j\in\R,\ \text{且点}\ a_j\ \text{互不相同}\right\}.
\end{equation*}

1. 证明若 $\mu\in\mathcal{D}$, 则 $\norm{\mu}=\sum_{j\in J}\abs{\alpha_j}$,
且 $[\mu\ge0]\iff[\alpha_j\ge0\ \forall j]$.令 $\mathcal{D}_1=\{\mu\in\mathcal{D};\ \norm{\mu}\le1\}$.

2. 证明满足 $\norm{\mu_0}\le1$ 的任何测度 $\mu_0\in\mathcal{M}(K)$ 属于
$\overline{\mathcal{D}_1}^{\sigma(E^\star,E)}$.\\
{[提示: 使用与第 B 部分相同的技巧.]}

3. 由此推出对给定的任何测度 $\mu_0\in\mathcal{M}(\overline\Omega)$, 存在 $\mathcal{D}$
中的序列 $(\nu_n)$ 使得 $\nu_n\overset{*}{\rightharpoonup}\mu_0$ 且对一切 $n$,
$\norm{\nu_n}=\norm{\mu_0}$.

4. 设 $\mu_0$ 为概率测度.证明存在 $\mathcal{D}$ 中的概率测度序列 $(\nu_n)$
使得 $\nu_n\overset{*}{\rightharpoonup}\mu_0$.

\runin{-- D --}

本部分的目的是证明每个 $\mu\in\mathcal{M}(K)$ 允许唯一的分解 $\mu=\mu_1-\mu_2$,
其中 $\mu_1,\mu_2\in\mathcal{M}(K)$, $\mu_1,\mu_2\ge0$, 且 $\norm{\mu_1}+\norm{\mu_2}=\norm{\mu}$
(测度 $\mu_1$ 与 $\mu_2$ 常记作 $\mu^+$ 与 $\mu^-$).

给定满足 $f\ge0$ 的 $f\in C(K)$, 令
\begin{equation*}
L(f)=\sup\{\dual{\mu}{g};\ g\in C(K)\ \text{且在}\ K\ \text{上}\ 0\le g\le f\}.
\end{equation*}

1. 验证 $0\le L(f)\le\norm{\mu}\norm{f}$, 对一切 $\lambda\ge0$ 有 $L(\lambda f)=\lambda L(f)$,
且对满足 $f_1\ge0$ 与 $f_2\ge0$ 的一切 $f_1,f_2\in C(K)$, $L(f_1+f_2)=L(f_1)+L(f_2)$.

给定任何 $f\in C(K)$, 令 $\mu_1(f)=L(f^+)-L(f^-)$, 其中 $f^+=\max\{f,0\}$ 与
$f^-=\max\{-f,0\}$.

2. 证明映射 $f\mapsto\mu_1(f)$ 在 $C(K)$ 上是线性的且对一切 $f\in C(K)$ 满足
$\abs{\mu_1(f)}\le\norm{\mu}\norm{f}$, 从而 $\mu_1\in\mathcal{M}(K)$.验证 $\mu_1\ge0$.

3. 令 $\mu_2=\mu_1-\mu$ 并验证 $\mu_2\ge0$.证明 $\norm{\mu}=\norm{\mu_1}+\norm{\mu_2}$.

4. 设 $\nu\in\mathcal{M}(K)$ 满足 $\nu\ge0$ 且 $\nu\ge\mu$.证明 $\nu\ge\mu_1$.
类似地若 $\nu\in\mathcal{M}(K)$ 且 $\nu\ge-\mu$, 证明 $\nu\ge\mu_2$.由此推出分解的唯一性.

\runin{-- E --}

证明当把空间 $E=C(\overline\Omega)$ 换成子空间
$E_0=\{f\in C(\overline\Omega);\ f=0\ \text{在}\ \partial\Omega\ \text{上}\}$ 时,
上述所有结果 (除第 B6 问外) 仍然成立.$E_0$ 的对偶常记作 $\mathcal{M}_0(\Omega)$
或 $\mathcal{M}(\Omega)$.

\runin{-- F --}

重访 Dunford--Pettis 定理.设 $(f_n)$ 为 $L^1(\Omega)$ 中的序列.回忆若 $(f_n)$
满足对一切 $\varepsilon>0$, 存在 $\delta>0$ 使得对满足 $\abs{A}<\delta$ 的一切可测集
$A\subset\Omega$ 及一切 $n$, $\int_A\abs{f_n}<\varepsilon$, 则称 $(f_n)$ 是等度可积的.
目标是证明每个等度可积序列 $(f_n)$ 允许子序列 $(f_{n_j})$ 弱收敛于某个 $f\in L^1(\Omega)$.

1. 证明 $(f_n)$ 在 $L^1(\Omega)$ 中有界.

2. 验证对一切 $n,k$, $\int_\Omega\abs{f_n-T_kf_n}\le\int_{[\abs{f_n}>k]}\abs{f_n}$.

3. 由此推出对一切 $\varepsilon>0$, 存在 $k>0$ 使得对一切 $n$, $\int_\Omega\abs{f_n-T_kf_n}\le\varepsilon$.

抽取子序列 (仍记作 $(f_n)$), 可以假设在 $\mathcal{M}(\overline\Omega)$ 中
$f_n\overset{*}{\rightharpoonup}\mu$ 弱$^\star$收敛于某个测度 $\mu$.

4. 证明对一切 $\varepsilon>0$, 存在 $g=g_\varepsilon\in L^\infty(\Omega)$ 使得
$\norm{\mu-g_\varepsilon}_{\mathcal{M}}\le\varepsilon$.

5. 由此推出 $\mu\in L^1(\Omega)$.

6. 证明 $f_n\rightharpoonup\mu$ 弱收敛于 $\sigma(L^1,L^\infty)$.
\end{problem}

\begin{problem}[第 1, 5 章]
\label{prob:25}
设 $H$ 为 Hilbert 空间, $C\subset H$ 为以 $0$ 为顶点的凸锥, 即 $0\in C$, 且对一切
$\lambda,\mu>0$ 与一切 $u,v\in C$, $\lambda u+\mu v\in C$.我们假设 $C$ 非空, 开,
且 $C\ne H$.

验证 $0\notin C$ 且 $0\in\overline C$.考虑集合
\begin{equation*}
\Gamma=\{u\in H;\ (u,v)\le0\ \forall v\in C\}.
\end{equation*}

1. 验证 $\Gamma$ 是以 $0$ 为顶点的凸锥, $\Gamma$ 闭, 且 $0\in\Gamma$.证明
$C=\{v\in H;\ (u,v)<0\ \forall u\in\Gamma\setminus\{0\}\}$, 并由此推出 $\Gamma$
不退化为 $\{0\}$.\\
{[提示: 使用 Hahn--Banach 定理.]}

2. 设 $\omega\in C$ 固定, 考虑集合 $K=\{u\in\Gamma;\ (u,\omega)=-1\}$.证明 $K$
是非空, 有界, 闭的凸集, 使得 $0\notin K$ 且
$\Gamma\setminus\{0\}=\bigcup_{\lambda>0}\lambda K$.画出示意图.\\
{[提示: 考虑以 $\omega$ 为中心, 包含在 $C$ 内的半径为 $\rho>0$ 的球.]}

3. 设 $a=P_K0$.证明 $a\in(-C)\cap\Gamma$.

4. 用简捷的方法直接证明 $(-C)\cap\Gamma\ne\emptyset$.

5. 设 $D\subset H$ 为非空开凸集, $x_0\notin D$.证明存在某个 $w_0\in D$ 使得
对一切 $w\in D$, $(w_0-x_0,w-x_0)>0$.给出几何解释.\\
{[提示: 考虑集合 $C=\bigcup_{\mu>0}\mu(D-x_0)$.]}
\end{problem}
\begin{problem}[第 1, 5 章]
\index{MoreauProxyingshe@Moreau Prox 映射 (Moreau prox map)}
\label{prob:26}
\runin{Moreau 意义下的 Prox 映射}
设 $H$ 为 Hilbert 空间, $\varphi:H\to(-\infty,+\infty]$ 为满足 $\varphi\not\equiv+\infty$
的凸下半连续函数.

1. 证明对每个 $f\in H$, 存在某个 $u\in D(\varphi)$ 使得
\begin{equation}
\frac12\abs{f-u}^2+\varphi(u)=\inf_{v\in H}\left\{\frac12\abs{f-v}^2+\varphi(v)\right\}\equiv I.
\tag{P}
\label{eq:prob26.P}
\end{equation}
[提示: 先验证 $I>-\infty$.然后使用 Cauchy 序列论证或 $H$ 自反的事实.]

2. 验证 $u$ 满足 \eqref{eq:prob26.P} 当且仅当
\begin{equation}
u\in D(\varphi)\quad\text{且}\quad(u,v-u)+\varphi(v)-\varphi(u)\ge(f,v-u)\quad\forall v\in D(\varphi).
\tag{Q}
\label{eq:prob26.Q}
\end{equation}

3. 证明若 $u$ 与 $\bar u$ 分别是对应于 $f$ 与 $\bar f$ 的 \eqref{eq:prob26.P} 的解,
则 $\abs{u-\bar u}\le\abs{f-\bar f}$.由此推出 \eqref{eq:prob26.P} 的解的唯一性.

4. 考察 $\varphi=I_K$ 为闭凸集 $K$ 的示性函数的特殊情形.

5. 设 $\varphi^\star$ 为 $\varphi$ 的共轭函数, 考虑问题
\begin{equation}
\frac12\abs{f-u^\star}^2+\varphi^\star(u^\star)=\inf_{v\in H}\left\{\frac12\abs{f-v}^2+\varphi^\star(v)\right\}=I^\star.
\tag{P$^\star$}
\label{eq:prob26.Pstar}
\end{equation}
证明 \eqref{eq:prob26.P} 的解 $u$ 与 \eqref{eq:prob26.Pstar} 的解 $u^\star$ 满足
\begin{equation*}
u+u^\star=f\quad\text{且}\quad I+I^\star=\frac12\abs{f}^2.
\end{equation*}

6. 给定 $f\in H$ 与 $\lambda>0$, 设 $u_\lambda$ 表示问题
\begin{equation}
\frac12\abs{f-u_\lambda}^2+\lambda\varphi(u_\lambda)=\inf_{v\in H}\left\{\frac12\abs{f-v}^2+\lambda\varphi(v)\right\}
\tag{P$_\lambda$}
\label{eq:prob26.Plambda}
\end{equation}
的解.证明 $\lim_{\lambda\to0}u_\lambda=P_{\overline{D(\varphi)}}f$ ($=f$ 在
$\overline{D(\varphi)}$ 上的投影).\\
{[提示: 从弱收敛开始或使用习题 5.3.]}

7. 设 $K=\{v\in D(\varphi);\ \varphi(v)=\inf_H\varphi\}$ 并假设 $K\ne\emptyset$.
验证 $K$ 是闭凸集并证明 $\lim_{\lambda\to+\infty}u_\lambda=P_Kf$.当 $K=\emptyset$
时, 当 $\lambda\to+\infty$ 时 $(u_\lambda)$ 发生什么?

8. 证明 $\lim_{\lambda\to+\infty}\dfrac1\lambda u_\lambda=-P_{\overline{D(\varphi^\star)}}0$.\\
{[提示: 从 $f=0$ 的情形开始并应用第 5 与第 6 问.]}
\end{problem}

\begin{problem}[第~
\ref{chap:5}~章]
\label{prob:27}
\runin{交替投影}
设 $H$ 为 Hilbert 空间, $K\subset H$ 为非空闭凸集.验证对一切 $u,v\in H$,
\begin{equation*}
\abs{P_Ku-P_Kv}^2\le(P_Ku-P_Kv,u-v)\le\abs{u-v}^2.
\end{equation*}

设 $K_1\subset H$ 与 $K_2\subset H$ 为两个非空闭凸集.令 $P_1=P_{K_1}$ 与
$P_2=P_{K_2}$.给定 $u\in H$, 由归纳法定义序列 $(u_n)$ 如下:
\begin{equation*}
u_0=u,\quad u_1=P_1u_0,\quad u_2=P_2u_1,\quad\dots,\quad u_{2n-1}=P_1u_{2n-2},\quad u_{2n}=P_2u_{2n-1},\quad\dots
\end{equation*}

\runin{-- A --}

本部分的目的是证明序列 $(u_{2n}-u_{2n-1})$ 收敛于 $P_K0$, 其中 $K=K_2-K_1$
(注意 $K$ 是凸集, 为什么?).

1. 给定 $v\in H$, 考虑以 $v_0=v$ 开始由上述相同迭代定义的序列 $(v_n)$.验证
\begin{align*}
\abs{u_{2n}-v_{2n}}^2&\le(u_{2n}-v_{2n},u_{2n-1}-v_{2n-1})\le\abs{u_{2n-1}-v_{2n-1}}^2,\\
\abs{u_{2n+1}-v_{2n+1}}^2&\le(u_{2n+1}-v_{2n+1},u_{2n}-v_{2n})\le\abs{u_{2n}-v_{2n}}^2.
\end{align*}

2. 由此推出序列 $(\abs{u_n-v_n})$ 单调不增, 从而收敛于一个极限, 记作 $\ell$.
证明
\begin{equation*}
\lim_{n\to\infty}\abs{(u_{2n}-v_{2n})-(u_{2n-1}-v_{2n-1})}^2=0.
\end{equation*}

3. 验证序列 $(\abs{u_{2n}-u_{2n-1}})$ 单调不增.令
$d=\dist(K_1,K_2)=\inf\{\abs{a_1-a_2};\ a_1\in K_1,a_2\in K_2\}$.我们断言
$\lim_{n\to\infty}\abs{u_{2n}-u_{2n-1}}=d$.

4. 给定 $\varepsilon>0$, 选取 $v\in K_2$ 使得 $\dist(v,K_1)\le d+\varepsilon$.
证明对一切 $n$, $\abs{v_{2n}-v_{2n-1}}\le d+\varepsilon$.

5. 由此推出 $\lim_{n\to\infty}\abs{u_{2n}-u_{2n-1}}=d$.令 $z=P_K0$.

6. 验证 $\abs{z}=d$ 且对一切 $w\in K_2-K_1$, $\abs{z}^2\le(z,w)$.

7. 证明序列 $(u_{2n}-u_{2n-1})$ 收敛于 $z$.\\
{[提示: 使用上述结果估计 $\abs{z-(u_{2n}-u_{2n-1})}^2$.]}

8. 给出几何解释.

\runin{-- B --}

在整个本题的其余部分中, 我们假设 $z=P_K0\in K_2-K_1$ (例如若集合 $K_1$ 或 $K_2$
之一有界, 这个假设就成立, 为什么?).

我们断言存在 $a_1\in K_1$ 与 $a_2\in K_2$ 满足 $a_2-a_1=z$, 使得
$u_{2n}\rightharpoonup a_2$ 且 $u_{2n-1}\rightharpoonup a_1$ 弱收敛.注意 $a_1$
与 $a_2$ 可能依赖于初值 $u_0=u$ 的选取.画出示意图.

1. 考虑配以其自然标量积的 Hilbert 空间 $\mathcal{H}=H\times H$.令
$\mathcal{K}=\{[b_1,b_2]\in\mathcal{H};\ b_1\in K_1,b_2\in K_2\ \text{且}\ b_2-b_1=z\}$.
验证 $\mathcal{K}$ 是非空闭凸集.

2. 设 $b=[b_1,b_2]\in\mathcal{K}$.确定对应于 $v_0=b_1$ 的序列 $(v_n)$.
由此推出序列 $(\abs{u_{2n-1}-b_1})$ 与 $(\abs{u_{2n}-b_2})$ 单调不增.

3. 令 $x_n=[u_{2n-1},u_{2n}]$, 证明序列 $(x_n)$ 满足下述性质: 对弱收敛于某个
元素 $\bar x\in\mathcal{H}$ 的每个子序列 $(x_{n_k})$, 有 $\bar x\in\mathcal{K}$.

4. 应用 Opial 引理 (见习题 5.25 第 3 问) 并得出结论.
\end{problem}

\begin{problem}[第~
\ref{chap:5}~章]
\label{prob:28}
\runin{投影与正交投影}
设 $H$ 为 Hilbert 空间.满足 $P^2=P$ 的算子 $P\in\LL(H)$ 称为\emph{投影算子}.
验证投影算子满足下述性质:
(a) $I-P$ 是投影算子,
(b) $N(I-P)=R(P)$ 且 $N(P)=R(I-P)$,
(c) $N(P)\cap N(I-P)=\{0\}$,
(d) $H=N(P)+N(I-P)$.

\runin{-- A --}

若存在闭线性子空间 $M$ 使得 $P=P_M$ (其中 $P_M$ 在推论  \ref{cor:5.4} 中定义),
则称算子 $P\in\LL(H)$ 为\emph{正交投影算子}.验证每个正交投影算子是投影算子.

1. 给定投影算子 $P$, 证明下述性质等价:
(a) $P$ 是正交投影算子,
(b) $P^\star=P$,
(c) $\norm{P}\le1$,
(d) $N(P)\perp N(I-P)$,
其中记号 $X\perp Y$ 意味着对一切 $x\in X$ 及一切 $y\in Y$, $(x,y)=0$.

2. 设 $T\in\LL(H)$ 为满足 $T^\star=T$ 且 $T^2=I$ 的算子.证明 $P=\frac12(I-T)$
是正交投影算子.证明逆命题.额外假设对一切 $u\in H$ 有 $(Tu,u)\ge0$, 证明 $T=I$.

\runin{-- B --}

在整个本部分中, $M$ 与 $N$ 表示 $H$ 的两个闭线性子空间.令 $P=P_M$ 与 $Q=P_N$.

1. 证明下述性质等价:
(a) $PQ=QP$,
(b) $PQ$ 是投影算子,
(c) $QP$ 是投影算子.
在这时, 验证: (i) $PQ$ 是到 $M\cap N$ 上的正交投影算子, (ii) $(P+Q-PQ)$ 是到
$M+N$ 上的正交投影算子.

2. 证明下述性质等价:
(a) $M\perp N$,
(b) $PQ=0$,
(c) $QP=0$,
(d) 对一切 $u\in H$, $\abs{Pu}^2+\abs{Qu}^2\le\abs{u}^2$,
(e) 对一切 $u\in H$, $\abs{Pu}\le\abs{u-Qu}$,
(f) 对一切 $u\in H$, $\abs{Qu}\le\abs{u-Pu}$,
(g) $P+Q$ 是投影算子.
在这时, 验证 $(P+Q)$ 是到 $M+N$ 上的正交投影算子 (注意 $M+N$ 是闭的; 为什么?).

3. 证明下述性质等价:
(a) $M\subset N$,
(b) $PQ=P$,
(c) $QP=P$,
(d) 对一切 $u\in H$, $\abs{Pu}\le\abs{Qu}$,
(e) $Q-P$ 是投影算子.
在这时, 验证 $Q-P$ 是到 $M^\perp\cap N$ 上的正交投影算子.
\end{problem}

\begin{problem}[第~
\ref{chap:5}~章]
\label{prob:29}
\runin{非线性压缩映射的迭代. Opial 与 Baillon 遍历定理}
设 $H$ 为 Hilbert 空间, 设 $T:H\to H$ 为非线性压缩映射, 即对一切 $u,v\in H$,
$\abs{Tu-Tv}\le\abs{u-v}$.我们假设不动点集 $K=\{u\in H;\ Tu=u\}$ 非空.验证
$K$ 是闭凸集.给定 $f\in H$, 令
\begin{align*}
\sigma_n&=\frac1n(f+Tf+T^2f+\dots+T^{n-1}f),\\
\mu_n&=\left(\frac{I+T}{2}\right)^nf.
\end{align*}
本问题的目标是证明下述结论:
\begin{enumerate}[label=(\Alph*), itemsep=0.1em]
	\item 序列 $(\sigma_n)$ 与 $(\mu_n)$ 各自弱收敛于 $T$ 的某个不动点.
	\item 若此外 $T$ 是奇的, 即对一切 $v\in H$, $T(-v)=-Tv$, 则 $(\sigma_n)$ 与 $(\mu_n)$ 强收敛.
\end{enumerate}
在开始本题之前, 建议先做习题 5.22 与 5.25.在 $T$ 为线性的特殊情形下, 另见习题 5.21.

\runin{-- A --}

令 $u_n=T^nf$.

1. 验证对每个 $v\in K$, 序列 $(\abs{u_n-v})$ 单调不增.由此推出序列 $(\sigma_n)$
与 $(T\sigma_n)$ 有界.

2. 证明对一切 $n\ge1$, $\abs{\sigma_n-T\sigma_n}\le\frac{1}{\sqrt{n}}\abs{f-T\sigma_n}$.\\
{[提示: 注意 $\abs{T\sigma_n-Tu_i}^2\le\abs{\sigma_n-u_i}^2$, 并对 $0\le i\le n-1$ 将这些不等式相加.]}

3. 由此推出序列 $(\sigma_n)$ 满足习题 5.25 的性质 (P).得出结论: $\sigma_n\rightharpoonup\sigma$
弱收敛, 其中 $\sigma\in K$.

令 $S=\frac12(I+T)$.

4. 证明对一切 $u,v\in H$,
\begin{equation*}
\abs{(u-Su)-(v-Sv)}^2+\abs{Su-Sv}^2\le\abs{u-v}^2.
\end{equation*}

5. 由此推出对每个 $v\in K$, $\sum_{n=0}^\infty\abs{\mu_n-\mu_{n+1}}^2\le\abs{f-v}^2$,
从而对一切 $n$, $\abs{\mu_n-S\mu_n}\le\frac{1}{\sqrt{n+1}}\abs{f-v}$.

6. 得出结论: $\mu_n\rightharpoonup\mu$ 弱收敛, 其中 $\mu\in K$.

\runin{-- B --}

在整个本题的其余部分中, 我们假设 $T$ 是奇的, 即对一切 $v\in H$, $T(-v)=-Tv$.

1. 证明对每个整数 $p$, 对一切 $u,v\in H$,
\begin{equation*}
2\abs{(u,v)-(T^pu,T^pv)}\le\abs{u}^2+\abs{v}^2-\abs{T^pu}^2-\abs{T^pv}^2.
\end{equation*}
[提示: 从对一切 $u,v\in H$ 有 $\abs{T^pu-T^pv}^2\le\abs{u-v}^2$ 的不等式开始.]

2. 由此推出对每个固定的整数 $i\ge0$, $\ell(i)=\lim_{n\to\infty}(u_n,u_{n+i})$ 存在.
证明这个收敛关于 $i$ 一致成立, 即对一切 $i$ 与一切 $n$,
\begin{equation}
\abs{(u_n,u_{n+i})-\ell(i)}\le\varepsilon_n,
\label{eq:prob29.1}
\end{equation}
其中 $\lim_{n\to\infty}\varepsilon_n=0$.

3. 同理证明对每个固定的整数 $i\ge0$, $m(i)=\lim_{n\to\infty}(\mu_n,\mu_{n+i})$ 存在.
证明 $m(0)=m(1)=m(2)=\dots$.\\
{[提示: 使用第 A5 问的结果.]}

4. 由此推出 $\mu_n\to\mu$ 强收敛.

我们现在断言 $\sigma_n\to\sigma$ 强收敛.

5. 令 $X_p=\frac1p\sum_{i=0}^{p-1}\ell(i)$.证明对一切 $n$ 与一切 $p$,
\begin{equation*}
\abs{(u_n,\sigma_{n+p})-X_p}\le\varepsilon_n+\frac{2n}{p}\abs{f}^2.
\end{equation*}
[提示: 使用 \eqref{eq:prob29.1}.]

6. 由此推出: (i) $X=\lim_{p\to\infty}X_p$ 存在, (ii) 对一切 $n$, $\abs{(u_n,\sigma)-X}\le\varepsilon_n$,
(iii) $\abs{\sigma}^2=X$.

7. 证明 $\abs{\sigma_n}^2\le\frac{2}{n^2}\sum_{i=0}^{n-1}(n-i)\ell(i)+\frac2n\sum_{i=0}^{n-1}\varepsilon_i$.

8. 由此推出 $\limsup_{n\to\infty}\abs{\sigma_n}^2\le X$ 并得出结论.
\end{problem}

\begin{problem}[第 3, 5 章]
\label{prob:30}
\runin{Stampacchia 定理的变体. von Neumann 极小化极大定理}
设 $H$ 为 Hilbert 空间.

\runin{-- A --}

设 $a(u,v):H\times H\to\R$ 为满足对一切 $v\in H$ 有 $a(v,v)\ge0$ 的连续双线性形式.
设 $K\subset H$ 为非空闭凸集.设 $f\in H$.假设存在某个 $v_0\in K$ 使得集合
$\{u\in K;\ a(u,v_0-u)\ge(f,v_0-u)\}$ 是有界的.

1. 证明存在某个 $u\in K$ 使得对一切 $v\in K$, $a(u,v-u)\ge(f,v-u)$.\\
{[提示: 令 $f_\varepsilon=f+\varepsilon v_0$, 并考虑双线性形式
$a_\varepsilon(u,v)=a(u,v)+\varepsilon(u,v)$ ($\varepsilon>0$).然后利用习题 5.14
令 $\varepsilon\to0$ 取极限.]}

2. 恢复 Stampacchia 定理.

3. 在 $K$ 有界且对一切 $u,v\in H$ 有 $a(u,v)=0$ 的情形下给出几何解释.

\runin{-- B --}

设 $b(u,v):H\times H\to\R$ 为连续且强制的双线性形式.设 $\varphi:H\to(-\infty,+\infty]$
为满足 $\varphi\not\equiv+\infty$ 的凸下半连续函数.

1. 证明存在唯一的 $u\in D(\varphi)$ 使得对一切 $v\in D(\varphi)$,
\begin{equation*}
b(u,v-u)+\varphi(v)-\varphi(u)\ge0.
\end{equation*}
[提示: 在空间 $H\times\R$ 中应用第 A1 问的结果, 取 $K=\operatorname{epi}\varphi$,
$f=[0,-1]$, $a(U,V)=b(u,v)$, 其中 $U=[u,\lambda],V=[v,\mu]$.注意 $a$ 不是强制的.]

2. 恢复 Stampacchia 定理.

\runin{-- C --}

设 $H_1$ 与 $H_2$ 为两个 Hilbert 空间, 设 $A\subset H_1,B\subset H_2$ 为两个非空, 有界, 闭的凸集.

1. 设 $F(\lambda,\mu):H_1\times H_2\to\R$ 为连续双线性形式.证明存在 $\bar\lambda\in A$
与 $\bar\mu\in B$ 使得对一切 $\lambda\in A$ 与一切 $\mu\in B$,
\begin{equation}
F(\bar\lambda,\mu)\le F(\bar\lambda,\bar\mu)\le F(\lambda,\bar\mu).
\label{eq:prob30.1}
\end{equation}
[提示: 在 $H=H_1\times H_2$, $\mathcal{K}=A\times B$ 上应用第 A1 问, 取
$a(u,v)=F(\bar\lambda,\mu)-F(\lambda,\bar\mu)$, 其中 $u=[\bar\lambda,\bar\mu],v=[\lambda,\mu]$.]

2. 由此推出
\begin{equation}
\min_{\lambda\in A}\max_{\mu\in B}F(\lambda,\mu)=\max_{\mu\in B}\min_{\lambda\in A}F(\lambda,\mu).
\label{eq:prob30.2}
\end{equation}
注意所有的极小与极大都可以达到 (为什么?).\\
{[提示: 验证没有任何附加假设下就有 $\max\min\le\min\max$; 使用 \eqref{eq:prob30.1}
证明反向不等式.]}

3. 证明 \eqref{eq:prob30.2} 蕴含满足 \eqref{eq:prob30.1} 的某个 $\bar\lambda\in A$
与 $\bar\mu\in B$ 的存在性.

\runin{-- D --}

设 $E$ 与 $F$ 为两个自反 Banach 空间; 设 $A\subset E$ 与 $B\subset F$ 为两个非空, 有界, 闭的凸集.设 $K:E\times F\to\R$ 为满足下述假设的函数:
\begin{enumerate}[label=(\alph*), itemsep=0.1em]
	\item 对每个固定的 $v\in B$, 函数 $u\mapsto K(u,v)$ 是凸的且下半连续.
	\item 对每个固定的 $u\in A$, 函数 $v\mapsto K(u,v)$ 是凹的且上半连续, 即函数 $v\mapsto-K(u,v)$ 是凸的且下半连续.
\end{enumerate}
我们的目标是证明
\begin{equation*}
\min_{u\in A}\max_{v\in B}K(u,v)=\max_{v\in B}\min_{u\in A}K(u,v).
\end{equation*}
我们将用反证法并假设存在常数 $\gamma$ 使得
\begin{equation*}
\max_{v\in B}\min_{u\in A}K(u,v)<\gamma<\min_{u\in A}\max_{v\in B}K(u,v).
\end{equation*}

1. 对每个 $u\in A$, 令 $B_u=\{v\in B;\ K(u,v)\ge\gamma\}$; 对每个 $v\in B$, 令
$A_v=\{u\in A;\ K(u,v)\le\gamma\}$.验证 $\bigcap_{u\in A}B_u=\emptyset$ 且
$\bigcap_{v\in B}A_v=\emptyset$.

2. 选取 $u_1,\dots,u_n\in A$ 与 $v_1,\dots,v_m\in B$ 使得 $\bigcap_{i=1}^nB_{u_i}=\emptyset$
且 $\bigcap_{j=1}^mA_{v_j}=\emptyset$ (说明理由).在 $H_1=\R^n,H_2=\R^m$,
$A'=\left\{\lambda=(\lambda_1,\dots,\lambda_n);\ \lambda_i\ge0\ \forall i,\ \sum_{i=1}^n\lambda_i=1\right\}$,
$B'=\left\{\mu=(\mu_1,\dots,\mu_m);\ \mu_j\ge0\ \forall j,\ \sum_{j=1}^m\mu_j=1\right\}$,
以及 $F(\lambda,\mu)=\sum_{i,j}\lambda_i\mu_jK(u_i,v_j)$ 上应用第 C1 问的结果.
令 $\bar u=\sum_i\bar\lambda_iu_i$ 与 $\bar v=\sum_j\bar\mu_jv_j$.证明对一切
$k=1,\dots,n$ 与一切 $\ell=1,\dots,m$,
\begin{equation*}
K(\bar u,v_\ell)\le K(u_k,\bar v).
\end{equation*}

3. 验证另一方面有 $\min_k K(u_k,\bar v)<\gamma$ 以及 $\max_\ell K(\bar u,v_\ell)>\gamma$.\\
{[提示: 用反证法.]}

4. 得出结论.
\end{problem}
\begin{problem}[第 3, 5 章]
\label{prob:31}
\runin{单调算子. Minty--Browder 定理}
设 $E$ 为自反 Banach 空间.(非线性) 映射 $A:D(A)\subset E\to E^\star$ 若满足
\begin{equation*}
\dual{Au-Av}{u-v}\ge0\quad\forall u,v\in D(A),
\end{equation*}
则称它是\emph{单调的} (这里 $D(A)$ 表示 $E$ 的任意子集).

\runin{-- A --}

设 $A:D(A)\subset E\to E^\star$ 为单调映射, 设 $K\subset E$ 为非空, 有界, 闭的凸集.
我们的目标是证明存在某个 $u\in K$ 使得
\begin{equation*}
\dual{Av}{u-v}\ge0\quad\forall v\in D(A)\cap K.
\end{equation*}
为此, 对每个 $v\in D(A)\cap K$, 令 $K_v=\{u\in K;\ \dual{Av}{v-u}\ge0\}$.须证明
$\bigcap_{v\in D(A)\cap K}K_v\ne\emptyset$; 我们用反证法并假设
$\bigcap_{v\in D(A)\cap K}K_v=\emptyset$.

1. 验证 $K_v$ 是闭凸集.

2. 由此推出存在 $v_1,v_2,\dots,v_n\in D(A)\cap K$ 使得 $\bigcap_{i=1}^nK_{v_i}=\emptyset$.
令
\begin{equation*}
B=\left\{\lambda=(\lambda_1,\dots,\lambda_n);\ \lambda_i\ge0\ \forall i,\ \sum_{i=1}^n\lambda_i=1\right\},
\end{equation*}
并考虑由 $F(\lambda,\mu)=\sum_{i,j=1}^n\lambda_i\mu_j\dual{Av_j}{v_i-v_j}$ 定义的双线性形式
$F:\R^n\times\R^n\to\R$.

3. 验证对一切 $\lambda\in\R^n$, $F(\lambda,\lambda)\le0$.

4. 证明存在某个 $\bar\lambda\in B$ 使得对一切 $\mu\in B$, $F(\bar\lambda,\mu)\le0$.\\
{[提示: 应用问题 30 第 C1 问.]}

5. 令 $\bar u=\sum_{i=1}^n\bar\lambda_iv_i$, 证明对一切 $j=1,\dots,n$,
$\dual{Av_j}{\bar u-v_j}\le0$.

6. 得出结论.

\runin{-- B --}

在整个本题的其余部分中, 我们假设 $D(A)=E$, $A:E\to E^\star$ 单调, 且 $A$ 连续.

1. 设 $K\subset E$ 为非空, 有界, 闭的凸集.证明存在某个 $u\in K$ 使得对一切
$w\in K$, $\dual{Au}{w-u}\ge0$.\\
{[提示: 考虑 $v_t=(1-t)u+tw$, 其中 $t\in(0,1),w\in K$.]}

2. 设 $K$ 为包含 $0$ 的闭凸集 ($K$ 不必有界).假设集合 $\{u\in K;\ \dual{Au}{u}\le0\}$
是有界的.证明存在某个 $u\in K$ 使得对一切 $v\in K$, $\dual{Au}{v-u}\ge0$.\\
{[提示: 对充分大的 $R$, 把 B1 应用于集合 $K_R=\{v\in K;\ \norm{v}\le R\}$.]}

3. 假设 $\lim_{\norm{v}\to\infty}\dfrac{\dual{Av}{v}}{\norm{v}}=+\infty$.证明 $A$ 是满射.

4. 这里假设 $E$ 是与 $E^\star$ 等同的 Hilbert 空间.证明 $I+A$ 是由 $E$ 到自身上的双射.
\end{problem}

\begin{problem}[第~
\ref{chap:5}~章]
\label{prob:32}
\runin{压缩映射的延拓. Kirszbraun--Valentine 定理的 Schoenberg 方法}
设 $H$ 为 Hilbert 空间, $I$ 为有限指标集.

\runin{-- A --}

设 $(y_i)_{i\in I}$ 为 $H$ 中的元素, $(c_i)_{i\in I}$ 为 $\R$ 中的元素.令
\begin{equation*}
\phi(u)=\max_{i\in I}\{\abs{u-y_i}^2-c_i\},\quad u\in H,
\end{equation*}
以及 $J(u)=\{i\in I;\ \abs{u-y_i}^2-c_i=\phi(u)\}$.

1. 验证 $\inf_{u\in H}\phi(u)$ 由某个唯一的元素 $u_0\in H$ 达到.

2. 证明对一切 $v\in H$, $\max_{i\in J(u_0)}(v,u_0-y_i)\ge0$.

3. 由此推出
\begin{equation}
u_0\in\operatorname{conv}\left(\bigcup_{i\in J(u_0)}\{y_i\}\right).
\label{eq:prob32.1}
\end{equation}

4. 反之, 若 $u_0\in H$ 满足 \eqref{eq:prob32.1}, 证明 $\phi(u_0)=\inf_{u\in H}\phi(u)$.

5. 把这一结果推广到 $\phi(u)=\max_{i\in I}\{f_i(u)\}$ 的情形, 其中每个 $f_i:H\to\R$
是凸的 $C^1$ 函数.

\runin{-- B --}

设 $(x_i)_{i\in I}$ 与 $(y_i)_{i\in I}$ 为 $H$ 中的元素, 满足对一切 $i,j\in I$,
$\abs{y_i-y_j}\le\abs{x_i-x_j}$.我们断言: 给定任何 $p\in H$, 存在某个
$q\in\operatorname{conv}\left(\bigcup_{i\in I}\{y_i\}\right)$ 使得对一切 $i\in I$,
$\abs{q-y_i}\le\abs{p-x_i}$.

1. 令 $P=\{\lambda=(\lambda_i)_{i\in I};\ \lambda_i\ge0\ \forall i\ \text{且}\ \sum_{i\in I}\lambda_i=1\}$.
证明对每个 $p\in H$ 与每个 $\lambda\in P$,
\begin{equation*}
\sum_{j\in I}\lambda_j\abs{\sum_{i\in I}\lambda_iy_i-y_j}^2\le\sum_{j\in I}\lambda_j\abs{p-x_j}^2.
\end{equation*}
[提示: 验证 $\sum_{j\in I}\lambda_j\abs{\sum_{i\in I}\lambda_iy_i-y_j}^2=\frac12\sum_{i,j\in I}\lambda_i\lambda_j\abs{y_i-y_j}^2$.]

2. 考虑函数 $\phi(u)=\max_{i\in I}\{\abs{u-y_i}^2-\abs{p-x_i}^2\}$.设 $u_0\in H$
使得 $\phi(u_0)=\inf_{u\in H}\phi(u)$.证明 $\phi(u_0)\le0$.\\
{[提示: 应用第 A3 与 B1 问.]}

3. 得出结论.

\runin{-- C --}

1. 把第 B 部分的结果推广到 $I$ 为无限指标集的情形.

2. 设 $D\subset H$ 为 $H$ 的任意子集, 设 $S:D\to H$ 为压缩映射, 即对一切 $u,v\in D$,
$\abs{Su-Sv}\le\abs{u-v}$.证明存在定义在整个 $H$ 上的压缩映射 $T$, 延拓 $S$ 且满足
$T(H)\subset\overline{\operatorname{conv}S(D)}$.\\
{[提示: 使用 Zorn 引理与第 C1 问.]}
\end{problem}

\begin{problem}[第 4, 6 章]
\label{prob:33}
\runin{$L^p$ 中的乘法算子}
设 $\Omega$ 为测度空间 (有限或无限测度).令 $E=L^p(\Omega)$, 其中 $1<p<\infty$.
设 $a:\Omega\to\R$ 为可测函数.考虑由
\begin{equation*}
D(A)=\{u\in L^p(\Omega);\ au\in L^p(\Omega)\}\quad\text{以及}\quad Au=au
\end{equation*}
定义的无界线性算子 $A:D(A)\subset E\to E$.

1. 证明 $D(A)$ 在 $E$ 中稠密.\\
{[提示: 给定 $u\in E$, 考虑序列 $u_n(x)=(1+n^{-1}\abs{a(x)})^{-1}u(x)$.]}

2. 证明 $A$ 是闭的.

3. 证明 $D(A)=E$ 当且仅当 $a\in L^\infty(\Omega)$.\\
{[提示: 应用闭图像定理.]}

4. 确定 $N(A)$ 与 $N(A)^\perp$.

5. 确定 $D(A^\star)$, $A^\star$, $N(A^\star)$ 与 $N(A^\star)^\perp$.

6. 证明 $A$ 是满射当且仅当存在 $\alpha>0$ 使得在 $\Omega$ 上几乎处处 $\abs{a(x)}\ge\alpha$.\\
{[提示: 使用第 3 问.]}

在下文中假设 $a\in L^\infty(\Omega)$.

7. 确定 $A$ 的特征值与谱.验证 $\sigma(A)\subset[\operatorname{ess\,inf}_\Omega a,\operatorname{ess\,sup}_\Omega a]$,
且 $\operatorname{ess\,inf}_\Omega a\in\sigma(A)$, $\operatorname{ess\,sup}_\Omega a\in\sigma(A)$.

8. 若 $\Omega$ 是 $\R^N$ 中的开集 (配以 Lebesgue 测度) 且 $a\in C(\Omega)\cap L^\infty(\Omega)$,
证明 $\sigma(A)=\overline{a(\Omega)}$.

9. 证明 $\sigma(A)=\{0\}$ 当且仅当在 $\Omega$ 上几乎处处 $a=0$.

10. 假设 $\Omega$ 没有原子.证明 $A$ 紧当且仅当在 $\Omega$ 上几乎处处 $a=0$.
\end{problem}

\begin{problem}[第 4, 6 章]
\label{prob:34}
\runin{Hardy 算子 $Tu(x)=\frac1x\int_0^xu(t)dt$ 的谱分析}

\runin{-- A --}

设 $E=C([0,1])$ 配以范数 $\norm{u}=\sup_{t\in[0,1]}\abs{u(t)}$.给定 $u\in E$, 在 $[0,1]$
上定义函数 $Tu$:
\begin{equation*}
Tu(x)=\begin{cases}
\dfrac1x\int_0^xu(t)dt & \text{若}\ x\in(0,1],\\
u(0) & \text{若}\ x=0.
\end{cases}
\end{equation*}
验证 $Tu\in E$ 且对一切 $u\in E$, $\norm{Tu}\le\norm{u}$, 从而 $T\in\LL(E)$.

1. 证明 $\operatorname{EV}(T)=(0,1]$ 并确定相应的特征函数.

2. 验证 $\norm{T}_{\LL(E)}=1$.$T$ 是由 $E$ 到自身中的紧算子吗?

3. 证明 $\sigma(T)=[0,1]$.当 $\lambda\in\rho(T)$ 时给出 $(T-\lambda I)^{-1}$ 的显式公式.
证明对每个不属于 $\{0,1\}$ 的 $\lambda\in\R$, $(T-\lambda I)$ 是由 $E$ 到 $E$ 上的满射.
验证 $T$ 与 $(T-I)$ 不是满射.

4. 在本问中把 $T$ 看作由 $E=C([0,1])$ 到 $F=L^q(0,1)$ ($1\le q<\infty$) 中的有界算子.
证明 $T\in\KK(E,F)$.\\
{[提示: 考虑由 $(T_\varepsilon u)(x)=\dfrac{1}{x+\varepsilon}\int_0^xu(t)dt$ ($\varepsilon>0$)
定义的算子, 并估计当 $\varepsilon\to0$ 时 $\norm{T_\varepsilon-T}_{\LL(E,F)}$.]}

\runin{-- B --}

在本部分中设 $E=C^1([0,1])$ 配以范数 $\norm{u}=\sup_{t\in[0,1]}\abs{u(t)}+\sup_{t\in[0,1]}\abs{u'(t)}$.
给定 $u\in C^1([0,1])$, 如第 A 部分中那样定义 $Tu$.

1. 验证若 $u\in C^1([0,1])$, 则 $Tu\in C^1([0,1])$ 且对一切 $u\in E$, $\norm{Tu}\le\norm{u}$.

2. 证明 $\operatorname{EV}(T)=(0,1/2]\cup\{1\}$.

3. 证明 $\sigma(T)=[0,1/2]\cup\{1\}$.

\runin{-- C --}

在本部分中设 $E=L^p(0,1)$ ($1<p<\infty$).给定 $u\in L^p(0,1)$, 对 $x\in(0,1]$ 由
$Tu(x)=\frac1x\int_0^xu(t)dt$ 定义 $Tu$.验证 $Tu\in C((0,1])$ 且对每个 $q<p$,
$Tu\in L^q(0,1)$.我们的目标是证明 $Tu\in L^p(0,1)$ 且
\begin{equation}
\norm{Tu}_{L^p(0,1)}\le\frac{p}{p-1}\norm{u}_{L^p(0,1)}\quad\forall u\in E.
\label{eq:prob34.1}
\end{equation}

1. 证明当 $u\in C_c((0,1))$ 时 \eqref{eq:prob34.1} 成立.\\
{[提示: 令 $\phi(x)=\int_0^xu(t)dt$; 验证 $\abs{\phi}^p\in C^1([0,1])$ 并计算其导数.
使用公式 $\int_0^1\abs{Tu(x)}^pdx=\frac{1}{p-1}\int_0^1\abs{\phi(x)}^pd\left(-\frac{1}{x^{p-1}}\right)$
并分部积分估计 $\norm{Tu}_{L^p}$.]}

2. 证明对每个 $u\in E$, \eqref{eq:prob34.1} 成立.

在下文中把 $T$ 看作由 $E$ 到自身中的有界算子.

3. 证明 $\operatorname{EV}(T)=\left(0,\frac{p}{p-1}\right)$.

4. 由此推出 $\norm{T}_{\LL(E)}=\frac{p}{p-1}$.$T$ 是由 $E$ 到自身中的紧算子吗?

5. 证明 $\sigma(T)=\left[0,\frac{p}{p-1}\right]$.

6. 确定 $T^\star$.

7. 在本问中把 $T$ 看作由 $E=L^p(0,1)$ 到 $F=L^q(0,1)$ ($1\le q<p<\infty$) 中的有界算子.
证明 $T\in\KK(E,F)$.
\end{problem}

\begin{problem}[第~
\ref{chap:6}~章]
\label{prob:35}
\runin{Cotlar 引理}
设 $H$ 为与其对偶空间等同的 Hilbert 空间.

\runin{-- A --}

假设 $T\in\LL(H)$, 从而 $T^\star\in\LL(H)$.

1. 证明 $\norm{T^\star T}=\norm{T}^2$.

2. 在本问中假设 $T$ 自伴.证明对每个整数 $N$, $\norm{T^N}=\norm{T}^N$.

3. 由此推出 (对一般的 $T\in\LL(H)$), 对每个整数 $N$, $\norm{(T^\star T)^N}=\norm{T}^{2N}$.

\runin{-- B --}

设 $(T_j)_{1\le j\le m}$ 为 $\LL(H)$ 中的算子的有限族.假设对一切 $j,k\in\{1,2,\dots,m\}$,
\begin{align}
\norm{T_j^\star T_k}^{1/2}&\le\omega(j-k), \label{eq:prob35.1}\\
\norm{T_kT_j^\star}^{1/2}&\le\omega(k-j), \label{eq:prob35.2}
\end{align}
其中 $\omega:\Z\to[0,\infty)$.令 $\sigma=\sum_{i=-(m-1)}^{m-1}\omega(i)$.
本问题的目标是证明
\begin{equation}
\norm{\sum_{j=1}^m T_j}\le\sigma.
\label{eq:prob35.3}
\end{equation}
令 $U=\sum_{j=1}^mT_j$, 并固定整数 $N$.

1. 证明对整数 $j_1,k_1,\dots,j_N,k_N\in\{1,2,\dots,m\}$ 的任何选取,
\begin{equation*}
\norm{T_{j_1}^\star T_{k_1}T_{j_2}^\star T_{k_2}\dots T_{j_N}^\star T_{k_N}}
\le\sigma\omega(j_1-k_1)\omega(k_1-j_2)\dots\omega(k_{N-1}-j_N)\omega(j_N-k_N).
\end{equation*}

2. 由此推出 $\sum_{j_1}\sum_{k_1}\dots\sum_{j_N}\sum_{k_N}\norm{T_{j_1}^\star T_{k_1}\dots T_{j_N}^\star T_{k_N}}\le m\sigma^{2N}$.

3. 证明 $\norm{(U^\star U)^N}\le m\sigma^{2N}$, 并由此推出 \eqref{eq:prob35.3} 成立.
\end{problem}
\begin{problem}[第~
\ref{chap:6}~章]
\label{prob:36}
\runin{关于 Riesz--Fredholm 理论的补充}
设 $E$ 为 Banach 空间, $T\in\KK(E)$.对每个整数 $k\ge1$, 令
\begin{equation*}
N_k=N((I-T)^k)\quad\text{以及}\quad R_k=R((I-T)^k).
\end{equation*}

1. 验证对一切 $k\ge1$, $R_{k+1}\subset R_k$, $R_k$ 闭, $T(R_k)\subset R_k$, 以及
$(I-T)R_k\subset R_{k+1}$.

2. 证明存在整数 $p\ge1$ 使得对一切 $k<p$, $R_{k+1}\ne R_k$ (若 $p=1$ 则无条件),
且对一切 $k\ge p$, $R_{k+1}=R_k$.

3. 验证对一切 $k\ge1$, $N_k\subset N_{k+1}$, $\dim N_k<\infty$, $T(N_k)\subset N_k$,
以及 $(I-T)N_{k+1}\subset N_k$.

4. 证明对一切 $k\ge1$, $\operatorname{codim}R_k=\dim N_k$, 并由此推出对一切 $k<p$,
$N_{k+1}\ne N_k$ (若 $p=1$ 则无条件), 且对一切 $k\ge p$, $N_{k+1}=N_k$.

5. 证明 $R_p\cap N_p=\{0\}$ 且 $R_p+N_p=E$.

6. 证明 $(I-T)$ 限制在 $R_p$ 上是由 $R_p$ 到自身上的双射.

7. 这里额外假设 $E$ 是 Hilbert 空间且 $T$ 自伴.证明 $p=1$.
\end{problem}

\begin{problem}[第~
\ref{chap:6}~章]
\index{Rayleighshang@Rayleigh 商 (Rayleigh quotient)}
\index{CourantFischerjixiaohuajidayuanli@Courant--Fischer 极小化极大原理 (Courant--Fischer min-max principle)}
\label{prob:37}
\runin{Courant--Fischer 极小化极大原理. Rayleigh--Ritz 方法}
设 $H$ 为无限维可分 Hilbert 空间.设 $T$ 为由 $H$ 到自身的自伴紧算子, 满足对一切
$x\in H$, $(Tx,x)\ge0$.用 $(\mu_k)$ ($k\ge1$) 记其按不增顺序排列的, 计入重数的
特征值:
\begin{equation*}
\mu_1\ge\mu_2\ge\dots\ge0.
\end{equation*}
设 $(e_j)$ 为由特征向量组成的相应标准正交基.设 $E_k$ 为由 $\{e_1,e_2,\dots,e_k\}$
张成的空间.对 $x\ne0$, 我们定义 \emph{Rayleigh 商}
\begin{equation*}
R(x)=\frac{(Tx,x)}{\abs{x}^2}.
\end{equation*}

1. 证明对一切 $k\ge1$, $\min_{\substack{x\in E_k\\x\ne0}}R(x)=\mu_k$.

2. 证明对一切 $k\ge2$, $\max_{\substack{x\in E_{k-1}^\perp\\x\ne0}}R(x)=\mu_k$,
且 $\max_{\substack{x\in H\\x\ne0}}R(x)=\mu_1$.

3. 设 $\Sigma$ 为 $H$ 的任何 $k$ 维子空间 ($k\ge1$).证明
$\min_{\substack{x\in\Sigma\\x\ne0}}R(x)\le\mu_k$.\\
{[提示: 若 $k\ge2$, 证明 $\Sigma\cap E_{k-1}^\perp\ne\{0\}$ 并应用第 2 问.]}

4. 由此推出对一切 $k\ge1$,
\begin{equation*}
\max_{\substack{\Sigma\subset H\\\dim\Sigma=k}}\min_{\substack{x\in\Sigma\\x\ne0}}R(x)=\mu_k.
\end{equation*}

5. 设 $\Sigma$ 为 $H$ 的任何 $(k-1)$ 维子空间 ($k\ge2$).证明
$\max_{\substack{x\in\Sigma^\perp\\x\ne0}}R(x)\ge\mu_k$.\\
{[提示: 证明 $\Sigma^\perp\cap E_k\ne\{0\}$.]}

6. 由此推出对一切 $k\ge1$,
\begin{equation*}
\min_{\substack{\Sigma\subset H\\\dim\Sigma=k-1}}\max_{\substack{x\in\Sigma^\perp\\x\ne0}}R(x)=\mu_k.
\end{equation*}

7. 这里假设 $N(T)=\{0\}$, 从而对一切 $x\ne0$, $R(x)\ne0$, 或等价地对一切 $k$,
$\mu_k>0$.证明对一切 $k\ge1$,
\begin{equation*}
\min_{\substack{\Sigma\subset H\\\dim\Sigma=k}}\max_{\substack{x\in\Sigma\\x\ne0}}\frac{1}{R(x)}=\frac{1}{\mu_k},
\end{equation*}
以及
\begin{equation*}
\max_{\substack{M\subset H\\\operatorname{codim}M=k-1}}\min_{\substack{x\in M\\x\ne0}}\frac{1}{R(x)}=\frac{1}{\mu_k},
\end{equation*}
其中 $\Sigma$ 与 $M$ 是 $H$ 的闭子空间.特别地, 对 $k=1$,
$\min_{\substack{x\in H\\x\ne0}}\frac{1}{R(x)}=\frac{1}{\mu_1}$; 此外对一切 $k\ge2$,
$\min_{\substack{x\in E_{k-1}^\perp\\x\ne0}}\frac{1}{R(x)}=\frac{1}{\mu_k}$.

8. 设 $V$ 为 $H$ 的闭子空间 (有限维或无限维).设 $P_V$ 为由 $H$ 到 $V$ 上的正交
投影算子, 考虑由 $S=P_V\circ T|_V$ 定义的算子 $S:V\to V$.验证 $S$ 是由 $V$ 到
自身的自伴紧算子, 且对一切 $x\in V$, $(Sx,x)\ge0$.

9. 用 $(\nu_k)$ ($k\ge1$) 记 $S$ 的按不增顺序排列的, 计入重数的特征值.证明对
满足 $1\le k\le\dim V$ 的一切 $k$,
\begin{equation*}
\max_{\substack{\Sigma\subset V\\\dim\Sigma=k}}\min_{\substack{x\in\Sigma\\x\ne0}}R(x)=\nu_k.
\end{equation*}
由此推出对满足 $1\le k\le\dim V$ 的一切 $k$, $\nu_k\le\mu_k$.

10. 现在考虑 $H$ 的闭子空间的递增序列 $V^{(n)}$, 满足 $\overline{\bigcup_nV^{(n)}}=H$.
令 $S^{(n)}=P_{V^{(n)}}\circ T|_{V^{(n)}}$, 并设 $(\nu_k^{(n)})$ 记如第 9 问中那样排列的
$S^{(n)}$ 的特征值.证明对每个固定的 $k$, 序列 $n\mapsto\nu_k^{(n)}$ 单调不减, 且当
$n\to\infty$ 时收敛于 $\mu_k$.
\end{problem}

\begin{problem}[第 2, 6, 11 章]
\index{FredholmNoethersuanzi@Fredholm--Noether 算子 (Fredholm--Noether operator)}
\label{prob:38}
\runin{Fredholm--Noether 算子}
设 $E$ 与 $F$ 为 Banach 空间, $T\in\LL(E,F)$.

\runin{-- A --}

第 A 部分的目标是证明下述条件等价:
\begin{equation}
\text{(a)}\ R(T)\ \text{闭且在}\ F\ \text{中有有限余维,}\quad
\text{(b)}\ N(T)\ \text{在}\ E\ \text{中有补空间}.
\tag{1}
\label{eq:prob38.1}
\end{equation}
\begin{equation}
\text{存在}\ S\in\LL(F,E)\ \text{与}\ K\in\KK(F,F)\ \text{使得}\ T\circ S=I_F+K.
\tag{2}
\label{eq:prob38.2}
\end{equation}
\begin{equation}
\text{存在}\ U\in\LL(F,E)\ \text{与}\ F\ \text{中的有限秩投影}\ P\ \text{使得}\ T\circ U=I_F-P.
\tag{3}
\label{eq:prob38.3}
\end{equation}
此外, 可以选取 $U$ 与 $P$ 使得 $\dim R(P)=\operatorname{codim}R(T)$.

1. 证明 \eqref{eq:prob38.1} $\Rightarrow$ \eqref{eq:prob38.3}.\\
{[提示: 设 $X$ 为 $N(T)$ 在 $E$ 中的补空间.则 $T|_X$ 是由 $X$ 到 $R(T)$ 上的双射.
用 $U_0$ 记其逆算子.设 $Q$ 为由 $F$ 到 $R(T)$ 上的投影算子, 令 $U=U_0\circ Q$.]}

2. 证明 \eqref{eq:prob38.2} $\Rightarrow$ \eqref{eq:prob38.3}.\\
{[提示: 使用习题 6.25.]}

3. 证明 \eqref{eq:prob38.3} $\Rightarrow$ \eqref{eq:prob38.1}.\\
{[提示: 为建立 \eqref{eq:prob38.1} 的 (a) 部分, 注意 $R(T)\supset R(I_F-P)$ 并应用
命题  \ref{prop:11.5}.同理证明 $R(U^\star)$ 闭, 从而 $R(U)$ 也闭.最后证明在 $E$ 中
存在有限维空间 $\Sigma_1$ 与 $\Sigma_2$ 使得 $N(T)+R(U)+\Sigma_1=E$ 且
$N(T)\cap R(U)\subset\Sigma_2$.然后应用命题  \ref{prop:11.7}.]}

4. 得出结论.

\runin{-- B --}

证明下述条件等价:
\begin{equation}
\text{(a)}\ R(T)\ \text{闭且有补空间,}\quad\text{(b)}\ \dim N(T)<\infty.
\tag{4}
\label{eq:prob38.4}
\end{equation}
\begin{equation}
\text{存在}\ \widetilde S\in\LL(F,E)\ \text{与}\ \widetilde K\in\KK(E,E)\ \text{使得}\ \widetilde S\circ T=I_E+\widetilde K.
\tag{5}
\label{eq:prob38.5}
\end{equation}
\begin{equation}
\text{存在}\ \widetilde U\in\LL(F,E)\ \text{与}\ E\ \text{中的有限秩投影}\ \widetilde P\ \text{使得}\ \widetilde U\circ T=I_E-\widetilde P.
\tag{6}
\label{eq:prob38.6}
\end{equation}

\runin{-- C --}

若算子 $T\in\LL(E,F)$ 满足
\begin{equation}
\text{(a)}\ R(T)\ \text{闭且有有限余维,}\quad\text{(b)}\ \dim N(T)<\infty,
\tag{FN}
\label{eq:prob38.FN}
\end{equation}
则称它为 \emph{Fredholm} 算子 (或 \emph{Noether} 算子) ($R(T)$ 闭的性质可由其他
假设推出; 见习题 2.27).满足 \eqref{eq:prob38.FN} 的算子类记作 $\Phi(E,F)$.$T$ 的
\emph{指标}按定义为
\begin{equation*}
\operatorname{ind}T=\dim N(T)-\operatorname{codim}R(T).
\end{equation*}

1. 设 $T\in\Phi(E,F)$.证明存在 $U\in\LL(F,E)$ 以及 $F$ (相应地 $E$) 中的有限秩投影
$P$ (相应地 $\widetilde P$), 使得
\begin{equation}
T\circ U=I_F-P,\qquad U\circ T=I_E-\widetilde P,
\label{eq:prob38.7}
\end{equation}
满足 $\dim R(P)=\operatorname{codim}R(T)$, $\dim R(\widetilde P)=\dim N(T)$.\\
{[提示: 使用第 A1 问中构造的算子 $U$.]}

满足 $T\circ V=I_F+K$ 且 $V\circ T=I_E+\widetilde K$ ($K\in\KK(F)$, $\widetilde K\in\KK(E)$)
的算子 $V\in\LL(F,E)$ 称为 $T$ 的\emph{拟逆} (或模紧算子的逆).

2. 证明任何拟逆 $V$ 属于 $\Phi(F,E)$.

3. 证明算子 $T\in\LL(E,F)$ 属于 $\Phi(E,F)$ 当且仅当 $R(T)$ 闭, $\dim N(T)<\infty$,
且 $\dim N(T^\star)<\infty$.此外,
\begin{equation*}
\operatorname{ind}T=\dim N(T)-\dim N(T^\star).
\end{equation*}
[提示: 应用命题  \ref{prop:11.14} 与命题 2.18.]
% 原书此处印作 "Propositions 11.14 and 2.18" (印刷页 494); 第~
% \ref{chap:2}~章并无编号 2.18
% 的命题 (2.18 是推论~
% \ref{cor:2.18}). 依技能规范保留原样并记录, 不擅自更正, 故此处不用
% \ref 以免产生未定义引用.

4. 设 $T\in\Phi(E,F)$.证明 $T^\star\in\Phi(F^\star,E^\star)$ 且
$\operatorname{ind}T^\star=-\operatorname{ind}T$.\\
{[提示: 应用命题  \ref{prop:11.13} 与定理  \ref{thm:2.19}.]}

5. 反之, 设 $T\in\LL(E,F)$ 满足 $T^\star\in\Phi(F^\star,E^\star)$.证明 $T\in\Phi(E,F)$.

6. 假设 $J\in\LL(E,F)$ 是双射且 $K\in\KK(E,F)$.证明 $T=J+K$ 属于 $\Phi(E,F)$ 且
$\operatorname{ind}T=0$.反之, 若 $T\in\Phi(E,F)$ 且 $\operatorname{ind}T=0$, 证明
$T$ 可以写成 $T=J+K$, 其中 $J$ 与 $K$ 如上 (甚至可以选取 $K$ 为有限秩算子).\\
{[提示: 应用定理  \ref{thm:6.6}, 证明 $I_E+J^{-1}\circ K$ 属于 $\Phi(E,E)$ 且指标为零.
关于逆命题, 考虑由 $N(T)$ 到 $R(T)$ 的补空间 $Y$ 上的同构.]}

7. 设 $T\in\Phi(E,F)$ 且 $K\in\KK(E,F)$.证明 $T+K\in\Phi(E,F)$.

8. 在前一问的假设下, 证明 $\operatorname{ind}(T+K)=\operatorname{ind}T$.\\
{[提示: 令 $\widetilde E=E\times Y$, $\widetilde F=F\times N(T)$, 以及由
$\widetilde T(x,y)=(Tx+Kx,0)$ 定义的 $\widetilde T:\widetilde E\to\widetilde F$.
证明 $\widetilde T=\widetilde J+\widetilde K$, 其中 $\widetilde J$ 是由 $\widetilde E$
到 $\widetilde F$ 上的双射, 且 $\widetilde K\in\KK(\widetilde E,\widetilde F)$.
然后应用第 6 问.]}

9. 设 $T\in\Phi(E,F)$.证明存在 $\varepsilon>0$ (依赖于 $T$), 使得对满足
$\norm{M}<\varepsilon$ 的每个 $M\in\LL(E,F)$, 有 $T+M\in\Phi(E,F)$.证明
$\operatorname{ind}(T+M)=\operatorname{ind}T$.\\
{[提示: 设 $V$ 为 $T$ 的拟逆.则若 $\norm{M}<\norm{V}^{-1}$, $W=I_E+(V\circ M)$
是双射.验证 $T+M=(T\circ W)+\text{紧算子}$; 然后应用前一问.]}

10. 设 $(H_t)_{t\in[0,1]}$ 为 $\LL(E,F)$ 中的算子族.假设 $t\mapsto H_t$ 从 $[0,1]$
到 $\LL(E,F)$ 连续, 且对一切 $t\in[0,1]$, $H_t\in\Phi(E,F)$.证明 $\operatorname{ind}H_t$
在 $[0,1]$ 上为常数.

11. 设 $E_1,E_2$ 与 $E_3$ 为 Banach 空间, 设 $T_1\in\Phi(E_1,E_2)$, $T_2\in\Phi(E_2,E_3)$.
证明 $T_2\circ T_1\in\Phi(E_1,E_3)$.

12. 沿用上面的相同记号, 证明
\begin{equation*}
\operatorname{ind}(T_2\circ T_1)=\operatorname{ind}T_1+\operatorname{ind}T_2.
\end{equation*}
[提示: 考虑用矩阵记号对 $t\in[0,1]$ 由
\begin{equation*}
H_t=\begin{pmatrix}I&0\\0&T_2\end{pmatrix}\begin{pmatrix}(1-t)I&tI\\-tI&(1-t)I\end{pmatrix}\begin{pmatrix}T_1&0\\0&I\end{pmatrix}
\end{equation*}
定义的算子族 $H_t:E_1\times E_2\to E_2\times E_3$, 其中 $I$ 是 $E_2$ 中的恒等算子.
验证 $t\mapsto H_t$ 从 $[0,1]$ 到 $\LL(E_1\times E_2,E_2\times E_3)$ 连续.利用前一问
证明对每个 $t$, $H_t\in\Phi(E_1\times E_2,E_2\times E_3)$.计算 $\operatorname{ind}H_0$
与 $\operatorname{ind}H_1$.]

13. 设 $T\in\Phi(E,F)$.计算 $T$ 的任何拟逆 $V$ 的指标.

\runin{-- D --}

在本部分中我们研究两个简单的例子.

1. 假设 $\dim E<\infty$ 且 $\dim F<\infty$.证明从 $E$ 到 $F$ 的任何线性算子 $T$
属于 $\Phi(E,F)$ 并计算其指标.

2. 设 $E=F=\ell^2$.考虑习题 6.18 中定义的移位算子 $S_r$ 与 $S_\ell$.证明对满足
$\lambda\ne\pm1$ 的每个 $\lambda\in\R$, 有 $S_r-\lambda I\in\Phi(\ell^2,\ell^2)$
且 $S_\ell-\lambda I\in\Phi(\ell^2,\ell^2)$.计算它们的指标.证明 $S_r\pm I$,
$S_\ell\pm I$ 不属于 $\Phi(\ell^2,\ell^2)$.\\
{[提示: 使用习题 6.18 的结果.]}
\end{problem}

\begin{problem}[第 5, 6 章]
\label{prob:39}
\runin{自伴非负算子的平方根}
设 $H$ 为 Hilbert 空间.设 $S\in\LL(H)$; 若对一切 $x\in H$, $(Sx,x)\ge0$, 则称 $S$ 是
\emph{非负的}, 记作 $S\ge0$.当 $S_1,S_2\in\LL(H)$ 时, 若 $S_1-S_2\ge0$, 则记
$S_1\ge S_2$ (或 $S_2\le S_1$).

\runin{-- A --}

1. 设 $S\in\LL(H)$ 满足 $S^\star=S$ 且 $0\le S\le I$.证明 $\norm{S^2}=\norm{S}^2\le1$,
且 $0\le S^2\le S\le I$.\\
{[提示: 使用习题 6.24.]}

2. 设 $S\in\LL(H)$ 满足 $S^\star=S$ 且 $S\ge0$.设 $P(t)=\sum a_kt^k$ 为对一切 $k$
满足 $a_k\ge0$ 的多项式.证明 $[P(S)]^\star=P(S)$ 且 $P(S)\ge0$.

3. 设 $(S_n)$ 为 $\LL(H)$ 中的序列使得对一切 $n$, $S_n^\star=S_n$, 且对一切 $n$,
$S_{n+1}\le S_n$.假设存在常数 $M$ 使得对一切 $n$, $\norm{S_n}\le M$.证明对每个
$x\in H$, 当 $n\to\infty$ 时 $S_nx$ 收敛于一个极限, 记作 $Sx$, 且 $S\in\LL(H)$
满足 $S^\star=S$.\\
{[提示: 设 $n\ge m$.使用习题 6.24 证明 $\abs{S_nx-S_mx}^2\le2M(S_mx-S_nx,x)$.]}

\runin{-- B --}

假设 $T\in\LL(H)$ 满足 $T^\star=T$, $T\ge0$, 且 $\norm{T}\le1$.考虑从 $S_0=I$ 开始
由下式定义的序列 $(S_n)$:
\begin{equation*}
S_{n+1}=S_n+\frac12(T-S_n^2),\quad n\ge0.
\end{equation*}

1. 证明对一切 $n\ge0$, $S_n^\star=S_n$.

2. 证明 $I-S_{n+1}=\frac12(I-S_n)^2+\frac12(I-T)$, 并由此推出对一切 $n$, $I-S_n\ge0$.

3. 证明对一切 $n$, $S_n\ge0$.\\
{[提示: 利用第 A1 与 B2 问由数学归纳法证明 $I-S_n\le I$.]}

4. 由此推出对一切 $n$, $\norm{S_n}\le1$.

5. 证明对一切 $n$, $S_n-S_{n+1}=\frac12[(I-S_n)+(I-S_{n-1})]\circ(S_{n-1}-S_n)$,
并由此推出对一切 $n$, $S_{n-1}-S_n\ge0$.\\
{[提示: 由数学归纳法证明 $(I-S_n)=P_n(I-T)$ 与 $(S_{n-1}-S_n)=Q_n(I-T)$, 其中 $P_n$
与 $Q_n$ 是具有非负系数的多项式.]}

6. 证明 $\lim_{n\to\infty}S_nx=Sx$ 存在.证明 $S\in\LL(H)$ 满足 $S^\star=S$, $S\ge0$,
$\norm{S}\le1$, 且 $S^2=T$.

\runin{-- C --}

1. 设 $U\in\LL(H)$ 满足 $U^\star=U$ 且 $U\ge0$.证明存在 $V\in\LL(H)$ 使得 $V^\star=V$,
$V\ge0$, 且 $V^2=U$.\\
{[提示: 对 $T=U/\norm{U}$ 应用第 B 部分的构造.]}

接着我们证明 $V$ 的唯一性.更精确地说, 若 $W$ 是满足 $W^\star=W$, $W\ge0$ 且 $W^2=U$
的任何算子 $W\in\LL(H)$, 则 $W=V$.算子 $V$ 称为 $U$ 的\emph{平方根}, 记作 $U^{1/2}$.

2. 证明上述构造的算子 $V$ 与每个与 $U$ 对易的算子 $X$ 对易 (即 $X\circ U=U\circ X$
蕴含 $X\circ V=V\circ X$).

3. 证明 $W$ 与 $U$ 对易, 并由此推出 $V$ 与 $W$ 对易.

4. 验证 $(V-W)\circ(V+W)=0$, 并由此推出在 $R(V+W)$ 上 $V=W$.证明
$N(V)=N(W)=N(U)=N(V+W)$.得出在 $H$ 上 $V=W$ 的结论.\\
{[提示: 注意在 $\overline{R(V+W)}=N(U)^\perp$ 上 $V=W$, 且在 $N(U)$ 上 $V=W=0$.]}

5. 证明 $\norm{U^{1/2}}=\norm{U}^{1/2}$.

6. 设 $U_1,U_2\in\LL(H)$ 满足 $U_1^\star=U_1,U_2^\star=U_2,U_1\ge0,U_2\ge0$, 且
$U_1\circ U_2=U_2\circ U_1$.证明 $U_1\circ U_2\ge0$.\\
{[提示: 引入 $U_1^{1/2}$ 与 $U_2^{1/2}$.]}

\runin{-- D --}

设 $U\in\KK(H)$ 满足 $U^\star=U$ 且 $U\ge0$.证明其平方根 $V$ 属于 $\KK(H)$.假设
$H$ 可分, 在由 $U$ 的特征向量组成的 Hilbert 基上计算 $V$.求 $V$ 的特征值.
\end{problem}

\begin{problem}[第 4, 5, 6 章]
\index{HilbertSchmidtsuanzi@Hilbert--Schmidt 算子 (Hilbert--Schmidt operator)}
\label{prob:40}
\runin{Hilbert--Schmidt 算子}

\runin{-- A --}

设 $E$ 与 $F$ 为可分 Hilbert 空间, 都与其对偶空间等同.$E$ 与 $F$ 上的范数用相同符号
$\abs{\cdot}$ 表示.设 $T\in\LL(E,F)$, 从而 $T^\star\in\LL(F,E)$.

1. 设 $(e_k)$ (相应地 $(f_k)$) 为 $E$ (相应地 $F$) 的任意标准正交基.证明
$\sum_{k=1}^\infty\abs{T(e_k)}^2<\infty$ 当且仅当 $\sum_{k=1}^\infty\abs{T^\star(f_k)}^2<\infty$,
且 $\sum_{k=1}^\infty\abs{T(e_k)}^2=\sum_{k=1}^\infty\abs{T^\star(f_k)}^2$.

2. 设 $(e_k)$ 与 $(\tilde e_k)$ 为 $E$ 的两个标准正交基.证明
$\sum_{k=1}^\infty\abs{T(e_k)}^2<\infty$ 当且仅当 $\sum_{k=1}^\infty\abs{T(\tilde e_k)}^2<\infty$,
且 $\sum_{k=1}^\infty\abs{T(e_k)}^2=\sum_{k=1}^\infty\abs{T(\tilde e_k)}^2$.

若存在 $E$ 的某个标准正交基 $(e_k)$ 使得 $\sum_{k=1}^\infty\abs{T(e_k)}^2<\infty$,
则称 $T\in\LL(E,F)$ 为 \emph{Hilbert--Schmidt} 算子, 记作 $T\in\mathcal{HS}(E,F)$.

3. 证明 $\mathcal{HS}(E,F)$ 是 $\LL(E,F)$ 的线性子空间, 且
\begin{equation*}
\norm{T}_{\mathcal{HS}}=\left(\sum_{k=1}^\infty\abs{T(e_k)}^2\right)^{1/2}
\end{equation*}
定义了 $\mathcal{HS}(E,F)$ 上的范数.设 $\norm{\cdot}$ 表示 $\LL(E,F)$ 上的通常范数.
证明对一切 $T\in\mathcal{HS}(E,F)$, $\norm{T}\le\norm{T}_{\mathcal{HS}}$.

4. 证明配以范数 $\norm{\cdot}_{\mathcal{HS}}$ 的 $\mathcal{HS}(E,F)$ 是 Banach 空间.
证明事实上它是 Hilbert 空间.

5. 证明 $\mathcal{HS}(E,F)\subset\KK(E,F)$.\\
{[提示: 给定 $x\in E$, 写出 $x=\sum_{k=1}^\infty x_ke_k$, 并令
$T_n(x)=\sum_{k=1}^nx_kT(e_k)$.证明当 $n\to\infty$ 时 $\norm{T_n-T}\to0$.]}

6. 证明由 $E$ 到 $F$ 中的任何有限秩算子属于 $\mathcal{HS}(E,F)$.

7. 设 $T\in\LL(E,F)$.证明 $T\in\mathcal{HS}(E,F)\iff T^\star\in\mathcal{HS}(F,E)$,
且 $\norm{T^\star}_{\mathcal{HS}(F,E)}=\norm{T}_{\mathcal{HS}(E,F)}$.

8. 假设 $T\in\KK(E,E)$ 满足 $T^\star=T$, 设 $(\lambda_k)$ 记 $T$ 的特征值序列.
证明 $T\in\mathcal{HS}(E,E)$ 当且仅当 $\sum_{k=1}^\infty\lambda_k^2<\infty$, 且
$\norm{T}_{\mathcal{HS}}^2=\sum_{k=1}^\infty\lambda_k^2$.构造一个在 $E=\ell^2$ 上的
算子 $T\in\KK(E,E)$ 使得 $T\notin\mathcal{HS}(E,E)$.

9. 设 $G$ 为另一个可分 Hilbert 空间.设 $T_1\in\LL(E,F)$ 且 $T_2\in\LL(F,G)$.
证明若 $T_1$ 或 $T_2$ 之一属于 $\mathcal{HS}$, 则 $T_2\circ T_1\in\mathcal{HS}(E,G)$.

10. 设 $T\in\mathcal{HS}(E,E)$ 并假设 $N(I+T)=\{0\}$.证明 $(I+T)$ 是双射, 且
$(I+T)^{-1}=I+S$, 其中 $S\in\mathcal{HS}(E,E)$.

11. 设 $(e_k)$ (相应地 $(f_\ell)$) 为 $E$ (相应地 $F$) 的标准正交基.考虑由
$T_{k,\ell}(x)=(x,e_k)f_\ell$ 定义的算子 $T_{k,\ell}:E\to F$.证明 $(T_{k,\ell})$
是 $\mathcal{HS}(E,F)$ 的标准正交基.

\runin{-- B --}

假设 $\Omega$ 是 $\R^N$ 中的开子集.在下文中取 $E=F=L^2(\Omega)$.设
$K\in L^2(\Omega\times\Omega)$, 考虑算子
\begin{equation}
(Tu)(x)=\int_\Omega K(x,y)u(y)dy.
\label{eq:prob40.1}
\end{equation}

1. 证明 $T\in\LL(E,E)$ 且 $\norm{T}_{\LL(E,E)}\le\norm{K}_{L^2(\Omega\times\Omega)}$.

2. 证明 $T\in\mathcal{HS}(E,E)$ 且 $\norm{T}_{\mathcal{HS}(E,E)}\le\norm{K}_{L^2(\Omega\times\Omega)}$.\\
{[提示: 设 $(e_j)$ 为 $L^2(\Omega)$ 的标准正交基.验证族 $e_{j,k}=e_j\otimes e_k$
(其中 $(e_j\otimes e_k)(x,y)=e_j(x)e_k(y)$) 是 $L^2(\Omega\times\Omega)$ 的标准正交基.
然后写出 $\norm{T(e_k)}_{L^2(\Omega)}^2=\sum_{j=1}^\infty\abs{(T(e_k),e_j)}^2=\sum_{j=1}^\infty\abs{(K,e_j\otimes e_k)}$.]}

3. 反之, 设 $T\in\mathcal{HS}(E,E)$.证明存在唯一的函数 $K\in L^2(\Omega\times\Omega)$
使得 \eqref{eq:prob40.1} 成立.$K$ 称为 $T$ 的\emph{积分核}.\\
{[提示: 令 $t_{j,k}=(Te_k,e_j)$ 并验证 $\sum_{j,k=1}^\infty\abs{t_{j,k}}^2<\infty$.
定义 $K=\sum_{j,k=1}^\infty t_{j,k}e_j\otimes e_k$ 并证明 \eqref{eq:prob40.1} 成立.]}

4. 假设 $\Omega=(0,1)$, $E=L^2(\Omega)$, 考虑算子 $(Tu)(x)=\int_0^xu(t)dt$.证明
$T\in\mathcal{HS}(E,E)$ 并计算 $\norm{T}_{\mathcal{HS}}$.
\end{problem}
\begin{problem}[第 1, 6 章]
\label{prob:41}
\runin{Krein--Rutman 定理}
设 $E$ 为 Banach 空间, $P\subset E$ 为包含 $0$ 的闭凸集.假设 $P$ 是以 $0$ 为顶点的
凸锥, 即对一切 $\lambda>0,\mu>0,x\in P,y\in P$, $\lambda x+\mu y\in P$.

假设
\begin{equation}
\operatorname{Int}P\ne\emptyset\quad\text{且}\quad P\ne E.
\label{eq:prob41.1}
\end{equation}
设 $T\in\KK(E)$ 满足
\begin{equation}
T(P\setminus\{0\})\subset\operatorname{Int}P.
\label{eq:prob41.2}
\end{equation}

\runin{-- A --}

1. 证明 $(\operatorname{Int}P)\cap(-P)=\emptyset$.\\
{[提示: 使用习题 1.7.]}

在下文中我们固定某个 $u\in\operatorname{Int}P$.

2. 证明存在 $\alpha>0$ 使得对一切 $x\in P$, $\norm{x+u}\ge\alpha$.\\
{[提示: 用反证法并得出 $-u\in P$.]}

3. 验证存在 $r>0$ 使得 $Tu-ru\in P$.

4. 假设某个 $x\in P$ 满足对某个 $\lambda\in\R$, $T(x+u)=\lambda x$.证明 $\lambda\ge r$.\\
{[提示: 在 $E$ 上引入由 $y\ge z\iff y-z\in P$ 定义的序关系是方便的.由数学归纳法证明
$(\lambda/r)^nx\ge u$, $n=1,2,\dots$.]}

5. 考虑非线性映射 $F(x)=T\left(\dfrac{x+u}{\norm{x+u}}\right)$, $x\in P$.证明
$F:P\to P$ 连续且对 $E$ 的某个紧子集 $K$ 有 $F(P)\subset K$.由此推出存在某个
$x_1\in P$ 使得 $T(x_1+u)=\lambda_1x_1$, 其中 $\lambda_1=\norm{x_1+u}\ge r$.\\
{[提示: 应用 Schauder 不动点定理; 见习题 6.26.]}

6. 由此推出对每个 $\varepsilon>0$, 存在 $x_\varepsilon\in P$ 使得
$T(x_\varepsilon+\varepsilon u)=\lambda_\varepsilon x_\varepsilon$, 其中
$\lambda_\varepsilon=\norm{x_\varepsilon+\varepsilon u}\ge r$.

7. 证明存在 $x_0\in\operatorname{Int}P$ 与 $\mu_0>0$ 使得 $Tx_0=\mu_0x_0$.\\
{[提示: 证明 $(x_\varepsilon)$ 有界.由此推出存在满足所要求性质的序列
$\varepsilon_n\to0$ 使得 $x_{\varepsilon_n}\to x_0$ 且 $\lambda_{\varepsilon_n}\to\mu_0$.]}

\runin{-- B --}

1. 给定两点 $a\in\operatorname{Int}P$ 与 $b\in E\setminus P$, 证明存在唯一的
$\sigma\in(0,1)$ 使得对一切 $t\in[0,\sigma)$, $(1-t)a+tb\in\operatorname{Int}P$;
$(1-\sigma)a+\sigma b\in P$; 且对一切 $t\in(\sigma,1]$, $(1-t)a+tb\notin P$.
然后令 $\tau(a,b)=\sigma/(1-\sigma)$, 满足 $0<\tau(a,b)<\infty$.

2. 设 $x\in P\setminus\{0\}$ 满足对某个 $\mu\in\R$, $Tx=\mu x$.证明 $\mu=\mu_0$
且对某个 $m>0$ 有 $x=mx_0$, 其中 $\mu_0$ 与 $x_0$ 是在第 A7 问中构造的.\\
{[提示: 反设对一切 $m>0$ 都有 $x\ne mx_0$.证明 $\mu>0,x\in\operatorname{Int}P$,
且 $-x\notin P$.令 $y=x_0-\tau_0x$, 其中 $\tau_0=\tau(x_0,-x)$.计算 $Ty$ 并得出
$\mu<\mu_0$.然后互换 $x_0$ 与 $x$ 的角色.]}

3. 设 $x\in E\setminus\{0\}$ 满足对某个 $\mu\in\R$, $Tx=\mu x$.证明或者 $\mu=\mu_0$
且 $x=mx_0$ ($m\in\R,m\ne0$), 或者 $\abs{\mu}<\mu_0$.\\
{[提示: 由第 2 问可以假设 $x\notin P$ 且 $-x\notin P$.若 $\mu>0$ 考虑 $\tau(x_0,x)$;
若 $\mu<0$ 则同时考虑 $\tau(x_0,x)$ 与 $\tau(x_0,-x)$.]}

4. 由此推出 $N(T-\mu_0I)=\R x_0$.换言之, 特征值 $\mu_0$ 的几何重数为 1.

5. 证明对所有 $k\ge2$, $N((T-\mu_0I)^k)=\R x_0$.换言之, 特征值 $\mu_0$ 的代数重数也是 1.\\
{[提示: 鉴于问题 36, 只须证明 $N((T-\mu_0I)^2)=\R x_0$.]}
\end{problem}

\begin{problem}[第~
\ref{chap:6}~章]
\label{prob:42}
\runin{关于不变子空间的 Lomonosov 定理}
设 $E$ 为无限维 Banach 空间, 设 $T\in\KK(E)$, $T\ne0$.第 A 部分的目标是证明
存在 $T$ 的非平凡闭不变子空间 $Z$, 即 $T(Z)\subset Z$, 其中 $Z\ne\{0\}$ 且 $Z\ne E$.

\runin{-- A --}

令
\begin{align*}
\mathcal{A}&=\operatorname{span}\{I,T,T^2,\dots\}\\
&=\left\{\sum_{i\in I}\lambda_iT^i;\ \lambda_i\in\R,\ I\ \text{是}\ \{0,1,2,\dots\}\ \text{的有限子集}\right\}.
\end{align*}
对每个 $y\in E$, 令 $\mathcal{A}y=\{Sy;\ S\in\mathcal{A}\}$.显然 $y\in\mathcal{A}y$,
因此对每个 $y\ne0$, $\mathcal{A}y\ne\{0\}$.此外 $\mathcal{A}y$ 是 $E$ 的子空间
且 $T(\mathcal{A}y)\subset\mathcal{A}y$, 从而 $T(\overline{\mathcal{A}y})\subset\overline{\mathcal{A}y}$.
若对某个 $y\ne0$ 有 $\overline{\mathcal{A}y}\ne E$, 则 $\overline{\mathcal{A}y}$
是 $T$ 的非平凡闭不变子空间.因此我们可以假设对一切 $y\in E$ ($y\ne0$),
\begin{equation}
\overline{\mathcal{A}y}=E.
\label{eq:prob42.1}
\end{equation}

由于 $T\ne0$, 我们可以固定某个 $x_0\in E$ 使得 $Tx_0\ne0$, 并固定 $r$ 使得
$0<r\le\dfrac{\norm{Tx_0}}{2\norm{T}}\le\dfrac{\norm{x_0}}{2}$.令
$C=\{x\in E;\ \norm{x-x_0}\le r\}$.

1. 验证 $0\notin C$, 且对一切 $x\in C$, $\norm{Tx-Tx_0}\le\frac12\norm{Tx_0}$,
从而对一切 $x\in C$, $\norm{Tx}\ge\frac12\norm{Tx_0}$.由此推出 $0\notin\overline{T(C)}$.

2. 证明对每个 $y\in E$ ($y\ne0$), 存在某个 $S\in\mathcal{A}$ (记作 $S_y$), 使得
$\norm{S_yy-x_0}\le r/2$.\\
{[提示: 使用假设 \eqref{eq:prob42.1}.]}

3. 由此推出对每个 $y\in E$ ($y\ne0$), 存在某个 $\varepsilon_y>0$ (依赖于 $y$),
使得对一切 $z\in B(y,\varepsilon_y)$, $\norm{S_yz-x_0}\le r$.

4. 考虑 $T(C)$ 被球 $B(y_j,\frac12\varepsilon_{y_j})$ ($j\in J,J$ 有限) 的有限覆盖.
对 $j\in J$ 与 $x\in E$, 令 $q_j(x)=\max\{0,\varepsilon_{y_j}-\norm{Tx-y_j}\}$ 以及
$q(x)=\sum_{j\in J}q_j(x)$.验证函数 $q_j$ 与 $q$ 在 $E$ 上连续.证明对一切 $x\in C$,
$q(x)\ge\min_{j\in J}\frac12\varepsilon_{y_j}>0$.对 $x\in C$, 令
\begin{equation*}
F(x)=\frac{1}{q(x)}\sum_{j\in J}q_j(x)S_{y_j}(Tx).
\end{equation*}

5. 证明 $F$ 从 $C$ 到 $E$ 连续, 且对一切 $x\in C$, $\norm{F(x)-x_0}\le r$.\\
{[提示: 使用第 3 问.]}

6. 证明 $F(C)\subset K$, 其中 $K$ 是 $C$ 的紧子集.由此推出存在 $\xi\in C$ 使得
$F(\xi)=\xi$.\\
{[提示: 应用 Schauder 不动点定理; 见习题 6.26.]}

7. 令 $U=\dfrac{1}{q(\xi)}\sum_{j\in J}q_j(\xi)(S_{y_j}\circ T)$, 其中 $\xi$ 如第 6 问中那样.
证明 $U\in\KK(E)$.由此推出 $Z=N(I-U)$ 是有限维的; 验证 $\xi\in Z$.

8. 证明 $T(Z)\subset Z$ 并得出结论.\\
{[提示: 证明 $U\in\mathcal{A}$, 由此推出 $T\circ U=U\circ T$.]}

9. 构造没有非平凡不变子空间的线性算子 $T:\R^2\to\R^2$.

\runin{-- B --}

现在建立上述结果的一个更强版本.假设 $T\in\KK(E)$ 且 $T\ne0$.设 $R\in\LL(E)$ 满足
$R\circ T=T\circ R$.证明 $R$ 允许非平凡闭不变子空间.\\
{[提示: 令 $\mathcal{B}=\operatorname{span}\{I,R,R^2,\dots\}$ 且
$\mathcal{B}y=\{Sy;\ S\in\mathcal{B}\}$.验证第 A 部分中的所有步骤在把 $\mathcal{A}$
换成 $\mathcal{B}$, 把 $\mathcal{A}y$ 换成 $\mathcal{B}y$ 之后仍然成立.]}
\end{problem}

\begin{problem}[第 2, 4, 5, 6 章]
\label{prob:43}
\runin{正规算子}
设 $H$ 为与其对偶空间等同的 Hilbert 空间.若算子 $T\in\LL(H)$ 满足
$T\circ T^\star=T^\star\circ T$, 则称它为\emph{正规算子}.

1. 证明 $T$ 正规当且仅当它满足对一切 $u\in H$, $\abs{Tu}=\abs{T^\star u}$.\\
{[提示: 计算 $\abs{T(u+v)}^2$.]}

在整个本题的其余部分中, 我们假设 $T$ 是正规的.

2. 假设 $u\in N(T-\lambda I)$ 且 $v\in N(T-\mu I)$, 其中 $\lambda\ne\mu$.证明
$(u,v)=0$.\\
{[提示: 利用第 1 问证明 $N(T^\star-\mu I)=N(T-\mu I)$, 并计算 $(Tu,v)$.]}

3. 证明 $\overline{R(T)}=\overline{R(T^\star)}=N(T)^\perp=N(T^\star)^\perp$.

4. 设 $f\in R(T)$.验证存在 $u\in R(T^\star)$ 满足 $f=Tu$.\\
{[提示: 注意 $H=\overline{R(T)}\oplus N(T)$.]}

5. 考虑 $R(T^\star)$ 中的序列 $(u_n)$ 使得当 $n\to\infty$ 时 $u_n\to u$.写出
$u_n=T^\star y_n$ ($y_n\in H$).证明当 $n\to\infty$ 时 $Ty_n$ 收敛于极限 $z\in H$,
满足 $T^\star z=f$.\\
{[提示: 使用第 1 问与 Cauchy 序列论证.]}

6. 由此推出 $R(T)=R(T^\star)$.\\
{[提示: 利用 $N(T)=N(T^\star)$ 的事实.]}

7. 证明 $\norm{T^2}=\norm{T}^2$.\\
{[提示: 写出 $\abs{Tu}^2\le\abs{T^\star Tu}\abs{u}=\abs{T^2u}\abs{u}$.]}

8. 由此推出对每个整数 $p\ge1$, $\norm{T^p}=\norm{T}^p$.\\
{[提示: 先考虑 $p=2^k$ 的情形.对一般的整数 $p$, 选取满足 $2^k\ge p$ 的任何 $k$,
并写出 $\norm{T}^{2^k}=\norm{T^{2^k}}=\norm{T^{2^k-p}T^p}$.]}

9. 证明 $N(T^2)=N(T)$, 并由此推出对每个整数 $p\ge1$, $N(T^p)=N(T)$.\\
{[提示: 注意若 $T^2u=0$, 则 $Tu\in N(T)\cap R(T)$.]}
\end{problem}

\begin{problem}[第 5, 6 章]
\label{prob:44}
\runin{等距与酉算子. 斜自伴算子. 极分解与 Cayley 变换}
设 $H$ 为与其对偶空间等同的 Hilbert 空间, 设 $T\in\LL(H)$.称
\begin{enumerate}[label=(\roman*), itemsep=0.1em]
	\item $T$ 为等距, 若对一切 $u\in H$, $\abs{Tu}=\abs{u}$,
	\item $T$ 为酉算子, 若 $T$ 是等距且又是满射,
	\item $T$ 为斜自伴的 (或反对称的), 若 $T^\star=-T$.
\end{enumerate}

\runin{-- A --}

1. 假设 $T$ 是等距.验证 $\norm{T}=1$.

2. 证明 $T\in\LL(H)$ 是等距当且仅当 $T^\star\circ T=I$.

3. 假设 $T\in\LL(H)$ 是等距.证明下述条件等价:
(a) $T$ 是酉算子,
(b) $T^\star$ 是单射,
(c) $T\circ T^\star=I$,
(d) $T^\star$ 是等距,
(e) $T^\star$ 是酉算子.

4. 举出一个不是酉算子的等距的例子.\\
{[提示: 使用习题 6.18.]}

5. 假设 $T$ 是等距.证明 $R(T)$ 闭且 $T\circ T^\star=P_{R(T)}$ ($=R(T)$ 上的正交投影算子).

6. 假设 $T$ 是等距.证明或者 $T$ 是酉算子此时 $\sigma(T)\subset\{-1,+1\}$, 或者 $T$ 不是酉算子此时 $\sigma(T)=[-1,+1]$.

7. 假设 $T\in\KK(H)$ 是等距.证明 $\dim H<\infty$.

8. 证明 $T\in\LL(H)$ 斜自伴当且仅当对一切 $u\in H$, $(Tu,u)=0$.

9. 假设 $T\in\LL(H)$ 斜自伴.证明 $\sigma(T)\subset\{0\}$.\\
{[提示: 使用 Lax--Milgram 定理.]}

10. 假设 $T\in\LL(H)$ 斜自伴.令 $U=(T+I)\circ(T-I)^{-1}$.验证 $U$ 良定,
$U=(T-I)^{-1}\circ(T+I)$, 且 $U\circ T=T\circ U$.证明 $U$ 是酉算子 ($U$ 称为 $T$
的 \emph{Cayley 变换}).

11. 反之, 设 $T\in\LL(H)$ 满足 $1\notin\sigma(T)$.假设 $U=(T+I)\circ(T-I)^{-1}$
是等距.证明 $T$ 斜自伴.

\runin{-- B --}

若算子 $T\in\LL(H)$ 满足
\begin{equation}
\text{存在由}\ N(T)\ \text{到}\ N(T^\star)\ \text{中的等距}\ J,
\tag{1}
\label{eq:prob44.1}
\end{equation}
则称它满足性质 \eqref{eq:prob44.1}.

第 B 部分的目标是证明满足性质 \eqref{eq:prob44.1} 的每个算子 $T\in\LL(H)$ 可以
分解为 $T=U\circ P$, 其中 $U\in\LL(H)$ 是等距且 $P\in\LL(H)$ 是自伴非负算子
(回忆非负意味着对一切 $u\in H$, $(Pu,u)\ge0$).这样的分解称为 $T$ 的\emph{极分解}.
此外, $P$ 在 $H$ 上唯一确定, $U$ 在 $N(T)^\perp$ 上唯一确定 (但在 $H$ 上不唯一).

1. 验证在下述情形下假设 \eqref{eq:prob44.1} 成立:
(i) $T$ 是单射,
(ii) $\dim H<\infty$,
(iii) $T$ 正规 (见问题 43),
(iv) $T=I-K$, 其中 $K\in\KK(H)$.

2. 举出不满足 \eqref{eq:prob44.1} 的例子.\\
{[提示: 使用习题 6.18.]}

3. 假设有极分解 $T=U\circ P$.证明 $P^2=T^\star\circ T$.

4. 由此推出 $P$ 在 $H$ 上唯一确定.\\
{[提示: 使用问题 39.]}

5. 设 $T=U\circ P$ 为 $T$ 的极分解.证明 $U$ 在 $N(T)^\perp$ 上唯一确定.

6. 假设 $T$ 允许极分解.证明 \eqref{eq:prob44.1} 成立.\\
{[提示: 令 $J=U|_{N(T)}$.]}

7. 证明满足 \eqref{eq:prob44.1} 的每个算子 $T\in\LL(H)$ 允许极分解.

8. 假设 $T$ 满足更强的假设
\begin{equation}
\text{存在由}\ N(T)\ \text{到}\ N(T^\star)\ \text{上的等距}\ J.
\tag{2}
\label{eq:prob44.2}
\end{equation}
证明 $T$ 允许极分解 $T=U\circ P$, 其中 $U$ 是酉算子.

9. 由此推出每个正规算子 $T\in\LL(H)$ 允许极分解 $T=U\circ P$, 其中 $U$ 是酉算子且
$U\circ P=P\circ U$.

10. 证明满足 \eqref{eq:prob44.2} 的每个算子 $T\in\LL(H)$ 可以分解为 $T=P\circ U$,
其中 $U\in\LL(H)$ 是酉算子, $P\in\LL(H)$ 是自伴非负算子.\\
{[提示: 对 $T^\star$ 应用第 8 问.]}

11. 证明满足 \eqref{eq:prob44.1} 的每个算子 $T\in\KK(H)$ 允许极分解 $T=U\circ P$,
其中 $P\in\KK(H)$.

12. 假设 $H$ 可分且 $T\in\KK(H)$ ($T$ 不必满足 \eqref{eq:prob44.1}).证明存在 $H$ 的
两个标准正交基 $(e_n)$ 与 $(f_n)$ 使得对一切 $u\in H$,
\begin{equation*}
Tu=\sum_{n=1}^\infty\alpha_n(u,e_n)f_n,
\end{equation*}
其中 $(\alpha_n)$ 是满足对一切 $n$ 有 $\alpha_n\ge0$ 且当 $n\to\infty$ 时 $\alpha_n\to0$
的序列.计算 $T^\star$.反之, 证明这种形式的任何算子必为紧算子.
\end{problem}

\begin{problem}[第~
\ref{chap:8}~章]
\label{prob:45}
\runin{强极大值原理}
考虑双线性形式
\begin{equation*}
a(u,v)=\int_0^1(pu'v'+quv),
\end{equation*}
其中 $p\in C^1([0,1])$, 在 $(0,1)$ 上 $p\ge\alpha>0$, 且 $q\in C([0,1])$.我们假设
$a$ 在 $H_0^1(0,1)$ 上强制 (但不对 $q$ 作任何符号假设).

给定 $f\in L^2(0,1)$, 设 $u\in H^2(0,1)$ 为
\begin{equation}
\begin{cases}
-(pu')'+qu=f & \text{在}\ (0,1)\ \text{上},\\
u(0)=u(1)=0
\end{cases}
\label{eq:prob45.1}
\end{equation}
的解.假设在 $(0,1)$ 上几乎处处 $f\ge0$ 且 $f\not\equiv0$.我们的目标是证明
\begin{equation}
u'(0)>0,\quad u'(1)<0\quad\text{且在}\ (0,1)\ \text{上}\ u(x)>0.
\label{eq:prob45.2}
\end{equation}

1. 假设 $\psi\in H^1(0,1)$ 满足对一切在 $(0,1)$ 上满足 $v\ge0$ 的 $v\in H_0^1(0,1)$,
$a(\psi,v)\le0$, 且 $\psi(0)\le0,\psi(1)\le0$.证明在 $(0,1)$ 上 $\psi\le0$.\\
{[提示: 取 $v=\psi^+$ 并使用习题 8.11.]}

考虑问题
\begin{equation}
\begin{cases}
-(p\zeta')'+q\zeta=0 & \text{在}\ (0,1)\ \text{上},\\
\zeta(0)=0,\quad\zeta(1)=1.
\end{cases}
\label{eq:prob45.5}
\end{equation}

2. 证明 \eqref{eq:prob45.5} 有唯一解 $\zeta$, 且在 $(0,1)$ 上 $\zeta\ge0$.

3. 验证在 $(0,1)$ 上 $u\ge0$, 并由此推出 $u'(0)\ge0$ 且 $u'(1)\le0$.

4. 证明 $p(1)\abs{u'(1)}=\int_0^1f\zeta$.\\
{[提示: 把 \eqref{eq:prob45.1} 乘以 $\zeta$, 把 \eqref{eq:prob45.5} 乘以 $u$.]}

令 $\phi(x)=(e^{Bx}-1)$, $B>0$.

5. 验证若 $B$ 充分大 (只依赖于 $p$ 与 $q$), 则在 $(0,1)$ 上
$-(p\phi')'+q\phi\le0$.在下文中固定使此式成立的 $B$.

6. 设 $A=(e^B-1)^{-1}$.证明在 $(0,1)$ 上 $\zeta\ge A\phi$.\\
{[提示: 对 $\psi=A\phi-\zeta$ 应用第 1 问.]}

7. 由此推出 $u'(1)<0$.\\
{[提示: 应用第 4 问.]}

8. 验证 $u'(0)>0$.\\
{[提示: 将 $t$ 换为 $1-t$.]}

9. 固定充分小的 $\delta\in(0,1/2)$ 使得对一切 $x\in(0,\delta)$, $\dfrac{u(x)}{x}\ge\frac12u'(0)$,
且对一切 $x\in(1-\delta,1)$, $\dfrac{u(x)}{1-x}\ge\frac12\abs{u'(1)}$.为什么这样的 $\delta$
存在? 设 $v$ 为在 $(\delta,1-\delta)$ 上 $-(pv')'+qv=0$, $v(\delta)=v(1-\delta)=\gamma$
的解, 其中 $\gamma=\dfrac\delta2\min\{u'(0),\abs{u'(1)}\}$.证明在 $(\delta,1-\delta)$
上 $u\ge v\ge0$.

10. 证明在 $(\delta,1-\delta)$ 上 $v>0$.\\
{[提示: 反设对某个 $x_0\in(\delta,1-\delta)$ 有 $v(x_0)=0$, 像习题 8.33 那样应用定理
 \ref{thm:7.3} (Cauchy--Lipschitz--Picard).]}

11. 由此推出对一切 $x\in(0,1)$, $u(x)>0$.

最后, 我们给出强极大值原理的一个更强的形式.

12. 证明存在常数 $a>0$ (只依赖于 $p$ 与 $q$) 使得
\begin{equation*}
u(x)\ge ax(1-x)\int_0^1f(t)t(1-t)dt.
\end{equation*}
[提示: 先从 $p\equiv1$ 且 $q\equiv k^2$ (正实数) 的情形开始; 使用 \eqref{eq:prob45.1}
的显式解.接着考虑 $p\equiv1$ 且不对 $q$ 作进一步假设的情形.最后通过变量代换把一般
情形归结为前述情形.]
\end{problem}

\begin{problem}[第~
\ref{chap:8}~章]
\label{prob:46}
\runin{上下解方法}
设 $h(t):[0,+\infty)\to[0,+\infty)$ 为连续单调不减函数.假设存在两个函数
$v,w\in C^2([0,1])$ 满足
\begin{equation}
\begin{cases}
0\le v\le w & \text{在}\ I=(0,1)\ \text{上},\\
-v''+v\le h(v) & \text{在}\ I\ \text{上},\\
v(0)=v(1)=0,\\
-w''+w\ge h(w) & \text{在}\ I\ \text{上},\\
w(0)\ge0,\quad w(1)\ge0
\end{cases}
\label{eq:prob46.1}
\end{equation}
($v$ 称为下解, $w$ 称为上解).目标是证明问题
\begin{equation}
\begin{cases}
-u''+u=h(u) & \text{在}\ I\ \text{上},\\
u(0)=u(1)=0,\\
v\le u\le w & \text{在}\ I\ \text{上}
\end{cases}
\label{eq:prob46.2}
\end{equation}
存在解 $u\in C^2([0,1])$.

考虑从 $u_0=w$ 开始由下式归纳定义的序列 $(u_n)_{n\ge1}$:
\begin{equation}
\begin{cases}
-u''_n+u_n=h(u_{n-1}) & \text{在}\ I\ \text{上},\ n\ge1,\\
u_n(0)=u_n(1)=0.
\end{cases}
\label{eq:prob46.3}
\end{equation}

1. 证明在 $I$ 上 $v\le u_1\le w$.\\
{[提示: 对 $(u_1-w)$ 与 $(u_1-v)$ 应用极大值原理.]}

2. 由归纳法证明对每个 $n\ge1$, 在 $I$ 上 $v\le u_n$ 且 $u_{n+1}\le u_n$.

3. 由此推出序列 $(u_n)$ 在 $L^2(I)$ 中收敛于极限 $u$, 且在 $L^2(I)$ 中 $h(u_n)\to h(u)$.

4. 证明 $u\in H_0^1(I)$, 且对一切 $\phi\in H_0^1(I)$,
\begin{equation*}
\int_0^1u'\phi'+\int_0^1u\phi=\int_0^1h(u)\phi.
\end{equation*}

5. 得出 $u\in C^2([0,1])$ 是 \eqref{eq:prob46.2} 的古典解.

在下文中我们选取 $h(t)=t^\alpha$, 其中 $0<\alpha<1$.目标是证明存在唯一的函数
$u\in C^2([0,1])$ 满足
\begin{equation}
\begin{cases}
-u''+u=u^\alpha & \text{在}\ I\ \text{上},\\
u(0)=u(1)=0,\\
u(x)>0 & \forall x\in I.
\end{cases}
\label{eq:prob46.4}
\end{equation}

6. 设 $v(x)=\varepsilon\sin(\pi x)$ 且 $w(x)\equiv1$.证明若 $\varepsilon$ 充分小,
假设 \eqref{eq:prob46.1} 满足.由此推出 \eqref{eq:prob46.4} 存在解.

现在转向唯一性问题.设 $u$ 为由上述方法从 $u_0\equiv1$ 出发得到的 \eqref{eq:prob46.4}
的解.设 $\tilde u\in C^2([0,1])$ 为 \eqref{eq:prob46.4} 的另一个解.

7. 证明在 $I$ 上 $\tilde u\le1$.\\
{[提示: 考虑 $\tilde u$ 达到其最大值的点 $x_0\in[0,1]$.]}

8. 证明从 $u_0\equiv1$ 开始由 \eqref{eq:prob46.3} 定义的序列 $(u_n)_{n\ge1}$ 满足在
$I$ 上 $\tilde u\le u_n$, 并由此推出在 $I$ 上 $\tilde u\le u$.

9. 证明 $\int_0^1(\tilde u^\alpha u-u^\alpha\tilde u)=0$.

10. 得出在 $I$ 上 $\tilde u=u$.\\
{[提示: 写出 $\tilde u^\alpha u-u^\alpha\tilde u=u\tilde u(\tilde u^{\alpha-1}-u^{\alpha-1})$,
并注意 $u^{\alpha-1}\le\tilde u^{\alpha-1}$.]}

我们现在给出存在性的另一种证明.对每个 $u\in H_0^1(I)$, 令
\begin{equation*}
F(u)=\frac12\int_0^1(u'^2+u^2)-\int_0^1g(u),
\end{equation*}
其中 $g(t)=\frac{1}{\alpha+1}(t^+)^\alpha+1$, $0<\alpha<1$, $t^+=\max(t,0)$.

11. 证明存在常数 $C$ 使得对一切 $u\in H_0^1(I)$,
$F(u)\ge\frac12\norm{u}_{H^1}^2-C\norm{u}_{H^1}^{\alpha+1}$.

12. 由此推出 $m=\inf_{v\in H_0^1(I)}F(v)>-\infty$, 且下确界可以达到.\\
{[提示: 设 $(u_n)$ 为极小化序列.验证子序列 $(u_{n_k})$ 在 $H_0^1(I)$ 中弱收敛于极限 $u$,
且 $\int_0^1g(u_{n_k})\to\int_0^1g(u)$.提醒读者注意泛函 $F$ 不是凸的; 为什么?]}

13. 证明 $m<0$.\\
{[提示: 证明对满足 $v^+\not\equiv0$ 的一切 $v\in H_0^1(I)$ 及充分小的 $\varepsilon$,
$F(\varepsilon v)<0$.]}

14. 验证对一切 $a,b\in\R$, $g(b)-g(a)\ge(a^+)^\alpha(b-a)$.

15. 设 $u\in H_0^1(I)$ 为 $F$ 在 $H_0^1(I)$ 上的极小元.证明对一切 $v\in H_0^1(I)$,
\begin{equation*}
\int_0^1(u'v'+uv)=\int_0^1(u^+)^\alpha v.
\end{equation*}
[提示: 写出 $F(u)\le F(u+tv)$, 应用第 14 问, 并令 $t\to0$.]

16. 由此推出 $u\in C^2([0,1])$ 是 $-u''+u=(u^+)^\alpha$ 在 $I$ 上, $u(0)=u(1)=0$
的解.证明在 $I$ 上 $u\ge0$ 且 $u\not\equiv0$.

17. 利用强极大值原理得出在 $I$ 上 $u>0$ (见问题 45).
\end{problem}
\begin{problem}[第~
\ref{chap:8}~章]
\label{prob:47}
\runin{Poincaré--Wirtinger 不等式}
设 $I=(0,1)$.

\runin{-- A --}

1. 证明对一切 $u\in W^{1,1}(I)$,
\begin{equation}
\norm{u-\bar u}_{L^\infty(I)}\le\norm{u'}_{L^1(I)},
\label{eq:prob47.1}
\end{equation}
其中 $\bar u=\int_I u$.\\
{[提示: 注意对某个 $x_0\in[0,1]$, $\bar u=u(x_0)$.]}

2. 证明 \eqref{eq:prob47.1} 中的常数 $1$ 是最优的, 即
\begin{equation}
\sup\{\norm{u-\bar u}_{L^\infty(I)};\ u\in W^{1,1}(I)\ \text{且}\ \norm{u'}_{L^1(I)}=1\}=1.
\label{eq:prob47.2}
\end{equation}
[提示: 考虑 $[0,1]$ 上的光滑函数序列 $(u_n)$, 满足在 $(0,1)$ 上对一切 $n$ 有
$u'_n\ge0$, $u_n(1)=1$, 以及在 $[0,1-1/n]$ 上 $u_n(x)=0$.]

3. 证明 \eqref{eq:prob47.2} 中的上确界达不到, 即不存在函数 $u\in W^{1,1}(I)$
使得 $\norm{u-\bar u}_{L^\infty(I)}=1$ 且 $\norm{u'}_{L^1(I)}=1$.

4. 证明对一切 $u\in W_0^{1,1}(I)$,
\begin{equation}
\norm{u}_{L^\infty(I)}\le\frac12\norm{u'}_{L^1(I)}.
\label{eq:prob47.3}
\end{equation}
[提示: 写出 $\abs{u(x)-u(0)}\le\int_0^x\abs{u'(t)}dt$ 以及
$\abs{u(x)-u(1)}\le\int_x^1\abs{u'(t)}dt$.]

5. 证明 $1/2$ 是 \eqref{eq:prob47.3} 中的最佳常数.它能达到吗? \\
{[提示: 固定 $a\in(0,1)$, 考虑在 $(0,a)$ 上递增, 在 $(a,1)$ 上递减且 $u(a)=1$ 的函数 $u\in W_0^{1,1}(I)$.]}

6. 由此推出下述不等式成立:
\begin{equation}
\norm{u-\bar u}_{L^q(I)}\le C\norm{u'}_{L^p(I)}\quad\forall u\in W^{1,p}(I),
\label{eq:prob47.4}
\end{equation}
以及
\begin{equation}
\norm{u}_{L^q(I)}\le C\norm{u'}_{L^p(I)}\quad\forall u\in W_0^{1,p}(I),
\label{eq:prob47.5}
\end{equation}
其中 $1\le q\le\infty$ 且 $1\le p\le\infty$.证明当 $1\le q\le\infty$ 且 $1<p\le\infty$ 时,
\eqref{eq:prob47.4} 与 \eqref{eq:prob47.5} 中的最佳常数可以达到.\\
{[提示: 在满足 $\norm{u-\bar u}_{L^q(I)}=1$ 的 $u\in W^{1,p}(I)$ 类 (相应地
$u\in W_0^{1,p}(I)$ 且 $\norm{u}_{L^q(I)}=1$) 中极小化 $\norm{u'}_{L^p(I)}$.]}

\runin{-- B --}

接下来的目标是求 $p=q=2$ 时 \eqref{eq:prob47.4} 中的最佳常数, 即
\begin{equation}
\norm{u-\bar u}_{L^2(I)}\le C\norm{u'}_{L^2(I)}\quad\forall u\in H^1(I).
\label{eq:prob47.6}
\end{equation}
令 $H=\{f\in L^2(I);\ \int_I f=0\}$ 与 $V=\{v\in H^1(I);\ \int_I v=0\}$.

1. 验证对每个 $f\in H$, 存在唯一的 $u\in V$ 使得对一切 $v\in V$, $\int_I u'v'=\int_I fv$.

2. 证明 $u\in H^2(I)$ 且满足在 $I$ 上几乎处处 $-u''=f$, $u'(0)=u'(1)=0$.

3. 证明由 $Tf=u$ 定义的算子 $T:H\to H$ 是自伴紧算子, 且对一切 $f\in H$, $\int_I fTf\ge0$.

4. 设 $\lambda_1$ 为 $T$ 的最大特征值.证明 \eqref{eq:prob47.6} 在 $C=\sqrt{\lambda_1}$
下成立, 且 $\sqrt{\lambda_1}$ 是 \eqref{eq:prob47.6} 中的最佳常数.\\
{[提示: 使用习题 6.24.]}

5. 显式计算 \eqref{eq:prob47.6} 中的最佳常数.

\runin{-- C --}

1. 证明对一切 $u\in W^{1,1}(I)$,
\begin{equation}
\norm{u-\bar u}_{L^1(I)}\le2\int_I\abs{u'(t)}t(1-t)dt.
\label{eq:prob47.7}
\end{equation}

2. 由此推出对一切 $u\in W^{1,1}(I)$,
\begin{equation}
\norm{u-\bar u}_{L^1(I)}\le\frac12\norm{u'}_{L^1(I)}.
\label{eq:prob47.8}
\end{equation}

3. 证明 \eqref{eq:prob47.8} 中的常数 $1/2$ 是最优的, 即
\begin{equation}
\sup\{\norm{u-\bar u}_{L^1(I)};\ u\in W^{1,1}(I)\ \text{且}\ \norm{u'}_{L^1(I)}=1\}=\frac12.
\label{eq:prob47.9}
\end{equation}

4. \eqref{eq:prob47.9} 中的上确界能达到吗?
\end{problem}

\begin{problem}[第~
\ref{chap:8}~章]
\label{prob:48}
\runin{一个非线性问题}
设 $j:[-1,+1]\to[0,+\infty)$ 为连续凸函数, 满足 $j\in C^2((-1,+1))$, $j(0)=0$,
$j'(0)=0$, 且
\begin{equation*}
\lim_{t\uparrow+1}j'(t)=+\infty,\qquad\lim_{t\downarrow-1}j'(t)=-\infty.
\end{equation*}
(一个应牢记的好例子是 $j(t)=1-\sqrt{1-t^2}$, $t\in[-1,+1]$.) 给定 $f\in L^2(0,1)$,
定义函数 $\phi:H_0^1(0,1)\to(-\infty,+\infty]$:
\begin{equation*}
\phi(v)=\begin{cases}
\dfrac12\int_0^1v'^2+\int_0^1j(v)-\int_0^1fv & \text{若}\ v\in H_0^1(0,1)\ \text{且}\ \norm{v}_{L^\infty}\le1,\\
+\infty & \text{其他}.
\end{cases}
\end{equation*}

1. 验证 $\phi$ 在 $H_0^1(0,1)$ 上是凸的, 下半连续的, 且 $\lim_{\norm{v}_{H_0^1}\to+\infty}\phi(v)=+\infty$.

2. 由此推出存在唯一的 $u\in H_0^1(0,1)$ 使得 $\phi(u)=\min_{v\in H_0^1(0,1)}\phi(v)$.

目标是证明若 $f\in L^\infty(0,1)$, 则 $\norm{u}_{L^\infty(0,1)}<1$, $u\in H^2(0,1)$,
且 $u$ 满足
\begin{equation}
\begin{cases}
-u''+j'(u)=f & \text{在}\ (0,1)\ \text{上},\\
u(0)=u(1)=0.
\end{cases}
\label{eq:prob48.1}
\end{equation}

3. 验证对一切 $t\in[-1,+1]$ 与一切 $a\in(-1,+1)$, $j(t)-j(a)\ge j'(a)(t-a)$.\\
{[提示: 使用 $j$ 的凸性.]}

固定 $a\in[0,1)$.

4. 令 $v=\min(u,a)$.证明 $v\in H_0^1(0,1)$ 且在 $[u\le a]$ 上几乎处处 $v'=u'$,
在 $[u>a]$ 上几乎处处 $v'=0$.\\
{[提示: 写出 $v=a-(a-u)^+$ 并使用习题 8.11.]}

5. 证明 $\frac12\int_{[u>a]}u'^2\le\int_{[u>a]}(f-j'(a))(u-a)$.\\
{[提示: 写出 $\phi(u)\le\phi(v)$, 其中 $v$ 在第 4 问中定义.然后使用第 3 问.]}

6. 选取 $a\in[0,1)$ 使得对一切 $x\in[0,1]$, $f(x)\le j'(a)$, 并证明对一切 $x\in[0,1]$,
$u(x)\le a$.\\
{[提示: 证明 $\int_0^1w'^2=0$, 其中 $w=(u-a)^+$ 属于 $H_0^1(0,1)$; 为什么?]}

7. 得出结论: $\norm{u}_{L^\infty(0,1)}<1$.\\
{[提示: 应用前面的论证, 将 $u$ 换为 $-u$, $j(t)$ 换为 $j(-t)$, $f$ 换为 $-f$.]}

8. 由此推出 $u$ 属于 $H^2(0,1)$ 且满足 \eqref{eq:prob48.1}.\\
{[提示: 写出 $\phi(u)\le\phi(u+\varepsilon v)$, 其中 $v\in H_0^1(0,1)$, $\varepsilon$ 很小.]}

9. 验证若 $f\in C([0,1])$, 则 $u\in C^2([0,1])$.

10. 反之, 证明满足 $\norm{u}_{L^\infty(0,1)}<1$ 且满足 \eqref{eq:prob48.1} 的任何函数
$u\in C^2([0,1])$ 是 $\phi$ 在 $H_0^1(0,1)$ 上的极小元.\\
{[提示: 在第 3 问中取 $t=v(x)$ 与 $a=u(x)$.]}

现在假设 $f\in L^2(0,1)$.令 $f_n=T_nf$, 其中 $T_n$ 是截断算子 (定义在第~
\ref{chap:4}~章定理~
\ref{thm:4.12}
之后).设 $u_n$ 为对应于 $f_n$ 的 \eqref{eq:prob48.1} 的解.

11. 证明当 $n\to\infty$ 时 $\norm{j'(u_n)}_{L^2(0,1)}\le C$.\\
{[提示: 把 \eqref{eq:prob48.1} 乘以 $j'(u_n)$.]}

12. 由此推出 $\norm{u_n}_{H^2(0,1)}\le C$.

13. 证明某个子序列 $(u_{n_k})$ 在 $H^2(0,1)$ 中弱收敛于极限 $u\in H^2(0,1)$, 满足在
$C^1([0,1])$ 中 $u_{n_k}\to u$.证明在 $(0,1)$ 上几乎处处 $\abs{u(x)}<1$, 且
$j'(u)\in L^2(0,1)$.\\
{[提示: 对序列 $j'(u_{n_k})^2$ 应用 Fatou 引理.]}

14. 证明 $j'(u_{n_k})$ 在 $L^2(0,1)$ 中弱收敛于 $j'(u)$, 并由此推出 \eqref{eq:prob48.1} 成立.\\
{[提示: 应用习题 4.16.]}

15. 若额外假设
\begin{equation*}
\liminf_{t\uparrow1}j'(t)(1-t)^{1/3}>0\quad\text{且}\quad\limsup_{t\downarrow-1}j'(t)(1+t)^{1/3}<0,
\end{equation*}
由此推出 $\norm{u}_{L^\infty(0,1)}<1$.\\
{[提示: 反设对某个 $x_0\in(0,1)$ 有 $u(x_0)=1$.验证对一切 $x\in(0,1)$,
$\abs{u'(x)}\le\abs{x-x_0}^{1/2}\norm{u''}_{L^2}$ 以及
$\abs{u(x)-1}\le\frac23\abs{x-x_0}^{3/2}\norm{u''}_{L^2}$.由此推出 $j'(u)\notin L^2(0,1)$.]}
\end{problem}

\begin{problem}[第~
\ref{chap:8}~章]
\label{prob:49}
\runin{Sturm--Liouville 算子特征值的极小化极大原理}
考虑 $(0,1)$ 上的 Sturm--Liouville 算子 $Au=-(pu')'+qu$, 带 Dirichlet 边界条件
$u(0)=u(1)=0$.假设 $p\in C^1([0,1])$, 对一切 $x\in[0,1]$ 有 $p(x)\ge\alpha>0$,
且 $q\in C([0,1])$.令
\begin{equation*}
a(u,v)=\int_0^1(pu'v'+quv)\quad\forall u,v\in H_0^1(0,1).
\end{equation*}
注意我们不对 $q$ 作进一步的假设, 从而双线性形式 $a$ 不必是强制的.固定充分大的 $M$
使得 $\tilde a(u,v)=a(u,v)+M\int_0^1uv$ 是强制的 (例如 $M>-\min_{x\in[0,1]}q(x)$).
设 $(\lambda_k)$ 为 $A$ 的特征值序列.空间 $H=H_0^1(0,1)$ 配以标量积 $\tilde a(u,v)$
(记作 $(u,v)_H$) 以及相应的范数 $\abs{u}_H=\tilde a(u,u)^{1/2}$.给定任何 $f\in L^2(0,1)$,
设 $u\in H_0^1(0,1)$ 为问题 $\tilde a(u,v)=\int_0^1fv$ ($\forall v\in H_0^1(0,1)$) 的
唯一解.令 $u=Tf$, 并把 $T$ 看作由 $H$ 到自身中的算子.

1. 证明 $T$ 自伴且紧.\\
{[提示: 回忆由 $H$ 到 $L^2(0,1)$ 中的恒等映射是紧的.]}

2. 设 $(\lambda_k)$ 为 $A$ 的特征值序列 (在定理  \ref{thm:8.22} 的意义下), 相应的特征
函数为 $(e_k)$, 设 $(\mu_k)$ 为 $T$ 的特征值序列.验证对一切 $k$, $\mu_k>0$, 并证明
对一切 $k$, $\lambda_k=\dfrac{1}{\mu_k}-M$ 且 $T(e_k)=\mu_ke_k$.

3. 证明对一切 $w\in H$, $(Tw,w)_H=\int_0^1w^2$, 并由此推出对一切 $w\in H$ ($w\ne0$),
\begin{equation*}
\frac{1}{R(w)}=\frac{a(w,w)}{\int_0^1w^2}+M,
\end{equation*}
其中 $R$ 是与 $T$ 相关联的 Rayleigh 商, 即 $R(w)=\dfrac{(Tw,w)_H}{\abs{w}_H^2}$
(见问题 37).

4. 证明
\begin{equation}
\lambda_1=\min_{\substack{w\in H_0^1\\w\ne0}}\left\{\frac{a(w,w)}{\int_0^1w^2}\right\},
\label{eq:prob49.1}
\end{equation}
且对一切 $k\ge2$,
\begin{equation*}
\lambda_k=\min\left\{\frac{a(w,w)}{\int_0^1w^2};\ w\in H_0^1(0,1),\ w\ne0,\ \int_0^1we_j=0\ \forall j=1,\dots,k-1\right\}.
\end{equation*}
[提示: 应用问题 37 第 2 问并证明 $(w,e_j)_H=0\iff\int_0^1we_j=0$.]

5. 证明对一切 $k\ge1$,
\begin{equation*}
\lambda_k=\min_{\substack{\Sigma\subset H_0^1(0,1)\\\dim\Sigma=k}}\max_{\substack{u\in\Sigma\\u\ne0}}\left\{\frac{a(u,u)}{\int_0^1u^2}\right\},
\end{equation*}
以及
\begin{equation*}
\lambda_k=\max_{\substack{M\subset H_0^1(0,1)\\\operatorname{codim}M=k-1}}\min_{\substack{u\in M\\u\ne0}}\left\{\frac{a(u,u)}{\int_0^1u^2}\right\},
\end{equation*}
其中 $\Sigma$ 与 $M$ 是 $H_0^1(0,1)$ 的闭子空间.\\
{[提示: 应用问题 37 第 7 问.]}

6. 对带 Neumann 边界条件的 Sturm--Liouville 算子证明类似的结果.

现在我们回到公式 \eqref{eq:prob49.1} 并讨论对应于第一个特征值 $\lambda_1$ 的特征函数
的进一步性质.特别地, 我们将看到存在一个生成与 $\lambda_1$ 相关联的特征子空间的正特征函数.

7. 设 $w_0\in H_0^1(0,1)$ 为 \eqref{eq:prob49.1} 的满足 $\int_0^1w_0^2=1$ 的极小元.
证明在 $(0,1)$ 上 $Aw_0=\lambda_1w_0$.

8. 令 $w_1=\abs{w_0}$.验证 $w_1$ 也是 \eqref{eq:prob49.1} 的极小元, 并由此推出在
$(0,1)$ 上 $Aw_1=\lambda_1w_1$.\\
{[提示: 使用习题 8.11.]}

9. 证明在 $(0,1)$ 上 $w_1>0$, $w'_1(0)>0$, 且 $w'_1(1)<0$.\\
{[提示: 对算子 $A+M$ 应用强极大值原理; 见问题 45.]}

10. 假设 $w\in H_0^1(0,1)$ 在 $(0,1)$ 上满足 $Aw=\lambda_1w$.证明 $w$ 是 $w_1$ 的倍数.\\
{[提示: 回忆特征值是单的; 见习题 8.33.找出不依赖特征值单性的另一种证明; 在 (2) 中使用
$w^2/w_1$ 作为检验函数.]}

11. 证明在 $(0,1)$ 上满足 $A\psi=\mu\psi$, $\psi\ge0$ 且 $\int_0^1\psi^2=1$ ($\mu\in\R$)
的任何函数 $\psi\in H_0^1(0,1)$ 必与 $w_1$ 重合.\\
{[提示: 若 $\mu\ne\lambda_1$, 验证 $\int_0^1\psi w_1=0$.由此推出 $\mu=\lambda_1$.]}
\end{problem}

\begin{problem}[第~
\ref{chap:8}~章]
\label{prob:50}
\runin{另一个非线性问题}
设 $q\in C([0,1])$ 并考虑双线性形式 $a(u,v)=\int_0^1(u'v'+quv)$ ($u,v\in H_0^1(0,1)$).
假设存在 $v_1\in H_0^1(0,1)$ 使得
\begin{equation}
a(v_1,v_1)<0.
\label{eq:prob50.1}
\end{equation}

1. 验证假设 \eqref{eq:prob50.1} 等价于 $\lambda_1(A)<0$, 其中 $\lambda_1(A)$ 是带零
Dirichlet 条件的算子 $Au=-u''+qu$ 的第一个特征值.

2. 验证
\begin{equation}
-\infty<m=\inf_{u\in H_0^1(0,1)}\left\{\frac12a(u,u)+\frac14\int_0^1\abs{u}^4\right\}<0.
\label{eq:prob50.3}
\end{equation}
[提示: 使用充分小的 $\varepsilon>0$ 的 $u=\varepsilon v_1$.]

3. 证明 \eqref{eq:prob50.3} 中的下确界由某个 $u_0$ 达到.\\
{[警告: \eqref{eq:prob50.3} 中的泛函不是凸的; 为什么?]}

我们的目标是证明 \eqref{eq:prob50.3} 恰好有两个极小元.

4. 证明 $u_0$ 属于 $C^2([0,1])$ 且满足
\begin{equation}
\begin{cases}
-u''+qu+u^3=0 & \text{在}\ (0,1)\ \text{上},\\
u(0)=u(1)=0.
\end{cases}
\label{eq:prob50.4}
\end{equation}

5. 令 $u_1=\abs{u_0}$.证明 $u_1$ 也是 \eqref{eq:prob50.3} 的极小元.由此推出 $u_1$
满足 \eqref{eq:prob50.4}.\\
{[提示: 应用习题 8.11.]}

6. 证明对一切 $x\in(0,1)$, $u_1(x)>0$, $u'_1(0)>0$, 且 $u'_1(1)<0$.\\
{[提示: 选取充分大的常数 $a$ 使得 $-u''_1+a^2u_1=f\ge0,f\not\equiv0$.然后使用强极大值原理.]}

7. 设 $u_0\in H_0^1(0,1)$ 再次为 \eqref{eq:prob50.3} 中的任意极小元.证明或者对一切
$x\in(0,1)$ 有 $u_0(x)>0$, 或者对一切 $x\in(0,1)$ 有 $u_0(x)<0$.\\
{[提示: 验证对一切 $x\in(0,1)$, $\abs{u_0(x)}>0$.]}

8. 设 $U_1$ 为满足在 $[0,1]$ 上 $U_1\ge0$ 且 $U_1\not\equiv0$ 的 \eqref{eq:prob50.4}
的任何解.令 $\rho_1=U_1^2$.考虑在集合
\begin{equation*}
K=\{\rho\in H_0^1(0,1);\ \rho\ge0\ \text{在}\ (0,1)\ \text{上且}\ \sqrt{\rho}\in H_0^1(0,1)\}
\end{equation*}
上定义的泛函 $\Phi(\rho)=\int_0^1\left[\abs{(\sqrt{\rho})'}^2+q\rho+\frac12\rho^2\right]$.
证明对一切 $\rho\in K$,
\begin{equation}
\Phi(\rho)-\Phi(\rho_1)\ge\frac12\int_0^1(\rho-\rho_1)^2.
\label{eq:prob50.5}
\end{equation}
[提示: 设 $u\in C_c^1((0,1))$.注意在 $(0,1)$ 上
$\dfrac{2U'_1uu'}{U_1}\le u'^2+\dfrac{U_1'^2u^2}{U_1^2}$, 由此推出 (利用分部积分)
对一切 $u\in H_0^1(0,1)$, $\int_0^1(u'^2-U_1'^2)\ge-\int_0^1\dfrac{U''_1}{U_1}(u^2-U_1^2)$.
然后应用方程 \eqref{eq:prob50.4} 建立 \eqref{eq:prob50.5}.]

9. 由此推出存在 \eqref{eq:prob50.4} 的恰好一个在 $[0,1]$ 上满足 $u\ge0$ 的非平凡解.
记作 $U_0$.\\
{[评注: 一般说来 \eqref{eq:prob50.4} 存在许多变号解.]}

10. 证明 \eqref{eq:prob50.3} 恰好存在两个极小元: $U_0$ 与 $-U_0$.
\end{problem}

\begin{problem}[第~
\ref{chap:8}~章]
\label{prob:51}
\runin{调和振子. Hermite 多项式}
设 $p\in C(\R)$ 满足在 $\R$ 上 $p\ge0$.考虑空间
\begin{equation*}
V=\left\{v\in H^1(\R);\ \int_{-\infty}^{+\infty}pv^2<\infty\right\},
\end{equation*}
它配以标量积 $(u,v)_V=\int_{-\infty}^{+\infty}(u'v'+uv+puv)$ 以及相应的范数
$\abs{u}_V=(u,u)_V^{1/2}$.

1. 验证 $V$ 是可分 Hilbert 空间.

2. 证明 $C_c^\infty(\R)$ 在 $V$ 中稠密.\\
{[提示: 设 $\zeta_n$ 为定理  \ref{thm:8.7} 证明中的截断函数序列.给定 $u\in V$, 考虑
$\zeta_nu$, 然后使用卷积.]}

考虑双线性形式 $a(u,v)=\int_{-\infty}^{+\infty}(u'v'+puv)$, $u,v\in V$.
在下文中我们假设存在常数 $\delta>0$ 与 $A>0$ 使得
\begin{equation}
p(x)\ge\delta\quad\forall x\in\R\ \text{满足}\ \abs{x}\ge A.
\label{eq:prob51.1}
\end{equation}

3. 证明 $a$ 在 $V$ 上是强制的.由此推出对每个 $f\in L^2(\R)$, 存在问题
\begin{equation}
a(u,v)=\int_{-\infty}^{+\infty}fv\quad\forall v\in V
\label{eq:prob51.2}
\end{equation}
的唯一解 $u\in V$.

4. 假设 $f\in L^2(\R)\cap C(\R)$, 证明 $u$ 满足
\begin{equation}
\begin{cases}
u\in C^2(\R),\\
-u''+pu=f & \text{在}\ \R\ \text{上},\\
\lim_{\abs{x}\to\infty}u(x)=0.
\end{cases}
\label{eq:prob51.3}
\end{equation}

5. 反之, 证明 \eqref{eq:prob51.3} 的任何解 $u$ 属于 $V$ 且满足 \eqref{eq:prob51.2}.\\
{[提示: 把方程 $-u''+pu=f$ 乘以 $\zeta_n^2u$ 并使用 $a$ 强制的事实.]}

在下文中假设
\begin{equation}
\lim_{\abs{x}\to\infty}p(x)=+\infty.
\label{eq:prob51.4}
\end{equation}

6. 给定 $f\in L^2(\R)$, 令 $u=Tf$, 其中 $u$ 是 \eqref{eq:prob51.2} 的解.证明
$T:L^2(\R)\to L^2(\R)$ 自伴且紧.\\
{[提示: 使用推论  \ref{cor:4.27} 验证 $V\subset L^2(\R)$ 紧嵌入.]}

7. 由此推出存在满足当 $n\to\infty$ 时 $\lambda_n\to\infty$ 的正数序列 $(\lambda_n)$,
以及满足下式的 $L^2(\R)$ 的 Hilbert 基 $(e_n)$:
\begin{equation}
\begin{cases}
e_n\in V\cap C^2(\R),\\
-e''_n+pe_n=\lambda_ne_n & \text{在}\ \R\ \text{上}.
\end{cases}
\label{eq:prob51.5}
\end{equation}

在下文中取 $p(x)=x^2$.

8. 验证 \eqref{eq:prob51.5} 允许形如 $e_n(x)=e^{-x^2/2}P_n(x)$ 的解, 其中 $\lambda_n=2n+1$,
$P_n(x)$ 是 $n$ 次多项式.
\end{problem}
